%%%
%%% X
%%%
\section*{X}\addcontentsline{toc}{section}{X}\addcontentsline{loh}{figure}{\#\#\#\#\#\#\#\# X}

%%%%%%%%%% 夕 %%%%%%%%%%
\subsection*{夕}\addcontentsline{loh}{figure}{夕 \dpy{xi1}}

\begin{EntryWithPhonetic}{夕}{xi1}{3}{⼣}[Kangxi 36]
  \definition*{s.}{Sobrenome: Xi}
  \definition{s.}{pôr do sol; crepúsculo | tarde; noite}
\end{EntryWithPhonetic}

\begin{EntryWithPhonetic}{夕阳}{xi1yang2}{3,6}{⼣,⾩}
  \definition{s.}{pôr do sol | ocaso; metaforicamente, refere"-se à velhice; ou a indústrias tradicionais e não competitivas, sem perspectivas de futuro}
  \seealsoref{日出}{ri4chu1}
  \synonymref{落日}{luo4ri4}
  \synonymref{斜阳}{xie2yang2}
  \antonymref{旭日}{xu4ri4}
\end{EntryWithPhonetic}

%%%%%%%%%% 吸 %%%%%%%%%%
\subsection*{吸}\addcontentsline{loh}{figure}{吸 \dpy{xi1}}

\begin{EntryWithPhonetic}{吸}{xi1}{6}{⼝}[HSK 4]
  \definition{v.}{inalar; inspirar; aspirar | sugar (líquidos) | absorver; sugar | atrair; atrair para si mesmo | aspirar; introdução de líquidos, gases, etc. no corpo}
  \antonymref{呼}{hu1}
\end{EntryWithPhonetic}

\begin{EntryWithPhonetic}{吸毒}{xi1du2}{6,9}{⼝,⽏}[HSK 6]
  \definition{s.}{droga}
  \definition{v.}{usar drogas viciantes; ser viciado em um narcótico; consumir drogas}
\end{EntryWithPhonetic}

\begin{EntryWithPhonetic}{吸管}{xi1guan3}{6,14}{⼝,⽵}[HSK 4]
  \definition[根,个,支]{s.}{tubo de sucção; sugador; canudo (para beber); refere"-se ao tubo fino usado para sugar bebidas | conta"-gotas; pipeta; cateter para bombeamento de líquidos usando pressão de ar}
\end{EntryWithPhonetic}

\begin{EntryWithPhonetic}{吸纳}{xi1na4}{6,7}{⼝,⽷}[HSK 7-9]
  \definition{v.}{absorver; atrair; extrair | aceitar | admitir}
  \synonymref{接纳}{jie1na4}
  \synonymref{吸取}{xi1qu3}
  \synonymref{吸收}{xi1shou1}
\end{EntryWithPhonetic}

\begin{EntryWithPhonetic}{吸取}{xi1qu3}{6,8}{⼝,⼜}[HSK 7-9]
  \definition{v.}{assimilar; absorver (informação, conhecimento, etc.); absorver e adotar; este termo tem um âmbito de uso mais restrito; seu objeto geralmente são conceitos abstratos como experiência ou lições aprendidas | absorver; ingerir (líquido, gás, etc.)}
  \synonymref{接收}{jie1shou1}
  \synonymref{吸纳}{xi1na4}
  \synonymref{吸收}{xi1shou1}
  \antonymref{灌输}{guan4shu1}
\end{EntryWithPhonetic}

\begin{EntryWithPhonetic}{吸收}{xi1shou1}{6,6}{⼝,⽁}[HSK 4]
  \definition{v.}{imbuir; absorver; assimilar; sugar;  chupar; (animais, plantas, etc.) extrair material de fora dos tecidos para o interior dos tecidos | absorver; chupar;  sugar alguma substância de fora para dentro | recrutar; alistar; inscrever"-se; matricular"-se; admitir; (organizações ou coletivos) aceitar novos membros | absorver; aproveitar e usar a experiência, o conhecimento, o dinheiro e outras coisas valiosas de outras pessoas | absorver; diminuir, atenuar ou eliminar determinados efeitos ou fenômenos}
\end{EntryWithPhonetic}

\begin{EntryWithPhonetic}{吸铁石}{xi1tie3shi2}{6,10,5}{⼝,⾦,⽯}
  \definition{s.}{imã}
  \seealsoref{磁铁}{ci2tie3}
\end{EntryWithPhonetic}

\begin{EntryWithPhonetic}{吸烟}{xi1/yan1}{6,10}{⼝,⽕}[HSK 4]
  \definition{v.+compl.}{fumar}
\end{EntryWithPhonetic}

\begin{EntryWithPhonetic}{吸引}{xi1yin3}{6,4}{⼝,⼸}[HSK 4]
  \definition{v.}{atrair; apelar para; chamar a atenção de outros objetos, forças ou pessoas para si mesmo}
\end{EntryWithPhonetic}

%%%%%%%%%% 西 %%%%%%%%%%
\subsection*{西}\addcontentsline{loh}{figure}{西 \dpy{xi1}}

\begin{EntryWithPhonetic}{西}{xi1}{6}{⾑}[HSK 1][Kangxi 146]
  \definition*{s.}{Espanha, abreviatura de 西班牙 | Paraíso Ocidental | Sobrenome: Xi}
  \definition{s.}{oeste; uma das quatro direções básicas, o lado onde o sol se põe | ocidental; refere"-se ao Ocidente (principalmente aos países europeus e americanos) | aqui e ali; significa 到处 ou 零散, 没有次序}
  \seealsoref{到处}{dao4chu4}
  \seealsoref{零散}{ling2san3}
  \seealsoref{没有次序}{mei2you3 ci4xu4}
  \seealsoref{西班牙}{xi1ban1ya2}
  \seealsoref{西边}{xi1bian5}
  \synonymref{西边}{xi1bian5}
  \antonymref{东}{dong1}
  \antonymref{中}{zhong4}
\end{EntryWithPhonetic}

\begin{EntryWithPhonetic}{西安}{xi1'an1}{6,6}{⾑,⼧}
  \definition*{s.}{Xi'an, Capital da Província de Shaanxi}
\end{EntryWithPhonetic}

\begin{EntryWithPhonetic}{西班牙}{xi1ban1ya2}{6,10,4}{⾑,⽟,⽛}
  \definition*{s.}{Espanha}
\end{EntryWithPhonetic}

\begin{EntryWithPhonetic}{西班牙文}{xi1ban1ya2wen2}{6,10,4,4}{⾑,⽟,⽛,⽂}
  \definition{s.}{espanhol, língua espanhola}
  \seealsoref{西文}{xi1wen2}
\end{EntryWithPhonetic}

\begin{EntryWithPhonetic}{西班牙语}{xi1ban1ya2yu3}{6,10,4,9}{⾑,⽟,⽛,⾔}[HSK 6]
  \definition[句]{s.}{espanhol | língua espanhola}
  \seealsoref{西语}{xi1yu3}
\end{EntryWithPhonetic}

\begin{EntryWithPhonetic}{西半球}{xi1ban4qiu2}{6,5,11}{⾑,⼗,⽟}
  \definition*{s.}{Hemisfério Ocidental}
  \antonymref{东半球}{dong1ban4qiu2}
\end{EntryWithPhonetic}

\begin{EntryWithPhonetic}{西北}{xi1bei3}{6,5}{⾑,⼔}[HSK 2]
  \definition{s.}{noroeste | noroeste da China; o Noroeste}
  \antonymref{东南}{dong1nan2}
\end{EntryWithPhonetic}

\begin{EntryWithPhonetic}{西边}{xi1bian5}{6,5}{⾑,⾡}[HSK 1]
  \definition{s.}{lado oeste; (oeste) Uma das quatro direções principais; uma das direções cardeais}
  \seealsoref{西}{xi1}
  \seealsoref{西部}{xi1bu4}
  \synonymref{西}{xi1}
  \synonymref{西部}{xi1bu4}
  \synonymref{西方}{xi1fang1}
  \antonymref{东边}{dong1bian1}
  \antonymref{东方}{dong1fang1}
  \antonymref{东面}{dong1mian4}
\end{EntryWithPhonetic}

\begin{EntryWithPhonetic}{西部}{xi1bu4}{6,10}{⾑,⾢}[HSK 3]
  \definition{s.}{(EUA) filme de faroeste; filme de \emph{cowboys}; um faroeste | filme da região ocidental (China) | parte ocidental; região oeste da China}
\end{EntryWithPhonetic}

\begin{EntryWithPhonetic}{西餐}{xi1can1}{6,16}{⾑,⾷}[HSK 2]
  \definition[份,顿,桌]{s.}{comida ocidental; comida de estilo ocidental, comida com garfo e faca (diferente da 中餐)}
  \seealsoref{中餐}{zhong1can1}
  \antonymref{中餐}{zhong1can1}
\end{EntryWithPhonetic}

\begin{EntryWithPhonetic}{西方}{xi1fang1}{6,4}{⾑,⽅}[HSK 2]
  \definition{s.}{oeste | o Ocidente; o Oeste; países europeus e americanos | Paraíso Ocidental, termo budista}
  \seealsoref{西部}{xi1bu4}
  \synonymref{西边}{xi1bian5}
  \synonymref{西部}{xi1bu4}
  \antonymref{东边}{dong1bian1}
  \antonymref{东方}{dong1fang1}
  \antonymref{东面}{dong1mian4}
\end{EntryWithPhonetic}

\begin{EntryWithPhonetic}{西瓜}{xi1gua5}{6,5}{⾑,⽠}[HSK 4]
  \definition[个,颗,粒]{s.}{melancia; fruto que é uma baga de formato grande, globular ou oval, com muita polpa aguada e doce}
\end{EntryWithPhonetic}

\begin{EntryWithPhonetic}{西红柿}{xi1hong2shi4}{6,6,9}{⾑,⽷,⽊}[HSK 5]
  \definition[种,只,株]{s.}{tomate}
\end{EntryWithPhonetic}

\begin{EntryWithPhonetic}{西兰花}{xi1lan2hua1}{6,5,7}{⾑,⼋,⾋}
  \definition{s.}{brócolis}
  \seealsoref{西蓝花}{xi1lan2hua1}
\end{EntryWithPhonetic}

\begin{EntryWithPhonetic}{西蓝花}{xi1lan2hua1}{6,13,7}{⾑,⾋,⾋}
  \variantof{西兰花}
\end{EntryWithPhonetic}

\begin{EntryWithPhonetic}{西面}{xi1mian4}{6,9}{⾑,⾯}
  \definition{s.}{oeste | lado oeste}
\end{EntryWithPhonetic}

\begin{EntryWithPhonetic}{西南}{xi1nan2}{6,9}{⾑,⼗}[HSK 2]
  \definition{s.}{sudoeste | o Sudoeste; Sudoeste da China}
  \antonymref{东北}{dong1bei3}
\end{EntryWithPhonetic}

\begin{EntryWithPhonetic}{西文}{xi1wen2}{6,4}{⾑,⽂}
  \definition{s.}{espanhol | língua espanhola}
  \seealsoref{西班牙文}{xi1ban1ya2wen2}
\end{EntryWithPhonetic}

\begin{EntryWithPhonetic}{西西}{xi1 xi1}{6,6}{⾑,⾑}
  \definition{num.}{Empréstimo Linguístico, Obsoleto: centímetro cúbico (cc) | mililitro; uma unidade métrica de volume, um mililitro equivale a um milésimo de litro}
\end{EntryWithPhonetic}

\begin{EntryWithPhonetic}{西药}{xi1yao4}{6,9}{⾑,⾋}
  \definition[片,粒]{s.}{medicina ocidental; refere"-se aos medicamentos usados na medicina ocidental, geralmente feitos por métodos sintéticos ou extraídos de produtos naturais, como comprimidos anti"-inflamatórios, aspirina, tintura de iodo, penicilina, etc.}
\end{EntryWithPhonetic}

\begin{EntryWithPhonetic}{西医}{xi1yi1}{6,7}{⾑,⼖}[HSK 2]
  \definition[名,位]{s.}{medicina ocidental; medicina introduzida na China a partir da Europa e da América | um médico treinado em medicina ocidental}
  \antonymref{中医}{zhong1yi1}
\end{EntryWithPhonetic}

\begin{EntryWithPhonetic}{西语}{xi1yu3}{6,9}{⾑,⾔}
  \definition{s.}{línguas ocidentais | espanhol | língua espanhola}
  \seealsoref{西班牙语}{xi1ban1ya2yu3}
\end{EntryWithPhonetic}

\begin{EntryWithPhonetic}{西藏}{xi1zang4}{6,17}{⾑,⾋}
  \definition*{s.}{Xizang; Região Autônoma do Tibete, 西藏自治区}
  \seealsoref{西藏自治区}{xi1zang4zi4zhi4qu1}
\end{EntryWithPhonetic}

\begin{EntryWithPhonetic}{西藏自治区}{xi1zang4zi4zhi4qu1}{6,17,6,8,4}{⾑,⾋,⾃,⽔,⼖}
  \definition*{s.}{Região Autônoma do Tibete}
\end{EntryWithPhonetic}

\begin{EntryWithPhonetic}{西装}{xi1zhuang1}{6,12}{⾑,⾐}[HSK 5]
  \definition[件,套,个]{s.}{terno; roupas de estilo ocidental; roupas ocidentais, divididas em masculinas e femininas}
\end{EntryWithPhonetic}

%%%%%%%%%% 希 %%%%%%%%%%
\subsection*{希}\addcontentsline{loh}{figure}{希 \dpy{xi1}}

\begin{EntryWithPhonetic}{希}{xi1}{7}{⼱}
  \definition*{s.}{Sobrenome: Xi}
  \definition{v.}{ter esperança}
\end{EntryWithPhonetic}

\begin{EntryWithPhonetic}{希望}{xi1wang4}{7,11}{⼱,⽉}[HSK 3]
  \definition[个,丝,点]{s.}{esperança; desejo; expectativa; a possibilidade de alcançar um determinado objetivo ou de ocorrer uma determinada situação ideal no futuro | aquilo em que a esperança é depositada; o objeto da esperança}
  \definition{v.}{ter esperança; desejar; esperar; pensar em alcançar algum objetivo ou que alguma situação ocorra}
\end{EntryWithPhonetic}

%%%%%%%%%% 昔 %%%%%%%%%%
\subsection*{昔}\addcontentsline{loh}{figure}{昔 \dpy{xi1}}

\begin{EntryWithPhonetic}{昔}{xi1}{8}{⽇}
  \definition{s.}{tempos antigos; o passado; era uma vez}
\end{EntryWithPhonetic}

\begin{EntryWithPhonetic}{昔日}{xi1ri4}{8,4}{⽇,⽇}[HSK 7-9]
  \definition{s.}{o passado; no passado; tempos passados; (em) dias (ou tempos) antigos; dias antes}
  \synonymref{从前}{cong2qian2}
  \synonymref{往日}{wang3ri4}
\end{EntryWithPhonetic}

%%%%%%%%%% 息 %%%%%%%%%%
\subsection*{息}\addcontentsline{loh}{figure}{息 \dpy{xi1}}

\begin{EntryWithPhonetic}{息}{xi1}{10}{⼼}
  \definition*{s.}{Sobrenome: Xi}
  \definition{s.}{respiração; o ar que entra e sai durante a respiração | notícias; mensagem | juros | interesse | os filhos de alguém}
  \definition{v.}{cessar; parar | descansar | crescer; multiplicar; procriar}
  \synonymref{休}{xiu1}
  \antonymref{作}{zuo4}
\end{EntryWithPhonetic}

\begin{EntryWithPhonetic}{息息相关}{xi1xi1-xiang1guan1}{10,10,9,6}{⼼,⼼,⽬,⼋}[HSK 7-9]
  \definition{expr.}{estreitamente relacionado; intimamente relacionado; inextricavelmente ligado; interconectado; estreitamente conectado; ``A respiração está interligada.'', é uma metáfora para um relacionamento íntimo; estreitamente ligado; inter"-relacionado; extremamente relevante}
  \seealsoref{息息相通}{xi1xi1-xiang1tong1}
  \synonymref{息息相通}{xi1xi1-xiang1tong1}
\end{EntryWithPhonetic}

\begin{EntryWithPhonetic}{息息相通}{xi1xi1-xiang1tong1}{10,10,9,10}{⼼,⼼,⽬,⾡}
  \definition{expr.}{mutuamente conectados; estreitamente ligado; inter"-relacionado; intimamente relacionado}
\end{EntryWithPhonetic}

%%%%%%%%%% 牺 %%%%%%%%%%
\subsection*{牺}\addcontentsline{loh}{figure}{牺 \dpy{xi1}}

\begin{EntryWithPhonetic}{牺}{xi1}{10}{⽜}
  \definition{s.}{um animal de cor uniforme para sacrifício; sacrifício; gado com pelagem pura usado para sacrifício}
\end{EntryWithPhonetic}

\begin{EntryWithPhonetic}{牺牲}{xi1sheng1}{10,9}{⽜,⽜}[HSK 6]
  \definition[份]{s.}{sacrifício; um animal abatido para sacrifício; refere"-se ao sacrifício da própria vida ou dos próprios interesses por um propósito justo, ou refere"-se ao preço pago por um determinado propósito}
  \definition{v.}{sacrificar"-se; morrer como mártir; dar a própria vida; sacrificar sua vida pela justiça | sacrificar; desistir; fazer algo às custas de; geralmente se refere a pagar um preço ou sofrer danos por alguém ou algo}
\end{EntryWithPhonetic}

%%%%%%%%%% 悉 %%%%%%%%%%
\subsection*{悉}\addcontentsline{loh}{figure}{悉 \dpy{xi1}}

\begin{EntryWithPhonetic}{悉}{xi1}{11}{⼼}
  \definition*{s.}{Sobrenome: Xi}
  \definition{adj.}{tudo; inteiro; total | detalhado}
  \definition{v.}{saber; aprender; ser informado de}
\end{EntryWithPhonetic}

\begin{EntryWithPhonetic}{悉尼}{xi1ni2}{11,5}{⼼,⼫}
  \definition*{s.}{Sidney, capital de Nova Gales do Sul, Austrália}
\end{EntryWithPhonetic}

\begin{EntryWithPhonetic}{悉数}{xi1shu3}{11,13}{⼼,⽁}
  \definition{adv.}{Literário: listar; enumerar em detalhes; contar por completo}
  \seeref{xi1shu4}
\end{EntryWithPhonetic}

\begin{EntryWithPhonetic}{悉数}{xi1shu4}{11,13}{⼼,⽁}
  \definition{adv.}{todos | cada um | toda a soma}
  \seeref{xi1shu3}
  \synonymref{全部}{quan2bu4}
  \synonymref{全面}{quan2mian4}
  \synonymref{全体}{quan2ti3}
  \synonymref{所有}{suo3you3}
  \synonymref{统统}{tong3tong3}
  \synonymref{一切}{yi2qie4}
  \synonymref{整个}{zheng3ge4}
  \synonymref{总共}{zong3gong4}
\end{EntryWithPhonetic}

\begin{EntryWithPhonetic}{悉心}{xi1xin1}{11,4}{⼼,⼼}
  \definition{v.}{dedicar toda a atenção; ter o máximo cuidado; usar toda a  força}
  \synonymref{精心}{jing1xin1}
  \synonymref{用心}{yong4/xin1}
  \antonymref{敷衍}{fu1yan3}
  \antonymref{马虎}{ma3hu5}
  \antonymref{疏忽}{shu1hu5}
  \antonymref{应付}{ying4fu5}
\end{EntryWithPhonetic}

%%%%%%%%%% 稀 %%%%%%%%%%
\subsection*{稀}\addcontentsline{loh}{figure}{稀 \dpy{xi1}}

\begin{EntryWithPhonetic}{稀}{xi1}{12}{⽲}[HSK 7-9]
  \definition{adj.}{raro; escasso; incomum | escasso; espalhado | aquoso; ralo | dispersos; as coisas estão muito distantes umas das outras; os espaços entre as partes das coisas são grandes | ralo; o líquido contém uma pequena quantidade de uma determinada substância}
  \definition{adv.}{(modificando certos adjetivos) muito; extremamente | tão; muito}
  \definition{s.}{algo líquido ou de consistência rala; aquoso}
  \synonymref{少}{shao3}
  \synonymref{疏}{shu1}
  \antonymref{稠}{chou2}
  \antonymref{密}{mi4}
  \antonymref{浓}{nong2}
\end{EntryWithPhonetic}

\begin{EntryWithPhonetic}{稀罕}{xi1han5}{12,7}{⽲,⽹}[HSK 7-9]
  \definition{adj.}{raro; escasso; incomum; raro e diferente do comum}
  \definition{s.}{novidade; raridade; coisas raras}
  \definition{v.}{estimar; valorizar como uma raridade; gostar de algo porque é incomum}
  \synonymref{罕见}{han3jian4}
  \synonymref{奇怪}{qi2guai4}
  \synonymref{特别}{te4bie2}
  \synonymref{稀奇}{xi1qi2}
  \synonymref{稀少}{xi1shao3}
  \synonymref{新鲜}{xin1xian1}
  \antonymref{平常}{ping2chang2}
  \antonymref{普遍}{pu3bian4}
\end{EntryWithPhonetic}

\begin{EntryWithPhonetic}{稀奇}{xi1qi2}{12,8}{⽲,⼤}[HSK 7-9]
  \definition{adj.}{estranho; raro; incomum; extraordinário; raro e inédito}
  \synonymref{别致}{bie2zhi4}
  \synonymref{罕见}{han3jian4}
  \synonymref{奇怪}{qi2guai4}
  \synonymref{奇迹}{qi2ji4}
  \synonymref{奇妙}{qi2miao4}
  \synonymref{特别}{te4bie2}
  \synonymref{稀罕}{xi1han5}
  \synonymref{新奇}{xin1qi2}
  \synonymref{新鲜}{xin1xian1}
  \synonymref{新颖}{xin1ying3}
  \antonymref{常见}{chang2jian4}
  \antonymref{平常}{ping2chang2}
  \antonymref{平凡}{ping2fan2}
  \antonymref{普通}{pu3tong1}
\end{EntryWithPhonetic}

\begin{EntryWithPhonetic}{稀少}{xi1shao3}{12,4}{⽲,⼩}[HSK 7-9]
  \definition{adj.}{poucos; raros; escassos; as coisas parecem menos}
  \synonymref{单独}{dan1du2}
  \synonymref{罕见}{han3jian4}
  \synonymref{特别}{te4bie2}
  \synonymref{稀奇}{xi1qi2}
  \synonymref{稀罕}{xi1han5}
  \antonymref{稠密}{chou2mi4}
  \antonymref{茂密}{mao4mi4}
  \antonymref{茂盛}{mao4sheng4}
  \antonymref{密集}{mi4ji2}
  \antonymref{普遍}{pu3bian4}
  \antonymref{众多}{zhong4duo1}
\end{EntryWithPhonetic}

%%%%%%%%%% 腊 %%%%%%%%%%
\subsection*{腊}\addcontentsline{loh}{figure}{腊 \dpy{xi1}}

\begin{EntryWithPhonetic}{腊}{xi1}{12}{⾁}
  \definition{s.}{carne seca}
  \seeref{la4}
\end{EntryWithPhonetic}

%%%%%%%%%% 锡 %%%%%%%%%%
\subsection*{锡}\addcontentsline{loh}{figure}{锡 \dpy{xi1}}

\begin{EntryWithPhonetic}{锡}{xi1}{13}{⾦}[HSK 7-9]
  \definition*{s.}{Sobrenome: Xi}
  \definition{s.}{estanho, Sn}
  \definition{v.}{Literário: conceder; outorgar; dar}
\end{EntryWithPhonetic}

%%%%%%%%%% 熄 %%%%%%%%%%
\subsection*{熄}\addcontentsline{loh}{figure}{熄 \dpy{xi1}}

\begin{EntryWithPhonetic}{熄}{xi1}{14}{⽕}
  \definition{v.}{extinguir; apagar | (fogo, luz, etc.) apagar"-se; extinguir"-se}
\end{EntryWithPhonetic}

\begin{EntryWithPhonetic}{熄火}{xi1/huo3}{14,4}{⽕,⽕}[HSK 7-9]
  \definition{v.+compl.}{(combustível, fogão, etc.) parar de queimar; extinguir"-se | (um motor) parar de funcionar; morrer | interromper a queima (de combustível); desligar; parar (um motor, etc.) | (competições) acabar; finalizar}
  \antonymref{发动}{fa1dong4}
\end{EntryWithPhonetic}

%%%%%%%%%% 熙 %%%%%%%%%%
\subsection*{熙}\addcontentsline{loh}{figure}{熙 \dpy{xi1}}

\begin{EntryWithPhonetic}{熙}{xi1}{14}{⽕}
  \definition{adj.}{Literário: brilhante; ensolarado | Literário: próspero |Literário: feliz; alegre; divertido}
  \definition{v.}{expor ao sol | levantar"-se;  ascender | jogar; brincar}
\end{EntryWithPhonetic}

\begin{EntryWithPhonetic}{熙来攘往}{xi1lai2-rang3wang3}{14,7,20,8}{⽕,⽊,⼿,⼻}
  \definition{expr.}{a agitação de grandes multidões; o vai e vem em meio a multidões; repleto de atividade; pessoas circulando freneticamente}
  \synonymref{车水马龙}{che1shui3ma3long2}
  \synonymref{川流不息}{chuan1liu2-bu4xi1}
  \synonymref{络绎不绝}{luo4yi4-bu4jue2}
  \synonymref{熙熙攘攘}{xi1xi1-rang3rang3}
\end{EntryWithPhonetic}

\begin{EntryWithPhonetic}{熙熙攘攘}{xi1xi1-rang3rang3}{14,14,20,20}{⽕,⽕,⼿,⼿}[HSK 7-9]
  \definition{expr.}{``Repleto de atividade.''; pessoas circulando freneticamente}
  \seealsoref{熙来攘往}{xi1lai2-rang3wang3}
  \synonymref{车水马龙}{che1shui3ma3long2}
  \synonymref{熙来攘往}{xi1lai2-rang3wang3}
\end{EntryWithPhonetic}

%%%%%%%%%% 蜥 %%%%%%%%%%
\subsection*{蜥}\addcontentsline{loh}{figure}{蜥 \dpy{xi1}}

\begin{EntryWithPhonetic}{蜥}{xi1}{14}{⾍}
  \definition{s.}{lagarto}
\end{EntryWithPhonetic}

\begin{EntryWithPhonetic}{蜥易}{xi1yi4}{14,8}{⾍,⽇}
  \variantof{蜥蜴}
\end{EntryWithPhonetic}

\begin{EntryWithPhonetic}{蜥蜴}{xi1yi4}{14,14}{⾍,⾍}
  \definition{s.}{lagarto (\emph{Lacertídeo})}
\end{EntryWithPhonetic}

%%%%%%%%%% 嘻 %%%%%%%%%%
\subsection*{嘻}\addcontentsline{loh}{figure}{嘻 \dpy{xi1}}

\begin{EntryWithPhonetic}{嘻}{xi1}{15}{⼝}
  \definition{interj.}{Literário: exclamação de admiração, espanto, etc.}
  \definition{s.}{Onomatopéia: o som de risos alegres; descreve o riso}
\end{EntryWithPhonetic}

\begin{EntryWithPhonetic}{嘻笑}{xi1xiao4}{15,10}{⼝,⽵}
  \variantof{嬉笑}
\end{EntryWithPhonetic}

%%%%%%%%%% 嬉 %%%%%%%%%%
\subsection*{嬉}\addcontentsline{loh}{figure}{嬉 \dpy{xi1}}

\begin{EntryWithPhonetic}{嬉}{xi1}{15}{⼥}
  \definition{v.}{jogar; brincar; praticar esporte}
\end{EntryWithPhonetic}

\begin{EntryWithPhonetic}{嬉笑}{xi1xiao4}{15,10}{⼥,⽵}[HSK 7-9]
  \definition{v.}{rir e brincar | estar rindo e brincando}
  \synonymref{嘲笑}{chao2xiao4}
  \antonymref{怒骂}{nu4ma4}
\end{EntryWithPhonetic}

%%%%%%%%%% 膝 %%%%%%%%%%
\subsection*{膝}\addcontentsline{loh}{figure}{膝 \dpy{xi1}}

\begin{EntryWithPhonetic}{膝}{xi1}{15}{⾁}
  \definition[个]{s.}{joelho | genu; a parte frontal da articulação que liga a coxa à parte inferior da perna}
\end{EntryWithPhonetic}

\begin{EntryWithPhonetic}{膝盖}{xi1gai4}{15,11}{⾁,⽫}[HSK 7-9]
  \definition[个,对,双]{s.}{joelho; a parte frontal da perna que pode dobrar}
\end{EntryWithPhonetic}

%%%%%%%%%% 歙 %%%%%%%%%%
\subsection*{歙}\addcontentsline{loh}{figure}{歙 \dpy{xi1}}

\begin{EntryWithPhonetic}{歙}{xi1}{16}{⽋}
  \definition{s.}{inalação}
  \definition{v.}{inalar}
  \seeref{she4}
\end{EntryWithPhonetic}

%%%%%%%%%% 习 %%%%%%%%%%
\subsection*{习}\addcontentsline{loh}{figure}{习 \dpy{xi2}}

\begin{EntryWithPhonetic}{习}{xi2}{3}{⼄}
  \definition*{s.}{Sobrenome: Xi}
  \definition{s.}{hábito; costume; prática usual; um comportamento que se desenvolve inconscientemente por meio de ações repetidas ao longo de um longo período de tempo}
  \definition{v.}{revisar; praticar; exercitar | acostumado a; familiarizado com; familiarizado com algo por meio de contato frequente | estudar; aprender (pássaro)}
\end{EntryWithPhonetic}

\begin{EntryWithPhonetic}{习惯}{xi2guan4}{3,11}{⼄,⼼}[HSK 2]
  \definition[个,种]{s.}{hábito; costume; prática usual; comportamentos, tendências ou tendências sociais que se desenvolvem gradualmente ao longo de um longo período de tempo e são difíceis de mudar}
  \definition{v.}{estar acostumado a; ter o hábito de}
  \seealsoref{惯}{guan4}
  \synonymref{风气}{feng1qi4}
  \synonymref{风俗}{feng1su2}
  \synonymref{惯}{guan4}
  \synonymref{民俗}{min2su2}
  \synonymref{习俗}{xi2su2}
\end{EntryWithPhonetic}

\begin{EntryWithPhonetic}{习尚}{xi2shang4}{3,8}{⼄,⼩}
  \definition{s.}{prática comum; costume; prática usual}
  \synonymref{风气}{feng1qi4}
  \synonymref{风尚}{feng1shang4}
\end{EntryWithPhonetic}

\begin{EntryWithPhonetic}{习俗}{xi2su2}{3,9}{⼄,⼈}[HSK 7-9]
  \definition[种,个]{s.}{costume; convenção; hábitos e costumes}
  \synonymref{风气}{feng1qi4}
  \synonymref{风俗}{feng1su2}
  \synonymref{民俗}{min2su2}
  \synonymref{习惯}{xi2guan4}
\end{EntryWithPhonetic}

%%%%%%%%%% 席 %%%%%%%%%%
\subsection*{席}\addcontentsline{loh}{figure}{席 \dpy{xi2}}

\begin{EntryWithPhonetic}{席}{xi2}{10}{⼱}[HSK 7-9]
  \definition*{s.}{Sobrenome: Xi}
  \definition[卷,张]{s.}{esteira | assento; lugar; caixa | assento (em uma assembleia legislativa) | festim; banquete; jantar}
\end{EntryWithPhonetic}

\begin{EntryWithPhonetic}{席卷}{xi2juan3}{10,8}{⼱,⼙}
  \definition{v.}{enrolar como um tapete; levar tudo embora; remover tudo; enrolar como um tapete para embrulhar tudo}
  \synonymref{遍布}{bian4bu4}
\end{EntryWithPhonetic}

\begin{EntryWithPhonetic}{席位}{xi2wei4}{10,7}{⼱,⼈}[HSK 7-9]
  \definition[个]{s.}{assento (em uma conferência, em uma assembleia legislativa, etc.)}
\end{EntryWithPhonetic}

%%%%%%%%%% 袭 %%%%%%%%%%
\subsection*{袭}\addcontentsline{loh}{figure}{袭 \dpy{xi2}}

\begin{EntryWithPhonetic}{袭}{xi2}{11}{⾐}
  \definition*{s.}{Sobrenome: Xi}
  \definition{clas.}{usado para conjuntos completos de roupas}
  \definition{v.}{fazer um ataque surpresa a; invadir | seguir o padrão de; continuar como antes; fazer o mesmo}
\end{EntryWithPhonetic}

\begin{EntryWithPhonetic}{袭击}{xi2ji1}{11,5}{⾐,⼐}[HSK 7-9]
  \definition{v.}{invadir; atacar de surpresa; fazer um ataque surpresa contra; em termos militares, isso se refere ao seu ataque surpresa | atingir; golpear; geralmente se refere a um ataque repentino}
  \synonymref{报复}{bao4fu4}
  \synonymref{冲击}{chong1ji1}
  \synonymref{挫折}{cuo4zhe2}
  \synonymref{打击}{da3ji1}
  \synonymref{攻击}{gong1ji1}
  \synonymref{进攻}{jin4gong1}
  \synonymref{抨击}{peng1ji1}
  \synonymref{障碍}{zhang4'ai4}
  \antonymref{保卫}{bao3wei4}
\end{EntryWithPhonetic}

%%%%%%%%%% 媳 %%%%%%%%%%
\subsection*{媳}\addcontentsline{loh}{figure}{媳 \dpy{xi2}}

\begin{EntryWithPhonetic}{媳}{xi2}{13}{⼥}
  \definition[个,位,名,些]{s.}{nora}
\end{EntryWithPhonetic}

\begin{EntryWithPhonetic}{媳妇}{xi2fu4}{13,6}{⼥,⼥}
  \definition[个]{s.}{esposa do filho; nora | cunhada; nora; a esposa de um irmão mais novo ou de um parente mais novo}
  \seeref{xi2fu5}
\end{EntryWithPhonetic}

\begin{EntryWithPhonetic}{媳妇}{xi2fu5}{13,6}{⼥,⼥}[HSK 7-9]
  \definition[个]{s.}{esposa; cônjuge}
  \seeref{xi2fu4}
\end{EntryWithPhonetic}

%%%%%%%%%% 洗 %%%%%%%%%%
\subsection*{洗}\addcontentsline{loh}{figure}{洗 \dpy{xi3}}

\begin{EntryWithPhonetic}{洗}{xi3}{9}{⽔}[HSK 1]
  \definition[个]{s.}{pequeno recipiente contendo água para enxaguar os pincéis de escrever | batismo}
  \definition{v.}{lavar; tomar banho; remover a sujeira do objeto com água ou outro solvente | batizar | eliminar; corrigir; reparar | saquear; matar e pilhar; matar ou roubar tudo, como se tivesse sido lavado | revelar filmes; imprimir fotos | apagar; limpar (uma gravação, etc.) | embaralhar (cartas, etc.)}
\end{EntryWithPhonetic}

\begin{EntryWithPhonetic}{洗涤}{xi3di2}{9,10}{⽔,⽔}
  \definition{s.}{lavar; enxaguar; limpar; lavar"-se, também metaforicamente, significa usar a própria alma para se tornar limpo e puro}
  \definition{v.}{enxaguar | lavar}
  \synonymref{净化}{jing4hua4}
  \synonymref{清洗}{qing1xi3}
  \antonymref{污秽}{wu1hui4}
  \antonymref{污染}{wu1ran3}
\end{EntryWithPhonetic}

\begin{EntryWithPhonetic}{洗涤剂}{xi3di2ji4}{9,10,8}{⽔,⽔,⼑}[HSK 7-9]
  \definition[种,瓶]{s.}{detergente; líquido para lavar roupa | agente de limpeza}
\end{EntryWithPhonetic}

\begin{EntryWithPhonetic}{洗涤间}{xi3di2 jian1}{9,10,7}{⽔,⽔,⾨}
  \definition{s.}{lavanderia}
\end{EntryWithPhonetic}

\begin{EntryWithPhonetic}{洗发皂}{xi3fa4 zao4}{9,5,7}{⽔,⼜,⽩}
  \definition{s.}{xampu}
\end{EntryWithPhonetic}

\begin{EntryWithPhonetic}{洗劫}{xi3jie2}{9,7}{⽔,⼒}
  \definition{v.}{saquear; pilhar; roubar}
\end{EntryWithPhonetic}

\begin{EntryWithPhonetic}{洗净}{xi3jing4}{9,8}{⽔,⼎}
  \definition{v.}{lavar até ficar limpo; limpar}
  \synonymref{干净}{gan1jing4}
  \synonymref{洁净}{jie2jing4}
  \synonymref{净化}{jing4hua4}
  \synonymref{清洗}{qing1xi3}
  \synonymref{洗涤}{xi3di2}
  \antonymref{污染}{wu1ran3}
\end{EntryWithPhonetic}

\begin{EntryWithPhonetic}{洗礼}{xi3li3}{9,5}{⽔,⽰}[HSK 7-9]
  \definition[次,场]{s.}{batismo; uma cerimônia religiosa realizada quando um cristão se converte ao cristianismo, na qual água é aspergida na testa da pessoa que está sendo batizada, ou seu corpo é imerso em água, para simbolizar a purificação dos pecados passados | teste severo; uma metáfora para treinamento e testes rigorosos}
  \definition{v.}{batizar}
  \antonymref{堕落}{duo4luo4}
\end{EntryWithPhonetic}

\begin{EntryWithPhonetic}{洗手}{xi3shou3}{9,4}{⽔,⼿}
  \definition{v.}{lavar as mãos; eufemismo, referente ao ato de ir ao banheiro | parar de praticar o mal e reformar"-se; refere"-se metaforicamente a ladrões, jogadores ou outros que se reformam e recomeçam a vida | desistir; lavar as mãos de algo; uma metáfora para o abandono de determinada profissão}
\end{EntryWithPhonetic}

\begin{EntryWithPhonetic}{洗手不干}{xi3shou3bu2gan4}{9,4,4,3}{⽔,⼿,⼀,⼲}
  \definition{v.}{(um ladrão, etc.) pare de fazer o mal e reformar"-se | lavar as mãos de algo | reformar os próprios hábitos | parar completamente de fazer algo}
  \synonymref{改邪归正}{gai3xie2-gui1zheng4}
\end{EntryWithPhonetic}

\begin{EntryWithPhonetic}{洗手池}{xi3shou3chi2}{9,4,6}{⽔,⼿,⽔}
  \definition{s.}{pia de banheiro | lavatório}
  \seealsoref{洗手盆}{xi3shou3pen2}
\end{EntryWithPhonetic}

\begin{EntryWithPhonetic}{洗手间}{xi3shou3jian1}{9,4,7}{⽔,⼿,⾨}[HSK 1]
  \definition[个,间,处,排]{s.}{banheiro; lavatório; lavabo}
  \synonymref{厕所}{ce4suo3}
  \synonymref{卫生间}{wei4sheng1jian1}
\end{EntryWithPhonetic}

\begin{EntryWithPhonetic}{洗手盆}{xi3shou3pen2}{9,4,9}{⽔,⼿,⽫}
  \definition{s.}{pia de banheiro | lavatório}
  \seealsoref{洗手池}{xi3shou3chi2}
\end{EntryWithPhonetic}

\begin{EntryWithPhonetic}{洗手乳}{xi3shou3 ru3}{9,4,8}{⽔,⼿,⼄}
  \definition{s.}{sabonete líquido para lavar as mãos}
  \seealsoref{洗手液}{xi3shou3ye4}
\end{EntryWithPhonetic}

\begin{EntryWithPhonetic}{洗手液}{xi3shou3ye4}{9,4,11}{⽔,⼿,⽔}
  \definition{s.}{sabonete líquido para as mãos | sabonete líquido}
  \seealsoref{洗手乳}{xi3shou3 ru3}
\end{EntryWithPhonetic}

\begin{EntryWithPhonetic}{洗脱}{xi3tuo1}{9,11}{⽔,⾁}
  \definition{v.}{expurgar | vindicar; exonerar; absolver | purificar | lavar}
\end{EntryWithPhonetic}

\begin{EntryWithPhonetic}{洗碗}{xi3wan3}{9,13}{⽔,⽯}
  \definition{v.}{lavar pratos}
\end{EntryWithPhonetic}

\begin{EntryWithPhonetic}{洗胃}{xi3wei4}{9,9}{⽔,⾁}
  \definition{s.}{Medicina: lavagem gástrica}
  \definition{v.}{ter o estômago lavado}
\end{EntryWithPhonetic}

\begin{EntryWithPhonetic}{洗衣粉}{xi3yi1fen3}{9,6,10}{⽔,⾐,⽶}[HSK 6]
  \definition[袋,包,勺]{s.}{sabão em pó; detergente para roupa (em pó); detergente em pó sintetizado quimicamente, específico para uso em lavanderia}
\end{EntryWithPhonetic}

\begin{EntryWithPhonetic}{洗衣机}{xi3yi1ji1}{9,6,6}{⽔,⾐,⽊}[HSK 2]
  \definition[台,款]{s.}{máquina de lavar roupa; eletrodomésticos para lavagem automática ou semiautomática de roupas}
\end{EntryWithPhonetic}

\begin{EntryWithPhonetic}{洗澡}{xi3/zao3}{9,16}{⽔,⽔}[HSK 2]
  \definition{v.+compl.}{tomar banho; tomar banho de chuveiro; lavar"-se}
\end{EntryWithPhonetic}

\begin{EntryWithPhonetic}{洗澡间}{xi3zao3jian1}{9,16,7}{⽔,⽔,⾨}
  \definition[间]{s.}{banheiro}
\end{EntryWithPhonetic}

%%%%%%%%%% 喜 %%%%%%%%%%
\subsection*{喜}\addcontentsline{loh}{figure}{喜 \dpy{xi3}}

\begin{EntryWithPhonetic}{喜}{xi3}{12}{⼝}
  \definition{adj.}{feliz; satisfeito; encantado}
  \definition[桩,件]{s.}{evento feliz (especialmente casamento); ocasião para celebração; algo para comemorar | gravidez | casamento ou coisas relacionadas a ele}
  \definition{v.}{gostar; fonte de; ter inclinação para | precisa; requer; combina melhor com; (um certo organismo) precisa ou é adequado para (um certo ambiente ou algo)}
\end{EntryWithPhonetic}

\begin{EntryWithPhonetic}{喜爱}{xi3'ai4}{12,10}{⼝,⽖}[HSK 4]
  \definition{v.}{gostar; amar; ter afeição por; estar interessado em; ter uma queda ou sentir interesse por pessoas ou coisas}
\end{EntryWithPhonetic}

\begin{EntryWithPhonetic}{喜出望外}{xi3chu1wang4wai4}{12,5,11,5}{⼝,⼐,⽉,⼣}[HSK 7-9]
  \definition{expr.}{extremamente feliz com algo; agradavelmente surpreso; ficar satisfeito além das expectativas; ficar especialmente feliz ao receber boas notícias inesperadas}
\end{EntryWithPhonetic}

\begin{EntryWithPhonetic}{喜好}{xi3hao4}{12,6}{⼝,⼥}[HSK 7-9]
  \definition[个]{s.}{\emph{hobby}; interesse; algo que eu realmente gosto; um \emph{hobby}}
  \definition{v.}{gostar; amar; ter afeição por; ser apaixonado por; estar interessado em algo}
  \synonymref{爱好}{ai4hao4}
  \synonymref{宠爱}{chong3'ai4}
  \synonymref{嗜好}{shi4hao4}
  \synonymref{特长}{te4chang2}
  \synonymref{喜爱}{xi3'ai4}
  \synonymref{喜欢}{xi3huan5}
  \antonymref{讨厌}{tao3/yan4}
\end{EntryWithPhonetic}

\begin{EntryWithPhonetic}{喜欢}{xi3huan5}{12,6}{⼝,⽋}[HSK 1]
  \definition{adj.}{feliz; encantado; exultante; cheio de alegria}
  \definition{v.}{gostar; amar; ter afeição por; estar interessado em; ter uma boa impressão ou interesse por alguém ou algo}
  \seealsoref{爱}{ai4}
  \seealsoref{爱好}{ai4hao4}
  \seealsoref{喜爱}{xi3'ai4}
  \seealsoref{欣赏}{xin1shang3}
  \synonymref{爱好}{ai4hao4}
  \synonymref{宠爱}{chong3'ai4}
  \synonymref{高兴}{gao1xing4}
  \synonymref{快乐}{kuai4le4}
  \synonymref{热爱}{re4'ai4}
  \synonymref{热衷}{re4zhong1}
  \synonymref{嗜好}{shi4hao4}
  \synonymref{喜爱}{xi3'ai4}
  \synonymref{喜好}{xi3hao4}
  \antonymref{仇恨}{chou2hen4}
  \antonymref{讨厌}{tao3/yan4}
  \antonymref{厌烦}{yan4fan2}
\end{EntryWithPhonetic}

\begin{EntryWithPhonetic}{喜酒}{xi3jiu3}{12,10}{⼝,⾣}[HSK 7-9]
  \definition[杯,顿]{s.}{banquete de casamento; bebidas alcoólicas em um casamento; o vinho ou banquete servido a parentes e amigos em um casamento}
\end{EntryWithPhonetic}

\begin{EntryWithPhonetic}{喜剧}{xi3ju4}{12,10}{⼝,⼑}[HSK 5]
  \definition[部,出]{s.}{comédia | comédia; uma das principais categorias do teatro; usa o exagero para satirizar e ridicularizar o feio; fenômenos retrógrados; destaca as contradições inerentes a esses fenômenos e seu conflito com coisas saudáveis; costuma provocar risadas; o final geralmente é feliz}
  \antonymref{悲剧}{bei1ju4}
\end{EntryWithPhonetic}

\begin{EntryWithPhonetic}{喜怒哀乐}{xi3-nu4-ai1-le4}{12,9,9,5}{⼝,⼼,⼝,⼃}[HSK 7-9]
  \definition{expr.}{``Alegria, raiva, tristeza e felicidade.''; as quatro expressões básicas de emoção incluem felicidade, raiva, tristeza e alegria}
\end{EntryWithPhonetic}

\begin{EntryWithPhonetic}{喜庆}{xi3qing4}{12,6}{⼝,⼴}[HSK 7-9]
  \definition{adj.}{feliz; alegre; jubilante; digno de elogios e celebração}
  \definition{s.}{um evento feliz; uma ocasião para celebrar; coisas que valem a pena curtir e celebrar}
  \definition{v.}{celebrar com alegria}
  \synonymref{庆贺}{qing4he4}
  \synonymref{庆祝}{qing4zhu4}
  \antonymref{凄凉}{qi1liang2}
\end{EntryWithPhonetic}

\begin{EntryWithPhonetic}{喜事}{xi3shi4}{12,8}{⼝,⼅}[HSK 7-9]
  \definition[桩]{s.}{evento feliz; ocasião alegre; algo para celebrar e se alegrar | casamento; matrimônio; refere"-se especificamente a casamentos}
\end{EntryWithPhonetic}

\begin{EntryWithPhonetic}{喜糖}{xi3tang2}{12,16}{⼝,⽶}[HSK 7-9]
  \definition[颗,包]{s.}{doces de casamento; doces servidos a parentes e amigos em um casamento}
\end{EntryWithPhonetic}

\begin{EntryWithPhonetic}{喜洋洋}{xi3yang2yang2}{12,9,9}{⼝,⽔,⽔}[HSK 7-9]
  \definition{adj.}{transbordando de alegria; radiante; muito feliz | radiante de alegria; descrevendo uma aparência muito alegre}
  \synonymref{呼啦啦}{hu1la1la1}
\end{EntryWithPhonetic}

\begin{EntryWithPhonetic}{喜悦}{xi3yue4}{12,10}{⼝,⼼}[HSK 7-9]
  \definition{adj.}{feliz; alegre; jubiloso; encantador; agradável}
  \synonymref{得意}{de2/yi4}
  \synonymref{高兴}{gao1xing4}
  \synonymref{欢乐}{huan1le4}
  \synonymref{欢喜}{huan1xi3}
  \synonymref{开心}{kai1/xin1}
  \synonymref{快乐}{kuai4le4}
  \synonymref{乐意}{le4yi4}
  \synonymref{兴奋}{xing1fen4}
  \synonymref{愉快}{yu2kuai4}
  \antonymref{悲哀}{bei1'ai1}
  \antonymref{悲伤}{bei1shang1}
  \antonymref{悲痛}{bei1tong4}
  \antonymref{伤心}{shang1/xin1}
  \antonymref{痛苦}{tong4ku3}
  \antonymref{辛酸}{xin1suan1}
  \antonymref{忧愁}{you1chou2}
\end{EntryWithPhonetic}

%%%%%%%%%% 戏 %%%%%%%%%%
\subsection*{戏}\addcontentsline{loh}{figure}{戏 \dpy{xi4}}

\begin{EntryWithPhonetic}{戏}{xi4}{6}{⼽}[HSK 5]
  \definition*{s.}{Sobrenome: Xi}
  \definition[场,部,出,台]{s.}{drama; peça; espetáculo; \emph{show}}
  \definition{v.}{brincar; praticar esportes; jogar | zombar; brincar; provocar}
\end{EntryWithPhonetic}

\begin{EntryWithPhonetic}{戏法}{xi4fa3}{6,8}{⼽,⽔}
  \definition{s.}{magia; ilusionismo; prestidigitação; truques | malabarismo}
  \synonymref{把戏}{ba3xi4}
  \synonymref{魔术}{mo2shu4}
\end{EntryWithPhonetic}

\begin{EntryWithPhonetic}{戏剧}{xi4ju4}{6,10}{⼽,⼑}[HSK 5]
  \definition[出,部]{s.}{drama; peça; teatro | roteiro; peça; cenário}
\end{EntryWithPhonetic}

\begin{EntryWithPhonetic}{戏剧般}{xi4ju4 ban1}{6,10,10}{⼽,⼑,⾈}
  \definition{adj.}{melodramático}
\end{EntryWithPhonetic}

\begin{EntryWithPhonetic}{戏剧编剧}{xi4ju4 bian1ju4}{6,10,12,10}{⼽,⼑,⽷,⼑}
  \definition{s.}{dramaturgo}
\end{EntryWithPhonetic}

\begin{EntryWithPhonetic}{戏剧化地}{xi4ju4hua4 di4}{6,10,4,6}{⼽,⼑,⼔,⼟}
  \definition{adv.}{teatralmente; dramaticamente}
\end{EntryWithPhonetic}

\begin{EntryWithPhonetic}{戏剧家}{xi4ju4jia1}{6,10,10}{⼽,⼑,⼧}
  \definition{s.}{dramaturgo; escritor}
\end{EntryWithPhonetic}

\begin{EntryWithPhonetic}{戏剧效果}{xi4ju4 xiao4guo3}{6,10,10,8}{⼽,⼑,⽁,⽊}
  \definition{s.}{efeito dramático}
\end{EntryWithPhonetic}

\begin{EntryWithPhonetic}{戏剧性}{xi4ju4xing4}{6,10,8}{⼽,⼑,⼼}
  \definition{adj.}{dramático; espetacular; inesperado; repentino}
  \synonymref{巧合}{qiao3he2}
\end{EntryWithPhonetic}

\begin{EntryWithPhonetic}{戏剧演出}{xi4ju4 yan3chu1}{6,10,14,5}{⼽,⼑,⽔,⼐}
  \definition{s.}{performance dramática; apresentação teatral}
\end{EntryWithPhonetic}

\begin{EntryWithPhonetic}{戏弄}{xi4nong4}{6,7}{⼽,⼶}
  \definition{v.}{brincar; provocar; zombar de; pregar peças em; fazer brincadeiras e provocar; zombar dos outros}
  \synonymref{嘲弄}{chao2nong4}
  \synonymref{嘲笑}{chao2xiao4}
  \synonymref{讥笑}{ji1xiao4}
  \synonymref{调侃}{tiao2kan3}
  \synonymref{戏谑}{xi4xue4}
  \antonymref{敬重}{jing4zhong4}
  \antonymref{尊敬}{zun1jing4}
  \antonymref{尊重}{zun1zhong4}
\end{EntryWithPhonetic}

\begin{EntryWithPhonetic}{戏曲}{xi4qu3}{6,6}{⼽,⽈}[HSK 6]
  \definition{s.}{drama; ópera chinesa; ópera tradicional; forma teatral tradicional | partes cantadas em 传奇 e zaju 杂剧}
  \seealsoref{传奇}{chuan2qi2}
  \seealsoref{杂剧}{za2ju4}
\end{EntryWithPhonetic}

\begin{EntryWithPhonetic}{戏耍}{xi4shua3}{6,9}{⼽,⽽}
  \definition{v.}{provocar; pregar peças em | divertir"-se; brincar com}
  \synonymref{嘲弄}{chao2nong4}
  \synonymref{游戏}{you2xi4}
\end{EntryWithPhonetic}

\begin{EntryWithPhonetic}{戏谑}{xi4xue4}{6,11}{⼽,⾔}
  \definition{v.}{provocar; zombar; fazer piadas com comentários engraçados e divertidos}
  \synonymref{戏弄}{xi4nong4}
  \antonymref{严肃}{yan2su4}
\end{EntryWithPhonetic}

\begin{EntryWithPhonetic}{戏院}{xi4yuan4}{6,9}{⼽,⾩}
  \definition{s.}{teatro; cinema}
  \synonymref{剧场}{ju4chang3}
\end{EntryWithPhonetic}

%%%%%%%%%% 系 %%%%%%%%%%
\subsection*{系}\addcontentsline{loh}{figure}{系 \dpy{xi4}}

\begin{EntryWithPhonetic}{系}{xi4}{7}{⽷}[HSK 3]
  \definition*{s.}{Sobrenome: Xi}
  \definition{s.}{sistema; série | departamento; faculdade; unidades administrativas de ensino divididas por disciplina nas instituições de ensino superior}
  \definition{v.}{relacionar"-se com; suportar; depender de | sentir"-se ansioso; estar preocupado | amarrar; prender | ser; expressa julgamento, equivalente a 是}
  \seeref{ji4}
  \seealsoref{是}{shi4}
\end{EntryWithPhonetic}

\begin{EntryWithPhonetic}{系列}{xi4lie4}{7,6}{⽷,⼑}[HSK 4]
  \definition{s.}{série; conjunto; conjunto de coisas relacionadas (matemática)}
\end{EntryWithPhonetic}

\begin{EntryWithPhonetic}{系囚}{xi4qiu2}{7,5}{⽷,⼞}
  \definition{s.}{prisioneiro}
\end{EntryWithPhonetic}

\begin{EntryWithPhonetic}{系统}{xi4tong3}{7,9}{⽷,⽷}[HSK 4]
  \definition{adj.}{sistemático; organizado}
  \definition[套,个]{s.}{sistema; relação de tipos semelhantes (ou seja, grupo de coisas semelhantes)}
\end{EntryWithPhonetic}

%%%%%%%%%% 细 %%%%%%%%%%
\subsection*{细}\addcontentsline{loh}{figure}{细 \dpy{xi4}}

\begin{EntryWithPhonetic}{细}{xi4}{8}{⽷}[HSK 4]
  \definition{adj.}{fino; delgado; esguio; esbelto | fino; em partículas pequenas; grãos pequenos | fino e macio;  um sussuro | fino; requintado; delicado | cuidadoso; detalhado; meticuloso | ínfimo; minúsculo; insignificante; diminuto | jovem; pequeno}
  \antonymref{粗}{cu1}
\end{EntryWithPhonetic}

\begin{EntryWithPhonetic}{细胞}{xi4bao1}{8,9}{⽷,⾁}[HSK 6]
  \definition[个]{s.}{célula; unidade estrutural e funcional básica de um organismo, com uma variedade de formas, composta principalmente pelo núcleo, citoplasma e membrana celular; as plantas também possuem paredes celulares fora da membrana celular}
\end{EntryWithPhonetic}

\begin{EntryWithPhonetic}{细节}{xi4jie2}{8,5}{⽷,⾋}[HSK 4]
  \definition[处]{s.}{detalhe; particularidade; aspectos secundários ou partes sutis de um enredo ou episódios secundários usados em uma obra literária para expressar o caráter de uma pessoa ou as características essenciais de uma coisa}
\end{EntryWithPhonetic}

\begin{EntryWithPhonetic}{细菌}{xi4jun1}{8,11}{⽷,⾋}[HSK 6]
  \definition[个]{s.}{germe; bactéria; um organismo muito pequeno, invisível aos olhos humanos}
\end{EntryWithPhonetic}

\begin{EntryWithPhonetic}{细菌战}{xi4jun1zhan4}{8,11,9}{⽷,⾋,⼽}
  \definition{s.}{guerra bacteriológica | guerra biológica}
\end{EntryWithPhonetic}

\begin{EntryWithPhonetic}{细腻}{xi4ni4}{8,13}{⽷,⾁}[HSK 7-9]
  \definition{adj.}{exótico; delicado; fino e suave; suave e delicado | pormenorizado; soberbo; requintado; (descrição, desempenho, etc.) detalhado e meticuloso}
  \synonymref{精细}{jing1xi4}
  \synonymref{精致}{jing1zhi4}
  \synonymref{细致}{xi4zhi4}
  \antonymref{粗糙}{cu1cao1}
  \antonymref{粗鲁}{cu1lu3}
\end{EntryWithPhonetic}

\begin{EntryWithPhonetic}{细微}{xi4wei1}{8,13}{⽷,⼻}[HSK 7-9]
  \definition{adj.}{leve; minúsculo; sutil; minuto}
  \synonymref{渺小}{miao3xiao3}
  \synonymref{轻微}{qing1wei1}
  \synonymref{微小}{wei1xiao3}
  \synonymref{细节}{xi4jie2}
  \antonymref{巨大}{ju4da4}
  \antonymref{庞大}{pang2da4}
  \antonymref{显著}{xian3zhu4}
  \antonymref{重大}{zhong4da4}
\end{EntryWithPhonetic}

\begin{EntryWithPhonetic}{细心}{xi4xin1}{8,4}{⽷,⼼}[HSK 7-9]
  \definition{adj.}{cuidadoso; atento; meticuloso}
  \seealsoref{认真}{ren4zhen1}
  \seealsoref{细致}{xi4zhi4}
  \synonymref{谨慎}{jin3shen4}
  \synonymref{精心}{jing1xin1}
  \synonymref{留心}{liu2/xin1}
  \synonymref{留意}{liu2/yi4}
  \synonymref{小心}{xiao3xin5}
  \synonymref{用心}{yong4/xin1}
  \synonymref{注意}{zhu4/yi4}
  \synonymref{仔细}{zi3xi4}
  \antonymref{粗心}{cu1xin1}
  \antonymref{大意}{da4yi5}
  \antonymref{含糊}{han2hu5}
  \antonymref{鲁莽}{lu3mang3}
  \antonymref{马虎}{ma3hu5}
\end{EntryWithPhonetic}

\begin{EntryWithPhonetic}{细致}{xi4zhi4}{8,10}{⽷,⾄}[HSK 4]
  \definition{adj.}{meticuloso; cuidadoso; minucioso | intrincado; delicado}
\end{EntryWithPhonetic}

%%%%%%%%%% 虾 %%%%%%%%%%
\subsection*{虾}\addcontentsline{loh}{figure}{虾 \dpy{xia1}}

\begin{EntryWithPhonetic}{虾}{xia1}{9}{⾍}[HSK 7-9]
  \definition[只,对]{s.}{camarão}
  \seeref{ha2}
\end{EntryWithPhonetic}

%%%%%%%%%% 瞎 %%%%%%%%%%
\subsection*{瞎}\addcontentsline{loh}{figure}{瞎 \dpy{xia1}}

\begin{EntryWithPhonetic}{瞎}{xia1}{15}{⽬}[HSK 7-9]
  \definition{adv.}{sem fundamento; tolamente; sem propósito}
  \definition{v.}{cegar | Dialeto: desperdiçar; estragar; perder | (bala, bomba, etc.) não explodir; não detonar | Dialeto: (fios, etc.) emaranhar"-se | (semente/grão) não germinar; não emergir do solo}
\end{EntryWithPhonetic}

%%%%%%%%%% 峡 %%%%%%%%%%
\subsection*{峡}\addcontentsline{loh}{figure}{峡 \dpy{xia2}}

\begin{EntryWithPhonetic}{峡}{xia2}{9}{⼭}
  \definition[个]{s.}{(frequentemente como parte do nome de um lugar) desfiladeiro | estreito | desfiladeiro; um lugar entre duas montanhas e um rio}
\end{EntryWithPhonetic}

\begin{EntryWithPhonetic}{峡谷}{xia2gu3}{9,7}{⼭,⾕}[HSK 7-9]
  \definition[段,条]{s.}{desfiladeiro; cânion; o rio atravessa um vale profundo e estreito, com penhascos íngremes em ambos os lados}
  \synonymref{山谷}{shan1gu3}
\end{EntryWithPhonetic}

%%%%%%%%%% 狭 %%%%%%%%%%
\subsection*{狭}\addcontentsline{loh}{figure}{狭 \dpy{xia2}}

\begin{EntryWithPhonetic}{狭}{xia2}{9}{⽝}
  \definition{adj.}{estreito}
  \antonymref{广}{guang3}
\end{EntryWithPhonetic}

\begin{EntryWithPhonetic}{狭隘}{xia2'ai4}{9,12}{⽝,⾩}[HSK 7-9]
  \definition{adj.}{estreito; largura pequena | de mente fechada; limitado e restrito; provinciano}
  \synonymref{狭小}{xia2xiao3}
  \synonymref{狭窄}{xia2zhai3}
  \antonymref{广阔}{guang3kuo4}
  \antonymref{宏大}{hong2da4}
  \antonymref{开阔}{kai1kuo4}
  \antonymref{宽广}{kuan1guang3}
  \antonymref{宽阔}{kuan1kuo4}
  \antonymref{宽容}{kuan1rong2}
  \antonymref{辽阔}{liao2kuo4}
\end{EntryWithPhonetic}

\begin{EntryWithPhonetic}{狭小}{xia2xiao3}{9,3}{⽝,⼩}[HSK 7-9]
  \definition{adj.}{estreito; estreito e pequeno}
  \definition{pref.}{esteno-}
  \synonymref{紧凑}{jin3cou4}
  \synonymref{狭隘}{xia2'ai4}
  \synonymref{狭窄}{xia2zhai3}
  \antonymref{广大}{guang3da4}
  \antonymref{广阔}{guang3kuo4}
  \antonymref{开阔}{kai1kuo4}
  \antonymref{宽广}{kuan1guang3}
  \antonymref{宽阔}{kuan1kuo4}
  \antonymref{宽敞}{kuan1chang5}
  \antonymref{辽阔}{liao2kuo4}
\end{EntryWithPhonetic}

\begin{EntryWithPhonetic}{狭义}{xia2yi4}{9,3}{⽝,⼂}[HSK 7-9]
  \definition{s.}{sentido restrito; uma definição mais restrita}
  \antonymref{广义}{guang3yi4}
\end{EntryWithPhonetic}

\begin{EntryWithPhonetic}{狭窄}{xia2zhai3}{9,10}{⽝,⽳}[HSK 7-9]
  \definition{adj.}{estreito; apertado; largura pequena | (mente, experiência, etc.) estreito; paroquial; de mente fechada; limitado | estreito; alcance pequeno}
  \synonymref{微小}{wei1xiao3}
  \synonymref{狭隘}{xia2'ai4}
  \synonymref{狭小}{xia2xiao3}
  \antonymref{广泛}{guang3fan4}
  \antonymref{广阔}{guang3kuo4}
  \antonymref{宏大}{hong2da4}
  \antonymref{开阔}{kai1kuo4}
  \antonymref{开朗}{kai1lang3}
  \antonymref{宽广}{kuan1guang3}
  \antonymref{宽阔}{kuan1kuo4}
  \antonymref{宽敞}{kuan1chang5}
  \antonymref{辽阔}{liao2kuo4}
\end{EntryWithPhonetic}

%%%%%%%%%% 下 %%%%%%%%%%
\subsection*{下}\addcontentsline{loh}{figure}{下 \dpy{xia4}}

\begin{EntryWithPhonetic}{下}{xia4}{3}{⼀}[HSK 1,2]
  \definition{clas.}{número de vezes usado para a ação | volume de um contêiner; quantidade de objetos que cabem em um utensílio | usado depois de 两 e 几 para expressar habilidade, capacidade, destreza}
  \definition{s.}{abaixo | próximo; último; segundo; referindo"-se ao que está por vir ou ao que vem em seguida | mais baixo; inferior; de baixo nível ou grau | próximo; último; segundo; em ordem ou em ordem cronológica | indica pertencer a uma determinada faixa, situação, condição, etc. | indica uma determinada época ou estação | usado após um número para indicar posição ou direção | para baixo (após uma preposição) | sob (depois de um substantivo) | para baixo (antes de um verbo)}
  \definition{v.}{desembarcar; descer; sair | cair (chuva, neve, etc.) | enviar; emitir; entregar | ir para | sair; partir; retirar"-se | lançar; colocar | descarregar; desmontar; tirar (fora) | formar (uma opinião, ideia, etc.); tomar decisões, fazer julgamentos, etc. | usar; aplicar | dar à luz (animais) | tomar; capturar; conquistar | ceder | terminar; deixar de lado; terminar o trabalho ou os estudos diários na hora prevista | para negação; ser inferior a; ser menor que}
  \seealsoref{几}{ji3}
  \seealsoref{两}{liang3}
  \seealsoref{下边}{xia4bian5}
  \seealsoref{下面}{xia4mian5}
  \synonymref{差}{chai4}
  \synonymref{次}{ci4}
  \synonymref{低}{di1}
  \synonymref{放}{fang4}
  \synonymref{攻}{gong1}
  \synonymref{克}{ke4}
  \synonymref{取}{qu3}
  \synonymref{少}{shao3}
  \synonymref{退}{tui4}
  \synonymref{卸}{xie4}
  \antonymref{安}{an1}
  \antonymref{多}{duo1}
  \antonymref{高}{gao1}
  \antonymref{好}{hao3}
  \antonymref{进}{jin4}
  \antonymref{良}{liang2}
  \antonymref{拿}{na2}
  \antonymref{上}{shang4}
  \antonymref{优}{you1}
  \antonymref{装}{zhuang1}
\end{EntryWithPhonetic}

\begin{EntryWithPhonetic}{下巴}{xia4ba5}{3,4}{⼀,⼰}
  \definition[个]{s.}{queixo | mandíbula inferior}
\end{EntryWithPhonetic}

\begin{EntryWithPhonetic}{下班}{xia4/ban1}{3,10}{⼀,⽟}[HSK 1]
  \definition{v.+compl.}{sair do trabalho; bater ponto; terminar o trabalho na hora prevista e sair do local de trabalho}
  \antonymref{开工}{kai1/gong1}
  \antonymref{上班}{shang4/ban1}
\end{EntryWithPhonetic}

\begin{EntryWithPhonetic}{下边}{xia4bian5}{3,5}{⼀,⾡}[HSK 1]
  \definition{s.}{abaixo; sob; por baixo | próximo em ordem; seguinte | nível inferior; subordinado | a parte inferior}
  \seealsoref{下}{xia4}
  \synonymref{底下}{di3xia5}
  \synonymref{下}{xia4}
  \synonymref{下方}{xia4fang1}
  \synonymref{下面}{xia4mian5}
  \antonymref{上边}{shang4bian5}
  \antonymref{上方}{shang4fang1}
  \antonymref{上面}{shang4mian5}
\end{EntryWithPhonetic}

\begin{EntryWithPhonetic}{下场}{xia4chang3}{3,6}{⼀,⼟}[HSK 7-9]
  \definition[个,种]{s.}{final ruim; estado lamentável; destino}
  \definition{v.}{sair de cena; encerrar a noite}
  \seealsoref{结局}{jie2ju2}
  \synonymref{结局}{jie2ju2}
  \synonymref{结束}{jie2shu4}
  \antonymref{上场}{shang4/chang3}
\end{EntryWithPhonetic}

\begin{EntryWithPhonetic}{下车}{xia4che1}{3,4}{⼀,⾞}[HSK 1]
  \definition{v.}{descer ou sair de (um ônibus, trem, carro etc.)}
  \antonymref{上车}{shang4che1}
\end{EntryWithPhonetic}

\begin{EntryWithPhonetic}{下次}{xia4ci4}{3,6}{⼀,⽋}[HSK 1]
  \definition{s.}{na próxima vez; na próxima oportunidade ou no próximo evento}
  \synonymref{再次}{zai4ci4}
\end{EntryWithPhonetic}

\begin{EntryWithPhonetic}{下蛋}{xia4dan4}{3,11}{⼀,⾍}
  \definition{v.}{botar ovos}
\end{EntryWithPhonetic}

\begin{EntryWithPhonetic}{下调}{xia4diao4}{3,10}{⼀,⾔}
  \definition{v.}{rebaixar; diminuir a regulamentação | passar para uma unidade inferior}
  \seeref{xia4tiao2}
\end{EntryWithPhonetic}

\begin{EntryWithPhonetic}{下跌}{xia4die1}{3,12}{⼀,⾜}[HSK 7-9]
  \definition{v.}{cair; despencar; descer}
  \synonymref{下降}{xia4jiang4}
  \synonymref{下落}{xia4luo4}
  \antonymref{上涨}{shang4zhang3}
\end{EntryWithPhonetic}

\begin{EntryWithPhonetic}{下方}{xia4fang1}{3,4}{⼀,⽅}
  \definition{s.}{parte inferior | abaixo | embaixo | mundo dos mortais}
  \definition{v.}{descer ao mundo dos mortais (deuses)}
  \seealsoref{上方}{shang4fang1}
  \antonymref{上方}{shang4fang1}
\end{EntryWithPhonetic}

\begin{EntryWithPhonetic}{下岗}{xia4/gang3}{3,7}{⼀,⼭}[HSK 7-9]
  \definition{v.+compl.}{ser despedido; perder o emprego; (funcionários) que deixam seus cargos devido à falência da empresa, redução de pessoal ou outros motivos | sair de serviço}
  \antonymref{全职}{quan2zhi2}
\end{EntryWithPhonetic}

\begin{EntryWithPhonetic}{下个月}{xia4ge4yue4}{3,3,4}{⼀,⼈,⽉}[HSK 4]
  \definition{s.}{próximo mês; mês que vem; refere"-se ao próximo mês do mês atual}
\end{EntryWithPhonetic}

\begin{EntryWithPhonetic}{下功夫}{xia4gong1fu5}{3,5,4}{⼀,⼒,⼤}[HSK 7-9]
  \definition{v.}{concentrar esforços; investir tempo e energia em algo para alcançar bons resultados}
\end{EntryWithPhonetic}

\begin{EntryWithPhonetic}{下海}{xia4/hai3}{3,10}{⼀,⽔}[HSK 7-9]
  \definition{v.+compl.}{ir pescar no mar; pescar para ganhar a vida; ir para o mar | tornar"-se profissional; isso se refere a atores de ópera amadores que se tornam atores profissionais | deixar o emprego original e abrir o próprio negócio; refere"-se a pessoas que não eram originalmente donas de negócios, mas que passaram a empreender}
  \antonymref{登陆}{deng1/lu4}
\end{EntryWithPhonetic}

\begin{EntryWithPhonetic}{下级}{xia4ji2}{3,6}{⼀,⽷}[HSK 7-9]
  \definition{s.}{classificação inferior | subordinado | nível inferior (posição)}
  \synonymref{部下}{bu4xia4}
  \synonymref{下属}{xia4shu3}
  \antonymref{上级}{shang4ji2}
  \antonymref{上司}{shang4si5}
\end{EntryWithPhonetic}

\begin{EntryWithPhonetic}{下降}{xia4jiang4}{3,8}{⼀,⾩}[HSK 4]
  \definition{v.}{cair; despencar; declinar; descer; diminuir; ir para baixo}
  \synonymref{低落}{di1luo4}
  \synonymref{降低}{jiang4di1}
  \synonymref{降落}{jiang4luo4}
  \synonymref{下跌}{xia4die1}
  \synonymref{下落}{xia4luo4}
  \synonymref{消沉}{xiao1chen2}
  \antonymref{回升}{hui2sheng1}
  \antonymref{爬升}{pa2sheng1}
  \antonymref{起飞}{qi3fei1}
  \antonymref{上升}{shang4sheng1}
  \antonymref{上涨}{shang4zhang3}
  \antonymref{增加}{zeng1jia1}
  \antonymref{增长}{zeng1zhang3}
\end{EntryWithPhonetic}

\begin{EntryWithPhonetic}{下决心}{xia4jue2xin1}{3,6,4}{⼀,⼎,⼼}[HSK 7-9]
  \definition{v.}{resolver; determinar; tomar uma decisão}
\end{EntryWithPhonetic}

\begin{EntryWithPhonetic}{下课}{xia4/ke4}{3,10}{⼀,⾔}[HSK 1]
  \definition{v.+compl.}{terminar a aula; sair da aula}
  \synonymref{下台}{xia4/tai2}
  \antonymref{就任}{jiu4ren4}
  \antonymref{上课}{shang4/ke4}
  \antonymref{上台}{shang4tai2}
\end{EntryWithPhonetic}

\begin{EntryWithPhonetic}{下来}{xia4lai5}{3,7}{⼀,⽊}[HSK 3]
  \definition{part.}{usado após o verbo, indica que a ação ou o comportamento se dirige para a posição do falante ou que a ação é contínua ou concluída | usado após um adjetivo, indica que uma determinada situação começou a ocorrer e continuará a se desenvolver}
  \definition{v.}{descer (para a minha localização) | (colheitas/frutas/vegetais, etc.) ser colhido; estar maduro o suficiente para ser colhido | (período de tempo) acabar; passar; chegar ao fim; indicar o fim de um período de tempo}
  \synonymref{下去}{xia4qu5}
  \antonymref{上去}{shang4qu5}
\end{EntryWithPhonetic}

\begin{EntryWithPhonetic}{下令}{xia4/ling4}{3,5}{⼀,⼈}[HSK 7-9]
  \definition{v.+compl.}{dar ordens; ordenar}
  \synonymref{命令}{ming4ling4}
\end{EntryWithPhonetic}

\begin{EntryWithPhonetic}{下楼}{xia4lou2}{3,13}{⼀,⽊}[HSK 4]
  \definition{v.}{descer as escadas}
\end{EntryWithPhonetic}

\begin{EntryWithPhonetic}{下落}{xia4luo4}{3,12}{⼀,⾋}[HSK 7-9]
  \definition{s.}{localização da pessoa ou coisa que está sendo procurada}
  \definition{v.}{cair; deixar cair}
  \synonymref{落下}{luo4xia4}
  \synonymref{下跌}{xia4die1}
  \synonymref{下降}{xia4jiang4}
  \synonymref{着落}{zhuo2luo4}
  \antonymref{爬升}{pa2sheng1}
  \antonymref{上升}{shang4sheng1}
  \antonymref{上涨}{shang4zhang3}
\end{EntryWithPhonetic}

\begin{EntryWithPhonetic}{下面}{xia4mian5}{3,9}{⼀,⾯}[HSK 3]
  \definition{s.}{em baixo; abaixo; parte de baixo | próximo; seguinte; a parte posterior; a parte posterior de um artigo ou discurso em relação ao que está sendo narrado no momento | subordinado; o nível inferior; homens nos níveis inferiores | por baixo}
  \synonymref{底下}{di3xia5}
  \synonymref{上方}{shang4fang1}
  \antonymref{上面}{shang4mian5}
\end{EntryWithPhonetic}

\begin{EntryWithPhonetic}{下期}{xia4qi1}{3,12}{⼀,⽉}[HSK 7-9]
  \definition{s.}{da próxima vez; o período ou intervalo de tempo do próximo evento ou atividade.}
\end{EntryWithPhonetic}

\begin{EntryWithPhonetic}{下棋}{xia4/qi2}{3,12}{⼀,⽊}[HSK 7-9]
  \definition{v.+compl.}{jogar xadrez; ter uma partida de xadrez}
\end{EntryWithPhonetic}

\begin{EntryWithPhonetic}{下去}{xia4qu5}{3,5}{⼀,⼛}[HSK 3]
  \definition{part.}{usado depois de verbos para indicar de alto a baixo | usado depois de um verbo para indicar continuação}
  \definition{v.}{descer; baixar (a partir da minha localização) | (após um verbo) continuar (fazendo algo); prosseguir | usado após o verbo, indica uma descida de um ponto alto para um ponto baixo | usado após o verbo, indica continuidade | usado após um adjetivo, indica que o grau continua aumentando}
  \synonymref{出来}{chu1lai5}
  \synonymref{下来}{xia4lai5}
  \antonymref{上来}{shang4lai5}
\end{EntryWithPhonetic}

\begin{EntryWithPhonetic}{下山}{xia4/shan1}{3,3}{⼀,⼭}[HSK 7-9]
  \definition{v.+compl.}{descer uma colina ou montanha | (Sol) pôr"-se; desaparecer abaixo do horizonte}
\end{EntryWithPhonetic}

\begin{EntryWithPhonetic}{下手}{xia4/shou3}{3,4}{⼀,⼿}[HSK 7-9]
  \definition{v.+compl.}{colocar a mão na massa (em uma tarefa, etc.); começar a fazer algo; empreender | fazer algo ruim a alguém}
  \synonymref{出手}{chu1/shou3}
  \synonymref{动手}{dong4/shou3}
  \synonymref{开始}{kai1shi3}
  \synonymref{入手}{ru4shou3}
  \synonymref{着手}{zhuo2shou3}
\end{EntryWithPhonetic}

\begin{EntryWithPhonetic}{下属}{xia4shu3}{3,12}{⼀,⼫}[HSK 7-9]
  \definition[个,位,名]{s.}{subordinado; subalterno}
  \synonymref{部属}{bu4shu3}
  \synonymref{部下}{bu4xia4}
  \synonymref{下级}{xia4ji2}
  \antonymref{上司}{shang4si5}
\end{EntryWithPhonetic}

\begin{EntryWithPhonetic}{下水道}{xia4shui3dao4}{3,4,12}{⼀,⽔,⾡}
  \definition{s.}{esgoto; sarjeta; ralo}
\end{EntryWithPhonetic}

\begin{EntryWithPhonetic}{下台}{xia4/tai2}{3,5}{⼀,⼝}[HSK 7-9]
  \definition{v.+compl.}{descer do palco, pódio ou da plataforma | perder o poder; deixar o cargo | (geralmente na forma negativa) sair de uma situação difícil ou embaraçosa; retirar"-se com dignidade; sair de fininho; saída elegante | renunciar; renunciar ao próprio direito (poder); perder o poder; livrar"-se de uma posição}
  \antonymref{上台}{shang4tai2}
\end{EntryWithPhonetic}

\begin{EntryWithPhonetic}{下调}{xia4tiao2}{3,10}{⼀,⾔}[HSK 7-9]
  \definition{v.}{regular para baixo; ajustar para baixo}
  \seeref{xia4diao4}
\end{EntryWithPhonetic}

\begin{EntryWithPhonetic}{下午}{xia4wu3}{3,4}{⼀,⼗}[HSK 1]
  \definition[个]{s.}{tarde; \emph{post meridiem} (p.m.); refere"-se ao período entre o meio"-dia e o pôr do sol}
  \synonymref{午后}{wu3hou4}
  \antonymref{上午}{shang4wu3}
\end{EntryWithPhonetic}

\begin{EntryWithPhonetic}{下午茶}{xia4wu3cha2}{3,4,9}{⼀,⼗,⾋}
  \definition[杯]{s.}{chá da tarde (normalmente chás com doces)}
\end{EntryWithPhonetic}

\begin{EntryWithPhonetic}{下弦}{xia4xian2}{3,8}{⼀,⼸}
  \definition{s.}{lua minguante ou em quarto minguante; último (ou terceiro) quarto (da lua)}
  \antonymref{上弦}{shang4xian2}
\end{EntryWithPhonetic}

\begin{EntryWithPhonetic}{下线}{xia4xian4}{3,8}{⼀,⽷}
  \definition{v.}{sair da sessão; desconectar"-se da \emph{Internet}; refere"-se à suspensão temporária das atividades de comunicação \emph{on"-line}, geralmente significando interromper temporariamente o bate"-papo \emph{on"-line} ou os jogos \emph{on"-line} | sair da linha de produção; isso se refere a automóveis, eletrodomésticos, etc., que foram montados na linha de produção e estão prontos para sair da fábrica}
  \synonymref{离开}{li2/kai1}
  \antonymref{登陆}{deng1/lu4}
  \antonymref{问世}{wen4shi4}
  \antonymref{在线}{zai4xian4}
\end{EntryWithPhonetic}

\begin{EntryWithPhonetic}{下限}{xia4xian4}{3,8}{⼀,⾩}
  \definition{s.}{limite mínimo ou mais recente permitido; limite inferior; limiar; mínimo prescrito; piso (nível)}
  \antonymref{上限}{shang4xian4}
\end{EntryWithPhonetic}

\begin{EntryWithPhonetic}{下乡}{xia4/xiang1}{3,3}{⼀,⼄}[HSK 7-9]
  \definition{v.+compl.}{ir para o campo, zona rural (especialmente para intelectuais)}
\end{EntryWithPhonetic}

\begin{EntryWithPhonetic}{下雪}{xia4/xue3}{3,11}{⼀,⾬}[HSK 2]
  \definition{v.+compl.}{nevar}
\end{EntryWithPhonetic}

\begin{EntryWithPhonetic}{下旬}{xia4xun2}{3,6}{⼀,⽇}[HSK 7-9]
  \definition[月]{s.}{última dezena do mês; último período de dez dias de um mês; do dia 21 até o final de cada mês}
  \seealsoref{中旬}{zhong1xun2}
  \antonymref{上旬}{shang4xun2}
\end{EntryWithPhonetic}

\begin{EntryWithPhonetic}{下一代}{xia4yi2dai4}{3,1,5}{⼀,⼀,⼈}[HSK 7-9]
  \definition{s.}{próxima geração | geração mais jovem}[《星际迷航: 下一代》===Star Trek: A Nova Geração | 关心下一代委员会。===Comitê para o Cuidado da Próxima Geração.]
\end{EntryWithPhonetic}

\begin{EntryWithPhonetic}{下意识}{xia4yi4shi2}{3,13,7}{⼀,⼼,⾔}[HSK 7-9]
  \definition{adv.}{instintivamente; subconscientemente}
  \definition{s.}{instinto; subconsciente; a psicologia refere"-se às reações psicológicas inconscientes causadas por determinadas condições}
\end{EntryWithPhonetic}

\begin{EntryWithPhonetic}{下游}{xia4you2}{3,12}{⼀,⽔}[HSK 7-9]
  \definition{s.}{a jusante; rio abaixo; trechos inferiores; o trecho do rio próximo à sua foz | para trás; a posição inferior; referindo"-se metaforicamente a uma posição invertida}
  \antonymref{上游}{shang4you2}
\end{EntryWithPhonetic}

\begin{EntryWithPhonetic}{下雨}{xia4/yu3}{3,8}{⼀,⾬}[HSK 1]
  \definition{v.+compl.}{chover}
  \synonymref{暴雨}{bao4yu3}
  \synonymref{大雨}{da4yu3}
\end{EntryWithPhonetic}

\begin{EntryWithPhonetic}{下载}{xia4zai3}{3,10}{⼀,⾞}[HSK 4]
  \definition{v.}{\emph{download}; baixar; salvar informações da \emph{Web} em um dispositivo, como um computador}
  \antonymref{传输}{chuan2shu1}
\end{EntryWithPhonetic}

\begin{EntryWithPhonetic}{下崽}{xia4 zai3}{3,12}{⼀,⼭}
  \definition{v.}{dar à luz à filhotes; parir}
\end{EntryWithPhonetic}

\begin{EntryWithPhonetic}{下周}{xia4zhou1}{3,8}{⼀,⼝}[HSK 2]
  \definition{s.}{próxima semana}
  \antonymref{上周}{shang4zhou1}
\end{EntryWithPhonetic}

\begin{EntryWithPhonetic}{下坠}{xia4zhui4}{3,7}{⼀,⼟}[HSK 7-9]
  \definition{v.}{cair; deixar cair | Medicina: fazer força (ao defecar); ter urgência (na gravidez); ter tenesmo}
  \synonymref{笨重}{ben4zhong4}
  \synonymref{下降}{xia4jiang4}
  \synonymref{坠落}{zhui4luo4}
  \antonymref{飞跃}{fei1yue4}
\end{EntryWithPhonetic}

%%%%%%%%%% 吓 %%%%%%%%%%
\subsection*{吓}\addcontentsline{loh}{figure}{吓 \dpy{xia4}}

\begin{EntryWithPhonetic}{吓}{xia4}{6}{⼝}[HSK 5]
  \definition{interj.}{demonstrar espanto; expressar insatisfação}
  \definition{v.}{ameaçar; intimidar; usar ameaças ou meios coercitivos para intimidar ou assustar}
\end{EntryWithPhonetic}

\begin{EntryWithPhonetic}{吓唬}{xia4hu5}{6,11}{⼝,⼝}[HSK 7-9]
  \definition{v.}{assustar; amedrontar; punir; intimidar; ameaçar}
  \synonymref{恐吓}{kong3he4}
  \synonymref{威胁}{wei1xie2}
  \synonymref{吓人}{xia4/ren2}
  \antonymref{安慰}{an1wei4}
\end{EntryWithPhonetic}

\begin{EntryWithPhonetic}{吓人}{xia4/ren2}{6,2}{⼝,⼈}[HSK 7-9]
  \definition{adj.}{assustador; aterrorizante; apavorante; amedrontador; horrível}
  \definition{v.+compl.}{assustar}
  \synonymref{吓唬}{xia4hu5}
\end{EntryWithPhonetic}

%%%%%%%%%% 夏 %%%%%%%%%%
\subsection*{夏}\addcontentsline{loh}{figure}{夏 \dpy{xia4}}

\begin{EntryWithPhonetic}{夏}{xia4}{10}{⼢}
  \definition*{s.}{Dinastia Xia (2070--1600 a.C.) | China; refere"-se à China | Sobrenome: Xia}
  \definition{s.}{verão}
\end{EntryWithPhonetic}

\begin{EntryWithPhonetic}{夏季}{xia4ji4}{10,8}{⼢,⼦}[HSK 4]
  \definition[个]{s.}{verão; segundo trimestre do ano, habitualmente chamado na China de período de três meses, do início do verão ao início do outono, também chamado de ``quarto, quinto e sexto'' meses do calendário lunar}
\end{EntryWithPhonetic}

\begin{EntryWithPhonetic}{夏令营}{xia4ling4ying2}{10,5,11}{⼢,⼈,⾋}[HSK 7-9]
  \definition[个]{s.}{acampamento de verão; os acampamentos de verão, que oferecem descanso e recreação de curta duração para adolescentes ou membros de grupos, estão geralmente localizados em florestas ou à beira"-mar}
\end{EntryWithPhonetic}

\begin{EntryWithPhonetic}{夏日}{xia4ri4}{10,4}{⼢,⽇}
  \definition{s.}{dia de verão | sol de verão}
  \synonymref{夏季}{xia4ji4}
  \synonymref{夏天}{xia4tian1}
\end{EntryWithPhonetic}

\begin{EntryWithPhonetic}{夏天}{xia4tian1}{10,4}{⼢,⼤}[HSK 2]
  \definition[个]{s.}{verão}
  \seealsoref{夏}{xia4}
  \synonymref{夏}{xia4}
  \synonymref{夏季}{xia4ji4}
  \synonymref{夏日}{xia4ri4}
\end{EntryWithPhonetic}

\begin{EntryWithPhonetic}{夏禹}{xia4yu3}{10,9}{⼢,⽱}
  \seealsoref{大禹}{da4yu3}
\end{EntryWithPhonetic}

%%%%%%%%%% 仙 %%%%%%%%%%
\subsection*{仙}\addcontentsline{loh}{figure}{仙 \dpy{xian1}}

\begin{EntryWithPhonetic}{仙}{xian1}{5}{⼈}
  \definition*{s.}{Sobrenome: Xian}
  \definition[个,位,名,些]{s.}{ser celestial; imortal; divindade}
  \synonymref{神}{shen2}
\end{EntryWithPhonetic}

\begin{EntryWithPhonetic}{仙鹤}{xian1he4}{5,15}{⼈,⿃}[HSK 7-9]
  \definition[只]{s.}{grou branco (mantido pelos imortais) | grou"-da"-manchúria (\emph{Grus japonensis})}
\end{EntryWithPhonetic}

\begin{EntryWithPhonetic}{仙女}{xian1nv3}{5,3}{⼈,⼥}[HSK 7-9]
  \definition[个,位,名,些]{s.}{fada; celestial feminina; na mitologia, existem jovens mulheres que são imortais e possuem poderes mágicos}
  \synonymref{少女}{shao4nv3}
\end{EntryWithPhonetic}

%%%%%%%%%% 先 %%%%%%%%%%
\subsection*{先}\addcontentsline{loh}{figure}{先 \dpy{xian1}}

\begin{EntryWithPhonetic}{先}{xian1}{6}{⼉}[HSK 1]
  \definition*{s.}{Sobrenome: Xian}
  \definition{adv.}{primeiro; antes; mais cedo; com antecedência | no momento; por enquanto; em um curto espaço de tempo; temporariamente}
  \definition{s.}{início; começo; em ordem cronológica ou de precedência | ancestral; geração mais velha; antepassado | tardio; falecido; morto (honrar os mortos)}
  \seealsoref{首先}{shou3xian1}
  \synonymref{前}{qian2}
  \synonymref{首先}{shou3xian1}
  \synonymref{先前}{xian1qian2}
  \synonymref{以前}{yi3qian2}
  \synonymref{早}{zao3}
  \antonymref{后}{hou4}
\end{EntryWithPhonetic}

\begin{EntryWithPhonetic}{先不先}{xian1bu4xian1}{6,4,6}{⼉,⼀,⼉}
  \definition{adv.}{Dialeto: antes de tudo | em primeiro lugar}
\end{EntryWithPhonetic}

\begin{EntryWithPhonetic}{先到先得}{xian1dao4 xian1de2}{6,8,6,11}{⼉,⼑,⼉,⼻}
  \definition{adv.}{por ordem de chegada; o primeiro a chegar é o primeiro e ser servido}
\end{EntryWithPhonetic}

\begin{EntryWithPhonetic}{先锋}{xian1feng1}{6,12}{⼉,⾦}[HSK 6]
  \definition{s.}{pioneiro; vanguarda; a vanguarda de uma batalha ou marcha; geralmente se refere a uma pessoa ou grupo que desempenha um papel de vanguarda}
\end{EntryWithPhonetic}

\begin{EntryWithPhonetic}{先后}{xian1hou4}{6,6}{⼉,⼝}[HSK 5]
  \definition{adv.}{sucessivamente; um após o outro}
  \definition{s.}{prioridade; ordem; cedo ou tarde; primeiro e último}
\end{EntryWithPhonetic}

\begin{EntryWithPhonetic}{先进}{xian1jin4}{6,7}{⼉,⾡}[HSK 3]
  \definition{adj.}{avançado; progressos rápidos e nível elevado, podendo servir de exemplo a seguir}
  \definition{s.}{indivíduos ou grupos avançados}
\end{EntryWithPhonetic}

\begin{EntryWithPhonetic}{先例}{xian1li4}{6,8}{⼉,⼈}[HSK 7-9]
  \definition{s.}{precedente | caso anterior tomado como exemplo; exemplos existentes}
\end{EntryWithPhonetic}

\begin{EntryWithPhonetic}{先烈}{xian1lie4}{6,10}{⼉,⽕}
  \definition{s.}{mártir; título honorífico para mártires}
  \synonymref{烈士}{lie4shi4}
\end{EntryWithPhonetic}

\begin{EntryWithPhonetic}{先期}{xian1qi1}{6,12}{⼉,⽉}
  \definition{s.}{prematuro | \emph{front"-end}}
  \definition{s.}{antes da data agendada; mais cedo; com antecedência}
  \synonymref{事先}{shi4xian1}
  \synonymref{预先}{yu4xian1}
\end{EntryWithPhonetic}

\begin{EntryWithPhonetic}{先前}{xian1qian2}{6,9}{⼉,⼑}[HSK 5]
  \definition[出]{s.}{antes; anteriormente; refere"-se ao passado ou a um certo tempo anterior}
\end{EntryWithPhonetic}

\begin{EntryWithPhonetic}{先生}{xian1sheng5}{6,5}{⼉,⽣}[HSK 1]
  \definition[个,位]{s.}{professor; títulos honoríficos para professores, médicos, etc. | marido; antigamente, referia"-se ao marido de outra pessoa ou ao próprio marido (ambos com pronomes pessoais como determinantes) | médico; títulos usados para se referir aos médicos no passado | refere"-se a pessoas cuja profissão envolve contar histórias, adivinhação, etc.; antigamente, era chamado de contador | senhor; \emph{sir}; título dado aos intelectuais}
  \seealsoref{老师}{lao3shi1}
  \synonymref{教师}{jiao4shi1}
  \synonymref{老师}{lao3shi1}
  \synonymref{师长}{shi1zhang3}
  \synonymref{医生}{yi1sheng1}
  \antonymref{女士}{nv3shi4}
  \antonymref{学生}{xue2sheng5}
\end{EntryWithPhonetic}

\begin{EntryWithPhonetic}{先天}{xian1tian1}{6,4}{⼉,⼤}[HSK 7-9]
  \definition{adj.}{congênito; inato; o estágio embrionário de um ser humano ou animal | a priori; inato; em filosofia, refere"-se ao a priori}
  \definition{s.}{período embrionário}
  \antonymref{后天}{hou4tian1}
\end{EntryWithPhonetic}

\begin{EntryWithPhonetic}{先验}{xian1yan4}{6,10}{⼉,⾺}
  \definition{adj.}{Filosofia: \emph{a priori}}
\end{EntryWithPhonetic}

\begin{EntryWithPhonetic}{先有}{xian1you3}{6,6}{⼉,⽉}
  \definition{adj.}{preexistente | anterior}
\end{EntryWithPhonetic}

%%%%%%%%%% 纤 %%%%%%%%%%
\subsection*{纤}\addcontentsline{loh}{figure}{纤 \dpy{xian1}}

\begin{EntryWithPhonetic}{纤}{xian1}{6}{⽷}
  \definition{adj.}{fino; minúsculo}
  \seeref{qian4}
  \synonymref{小}{xiao3}
\end{EntryWithPhonetic}

\begin{EntryWithPhonetic}{纤维}{xian1wei2}{6,11}{⽷,⽷}[HSK 7-9]
  \definition[种,根,缕,束]{s.}{fibra; grampo; substâncias finas, longas e filamentosas encontradas nos corpos de plantas e animais; filamentos sintetizados artificialmente}
  \synonymref{化纤}{hua4xian1}
\end{EntryWithPhonetic}

%%%%%%%%%% 掀 %%%%%%%%%%
\subsection*{掀}\addcontentsline{loh}{figure}{掀 \dpy{xian1}}

\begin{EntryWithPhonetic}{掀}{xian1}{11}{⼿}[HSK 7-9]
  \definition{v.}{levantar (uma cobertura, etc.); descobrir | virar; levantar; tombar}
  \synonymref{揭}{jie1}
  \antonymref{盖}{gai4}
\end{EntryWithPhonetic}

\begin{EntryWithPhonetic}{掀起}{xian1qi3}{11,10}{⼿,⾛}[HSK 7-9]
  \definition{v.}{zarpar | levantar; erguer; descobrir | surgir; provocar uma onda | iniciar; dar início (a um movimento, etc.)}
  \antonymref{放下}{fang4xia4}
\end{EntryWithPhonetic}

%%%%%%%%%% 鲜 %%%%%%%%%%
\subsection*{鲜}\addcontentsline{loh}{figure}{鲜 \dpy{xian1}}

\begin{EntryWithPhonetic}{鲜}{xian1}{14}{⿂}[HSK 4]
  \definition*{s.}{Sobrenome: Xian}
  \definition{adj.}{fresco; novo; fresco (experiência, comida etc.) |brilhante; de cores vivas | saboroso; delicioso | exuberante; luxuriante}
  \definition{s.}{aves e animais recém"-abatidos; vegetais recém"-colhidos; frutas, etc. | alimentos aquáticos; geralmente, peixes vivos, camarões, etc., para alimentação}
  \seeref{xian3}
\end{EntryWithPhonetic}

\begin{EntryWithPhonetic}{鲜花}{xian1hua1}{14,7}{⿂,⾋}[HSK 4]
  \definition[朵,束,支]{s.}{flor; flores frescas; flores bonitas e frescas}
\end{EntryWithPhonetic}

\begin{EntryWithPhonetic}{鲜活}{xian1huo2}{14,9}{⿂,⽔}[HSK 7-9]
  \definition{adj.}{(produtos aquáticos, flores, etc.) fresco e vivo | vivaz; animado | (ingredientes alimentares) vivo ou fresco; carnes, ovos e laticínios frescos ou processados recentemente e que não estejam estragados | vívido; fresco e vibrante}
  \synonymref{生动}{sheng1dong4}
  \synonymref{水灵}{shui3ling2}
  \synonymref{新鲜}{xin1xian1}
  \antonymref{腐烂}{fu3lan4}
\end{EntryWithPhonetic}

\begin{EntryWithPhonetic}{鲜美}{xian1mei3}{14,9}{⿂,⽺}[HSK 7-9]
  \definition{adj.}{saboroso; delicioso | fresco e bonito}
  \synonymref{美味}{mei3wei4}
  \synonymref{新鲜}{xin1xian1}
  \antonymref{腐烂}{fu3lan4}
\end{EntryWithPhonetic}

\begin{EntryWithPhonetic}{鲜明}{xian1ming2}{14,8}{⿂,⽇}[HSK 4]
  \definition{adj.}{brilhante (cor) | distinto; bem definido; nítido; claro; característico}
\end{EntryWithPhonetic}

\begin{EntryWithPhonetic}{鲜血}{xian1xue4}{14,6}{⿂,⾎}[HSK 7-9]
  \definition[滴,股,滩]{s.}{sangue fresco; sangue vermelho vivo}
\end{EntryWithPhonetic}

\begin{EntryWithPhonetic}{鲜艳}{xian1yan4}{14,10}{⿂,⾊}[HSK 5]
  \definition{adj.}{de cores alegres; de cores brilhantes}
\end{EntryWithPhonetic}

%%%%%%%%%% 闲 %%%%%%%%%%
\subsection*{闲}\addcontentsline{loh}{figure}{闲 \dpy{xian2}}

\begin{EntryWithPhonetic}{闲}{xian2}{7}{⾨}[HSK 5]
  \definition{adj.}{ocioso; não ocupado; desocupado; sem coisas para fazer; sem atividades; tempo livre | desocupado; (casa, objeto, etc.) não em uso; ocioso | não oficial; não sério; não relacionado ao negócio}
  \definition{s.}{lazer; tempo livre}
\end{EntryWithPhonetic}

%%%%%%%%%% 弦 %%%%%%%%%%
\subsection*{弦}\addcontentsline{loh}{figure}{弦 \dpy{xian2}}

\begin{EntryWithPhonetic}{弦}{xian2}{8}{⼸}[HSK 7-9]
  \definition*{s.}{Sobrenome: Xian}
  \definition[根]{s.}{corda do arco; corda | Dialeto: mola (de um relógio, etc.) | corda de um instrumento musical | acorde | Matemática: hipotenusa}
\end{EntryWithPhonetic}

%%%%%%%%%% 咸 %%%%%%%%%%
\subsection*{咸}\addcontentsline{loh}{figure}{咸 \dpy{xian2}}

\begin{EntryWithPhonetic}{咸}{xian2}{9}{⼝}[HSK 4]
  \definition*{s.}{Sobrenome: Xian}
  \definition{adj.}{salgado; em conserva; sabor salgado}
  \definition{adv.}{todos; indica a totalidade de um intervalo, equivalente a 全 e 都}
  \seealsoref{都}{dou1}
  \seealsoref{全}{quan2}
\end{EntryWithPhonetic}

\begin{EntryWithPhonetic}{咸菜}{xian2cai4}{9,11}{⼝,⾋}
  \definition{s.}{legumes salgados; \emph{pickles}}
\end{EntryWithPhonetic}

\begin{EntryWithPhonetic}{咸淡}{xian2dan4}{9,11}{⼝,⽔}
  \definition{s.}{água salobra | grau de salinidade | salgado e sem sal (sabores)}
\end{EntryWithPhonetic}

\begin{EntryWithPhonetic}{咸肉}{xian2rou4}{9,6}{⼝,⾁}
  \definition{s.}{carne salgada; \emph{bacon} | carne curada com sal}
\end{EntryWithPhonetic}

\begin{EntryWithPhonetic}{咸涩}{xian2se4}{9,10}{⼝,⽔}
  \definition{adj.}{ácido | salgado e amargo}
\end{EntryWithPhonetic}

\begin{EntryWithPhonetic}{咸水}{xian2shui3}{9,4}{⼝,⽔}
  \definition{s.}{água salgada | salmoura}
\end{EntryWithPhonetic}

\begin{EntryWithPhonetic}{咸盐}{xian2yan2}{9,10}{⼝,⽫}
  \definition{s.}{Dialeto: sal de mesa; sal | Coloquial: sal}
\end{EntryWithPhonetic}

\begin{EntryWithPhonetic}{咸鱼}{xian2yu2}{9,8}{⼝,⿂}
  \definition{adj.}{preguiçoso; folgado; refere"-se à perda de esperança na vida e à falta de sentido que as pessoas encontram nos tempos modernos}
  \definition{s.}{peixe salgado; peixe curado com sal}
\end{EntryWithPhonetic}

%%%%%%%%%% 衔 %%%%%%%%%%
\subsection*{衔}\addcontentsline{loh}{figure}{衔 \dpy{xian2}}

\begin{EntryWithPhonetic}{衔}{xian2}{11}{⾏}
  \definition{s.}{patente; título}
  \definition{v.}{manter na boca | abrigar; suportar | receber ou executar (ordens, etc.) | conectar"-se; participar | receber; encomendar | aceitar; assumir}
\end{EntryWithPhonetic}

\begin{EntryWithPhonetic}{衔接}{xian2jie1}{11,11}{⾏,⼿}[HSK 7-9]
  \definition{v.}{juntar"-se; conectar"-se; estabelecer vínculo}
  \seealsoref{对接}{dui4jie1}
  \synonymref{接连}{jie1lian2}
  \synonymref{连接}{lian2jie1}
  \synonymref{相连}{xiang1lian2}
  \antonymref{脱节}{tuo1/jie2}
\end{EntryWithPhonetic}

%%%%%%%%%% 嫌 %%%%%%%%%%
\subsection*{嫌}\addcontentsline{loh}{figure}{嫌 \dpy{xian2}}

\begin{EntryWithPhonetic}{嫌}{xian2}{13}{⼥}[HSK 6]
  \definition{s.}{suspeita; suspeição; cisma | inimizade; rancor; má vontade; ressentimento}
  \definition{v.}{importar"-se com; não gostar e evitar; reclamar de}
\end{EntryWithPhonetic}

\begin{EntryWithPhonetic}{嫌弃}{xian2qi4}{13,7}{⼥,⼶}[HSK 7-9]
  \definition{v.}{ignorar; não gostar e evitar; antipatia e indisposição para se aproximar}
  \seealsoref{嫌}{xian2}
  \antonymref{宠爱}{chong3'ai4}
  \antonymref{亲近}{qin1jin4}
  \antonymref{羡慕}{xian4mu4}
  \antonymref{拥护}{yong1hu4}
\end{EntryWithPhonetic}

\begin{EntryWithPhonetic}{嫌疑}{xian2yi2}{13,14}{⼥,⽦}[HSK 7-9]
  \definition{s.}{suspeita; a possibilidade de ser suspeito de certos comportamentos}
  \synonymref{猜疑}{cai1yi2}
  \synonymref{怀疑}{huai2yi2}
  \synonymref{可疑}{ke3yi2}
  \synonymref{疑惑}{yi2huo4}
  \antonymref{信任}{xin4ren4}
\end{EntryWithPhonetic}

%%%%%%%%%% 显 %%%%%%%%%%
\subsection*{显}\addcontentsline{loh}{figure}{显 \dpy{xian3}}

\begin{EntryWithPhonetic}{显}{xian3}{9}{⽇}[HSK 5]
  \definition*{s.}{Sobrenome: Xian}
  \definition{adj.}{aparente; óbvio; perceptível | ilustre e influente | evidente; óbvio}
  \definition{v.}{mostrar; exibir; manifestar | aparecer; mostrar; revelar}
\end{EntryWithPhonetic}

\begin{EntryWithPhonetic}{显出}{xian3chu1}{9,5}{⽇,⼐}[HSK 6]
  \definition{v.}{mostrar; revelar | dar provas; expressar; exibir}
\end{EntryWithPhonetic}

\begin{EntryWithPhonetic}{显得}{xian3de5}{9,11}{⽇,⼻}[HSK 3]
  \definition{v.}{parecer; aparecer; manifestar (alguma situação)}
\end{EntryWithPhonetic}

\begin{EntryWithPhonetic}{显而易见}{xian3'er2yi4jian4}{9,6,8,4}{⽇,⽽,⽇,⾒}[HSK 7-9]
  \definition{expr.}{óbvio; aparente; tão claro quanto o nariz no rosto; a questão e o raciocínio são muito óbvios e fáceis de entender}
\end{EntryWithPhonetic}

\begin{EntryWithPhonetic}{显赫}{xian3he4}{9,14}{⽇,⾚}[HSK 7-9]
  \definition{adj.}{ilustre; célebre; proeminente; excepcional, notável, respeitado e famoso}
  \synonymref{显著}{xian3zhu4}
  \synonymref{卓越}{zhuo2yue4}
\end{EntryWithPhonetic}

\begin{EntryWithPhonetic}{显然}{xian3ran2}{9,12}{⽇,⽕}[HSK 3]
  \definition{adj.}{claro; evidente; óbvio; fatos, verdades e outras coisas que são fáceis de descobrir, perceber ou sentir claramente}
\end{EntryWithPhonetic}

\begin{EntryWithPhonetic}{显示}{xian3shi4}{9,5}{⽇,⽰}[HSK 3]
  \definition{v.}{mostrar; manifestar"-se claramente | exibir; ostentar}
\end{EntryWithPhonetic}

\begin{EntryWithPhonetic}{显示器}{xian3shi4qi4}{9,5,16}{⽇,⽰,⼝}[HSK 7-9]
  \definition[台,个]{s.}{tela; monitor; \emph{display}}
\end{EntryWithPhonetic}

\begin{EntryWithPhonetic}{显现}{xian3xian4}{9,8}{⽇,⾒}[HSK 7-9]
  \definition{v.}{manifestar"-se (ou revelar"-se); aparecer; mostrar"-se | manifestar"-se; revelar"-se; visualizar}
  \synonymref{暴露}{bao4lu4}
  \synonymref{呈现}{cheng2xian4}
  \synonymref{出现}{chu1xian4}
  \synonymref{流露}{liu2lu4}
  \synonymref{清楚}{qing1chu5}
  \synonymref{透露}{tou4lu4}
  \synonymref{显示}{xian3shi4}
  \synonymref{展示}{zhan3shi4}
  \synonymref{展现}{zhan3xian4}
  \antonymref{笼罩}{long3zhao4}
  \antonymref{隐蔽}{yin3bi4}
  \antonymref{隐藏}{yin3cang2}
\end{EntryWithPhonetic}

\begin{EntryWithPhonetic}{显眼}{xian3yan3}{9,11}{⽇,⽬}[HSK 7-9]
  \definition{adj.}{vistoso; chamativo; excepcional; atraente; conspícuo; facilmente visível; que atrai a atenção}
  \synonymref{显得}{xian3de5}
  \synonymref{醒目}{xing3mu4}
  \antonymref{隐蔽}{yin3bi4}
\end{EntryWithPhonetic}

\begin{EntryWithPhonetic}{显著}{xian3zhu4}{9,11}{⽇,⽬}[HSK 4]
  \definition{adj.}{notável; significativo; notável; extraordinário; muito óbvio; muito claramente demonstrado; muito facilmente visto ou sentido}
\end{EntryWithPhonetic}

%%%%%%%%%% 险 %%%%%%%%%%
\subsection*{险}\addcontentsline{loh}{figure}{险 \dpy{xian3}}

\begin{EntryWithPhonetic}{险}{xian3}{9}{⾩}[HSK 6]
  \definition{adj.}{perigoso; arriscado | sinistro; cruel; venenoso}
  \definition{adv.}{por um fio de cabelo; por centímetros; quase}
  \definition{s.}{lugar de difícil acesso; lugar perigoso e difícil de atravessar; passagem estreita; desfiladeiro | abreviação de seguro, 保险 | perigo; risco}
  \seealsoref{保险}{bao3xian3}
\end{EntryWithPhonetic}

%%%%%%%%%% 猃 %%%%%%%%%%
\subsection*{猃}\addcontentsline{loh}{figure}{猃 \dpy{xian3}}

\begin{EntryWithPhonetic}{猃}{xian3}{10}{⽝}
  \definition{s.}{(arcaico) um tipo de cão com focinho longo}
\end{EntryWithPhonetic}

\begin{EntryWithPhonetic}{猃狁}{xian3 yun3}{10,7}{⽝,⽝}
  \definition*{s.}{Termo da dinastia Zhou para uma tribo nômade do norte mais tarde chamou o Xiongnu (匈奴) nas dinastias Qin e Han}
  \definition{s.}{um povo nômade que outrora habitava a região norte da China}
  \seealsoref{匈奴}{xiong1nu2}
  \seealsoref{狁}{yun3}
\end{EntryWithPhonetic}

%%%%%%%%%% 鲜 %%%%%%%%%%
\subsection*{鲜}\addcontentsline{loh}{figure}{鲜 \dpy{xian3}}

\begin{EntryWithPhonetic}{鲜}{xian3}{14}{⿂}
  \definition{adj.}{raro; pouco; pequeno}
  \definition{adv.}{raramente}
  \seeref{xian1}
\end{EntryWithPhonetic}

%%%%%%%%%% 见 %%%%%%%%%%
\subsection*{见}\addcontentsline{loh}{figure}{见 \dpy{xian4}}

\begin{EntryWithPhonetic}{见}{xian4}{4}{⾒}[Kangxi 147]
  \definition{v.}{aparecer; também escrito como 现}
  \seeref{jian4}
  \seealsoref{现}{xian4}
\end{EntryWithPhonetic}

%%%%%%%%%% 县 %%%%%%%%%%
\subsection*{县}\addcontentsline{loh}{figure}{县 \dpy{xian4}}

\begin{EntryWithPhonetic}{县}{xian4}{7}{⼛}[HSK 4]
  \definition[个]{s.}{condado; unidade de divisão administrativa}
\end{EntryWithPhonetic}

%%%%%%%%%% 现 %%%%%%%%%%
\subsection*{现}\addcontentsline{loh}{figure}{现 \dpy{xian4}}

\begin{EntryWithPhonetic}{现}{xian4}{8}{⾒}
  \definition{adj.}{(dinheiro, etc.) em mãos}
  \definition{adv.}{recente; de improviso; naquela época; temporariamente}
  \definition{s.}{presente; atual; existente | dinheiro; dinheiro de pronto}
  \definition{v.}{mostrar; revelar; aparecer; tornar"-se visível}
  \seealsoref{见}{xian4}
\end{EntryWithPhonetic}

\begin{EntryWithPhonetic}{现场}{xian4chang3}{8,6}{⾒,⼟}[HSK 3]
  \definition[个,处]{s.}{local onde ocorreu o acidente, incidente ou desastre | local; ponto; local onde se realizam diretamente atividades como produção, apresentações e competições}
\end{EntryWithPhonetic}

\begin{EntryWithPhonetic}{现成}{xian4cheng2}{8,6}{⾒,⼽}[HSK 7-9]
  \definition{adj.}{pronto para uso; facilmente disponível; já existia; não é necessário nenhum preparo adicional}
  \synonymref{现有}{xian4you3}
  \synonymref{原有}{yuan2you3}
  \antonymref{发明}{fa1ming2}
\end{EntryWithPhonetic}

\begin{EntryWithPhonetic}{现出}{xian4/chu1}{8,5}{⾒,⼐}
  \definition{v.+compl.}{revelar; exibir}
  \seealsoref{出现}{chu1xian4}
  \synonymref{出现}{chu1xian4}
\end{EntryWithPhonetic}

\begin{EntryWithPhonetic}{现代}{xian4dai4}{8,5}{⾒,⼈}[HSK 3]
  \definition*{s.}{Hyundai, empresa sul"-coreana}
  \definition{adj.}{moderno; contemporâneo; com características, estilo e conceitos modernos, refletindo a vanguarda, a moda e a inovação da atualidade}
  \definition{s.}{tempos modernos; era contemporânea; atualmente, na divisão cronológica da história da China, refere"-se principalmente ao período desde o Movimento 4 de Maio até os dias atuais}
\end{EntryWithPhonetic}

\begin{EntryWithPhonetic}{现货}{xian4huo4}{8,8}{⾒,⾙}
  \definition{s.}{produtos em estoque; mercadorias que podem ser entregues no momento; item disponível}
\end{EntryWithPhonetic}

\begin{EntryWithPhonetic}{现货的}{xian4huo4 de5}{8,8,8}{⾒,⾙,⽩}
  \definition{s.}{produtos à vista; produtos em estoque}
\end{EntryWithPhonetic}

\begin{EntryWithPhonetic}{现金}{xian4jin1}{8,8}{⾒,⾦}[HSK 3]
  \definition[笔]{s.}{dinheiro; dinheiro vivo; moeda que pode ser usada diretamente | reserva de dinheiro em um banco; o dinheiro guardado no cofre do banco}
\end{EntryWithPhonetic}

\begin{EntryWithPhonetic}{现任}{xian4ren4}{8,6}{⾒,⼈}[HSK 7-9]
  \definition{adj.}{atualmente em exercício; titular do cargo; atualmente no escritório | atualmente ocupa o cargo de; atualmente em exercício como}
  \definition{v.}{ocupar um cargo; estar na posição de}
  \synonymref{此刻}{ci3ke4}
  \synonymref{当今}{dang1jin1}
  \synonymref{如今}{ru2jin1}
  \synonymref{现在}{xian4zai4}
  \antonymref{前任}{qian2ren4}
\end{EntryWithPhonetic}

\begin{EntryWithPhonetic}{现实}{xian4shi2}{8,8}{⾒,⼧}[HSK 3]
  \definition{adj.}{real; efetivo; verdadeiro; de acordo com circunstâncias objetivas}
  \definition[个]{s.}{realidade; factualidade; coisas que existem objetivamente}
\end{EntryWithPhonetic}

\begin{EntryWithPhonetic}{现象}{xian4xiang4}{8,11}{⾒,⾗}[HSK 3]
  \definition[个,种]{s.}{aparência (das coisas); fenômeno; a forma externa e as relações manifestadas pelas coisas em seu desenvolvimento e mudança}
\end{EntryWithPhonetic}

\begin{EntryWithPhonetic}{现行}{xian4xing2}{8,6}{⾒,⾏}[HSK 7-9]
  \definition{adj.}{atualmente em vigor; em operação; em funcionamento; atualmente válido | (um criminoso) ativo}
  \synonymref{流行}{liu2xing2}
  \synonymref{现在}{xian4zai4}
\end{EntryWithPhonetic}

\begin{EntryWithPhonetic}{现有}{xian4you3}{8,6}{⾒,⽉}[HSK 5]
  \definition{adj.}{agora disponível; existente}
  \definition{v.}{estar disponível agora; existir | Literário: ter em mãos; ter em posse}
\end{EntryWithPhonetic}

\begin{EntryWithPhonetic}{现在}{xian4zai4}{8,6}{⾒,⼟}[HSK 1]
  \definition{adv.}{agora; no momento; atualmente; neste momento, quando se fala, às vezes inclui um período de tempo mais ou menos longo antes ou depois da fala (diferente de 过去 ou 将来)}
  \seealsoref{过去}{guo4qu5}
  \seealsoref{将来}{jiang1lai2}
  \seealsoref{如今}{ru2jin1}
  \seealsoref{眼前}{yan3qian2}
  \synonymref{此刻}{ci3ke4}
  \synonymref{当今}{dang1jin1}
  \synonymref{当下}{dang1xia4}
  \synonymref{如今}{ru2jin1}
  \synonymref{正在}{zheng4zai4}
  \antonymref{从前}{cong2qian2}
  \antonymref{当初}{dang1chu1}
  \antonymref{过去}{guo4qu5}
  \antonymref{将来}{jiang1lai2}
  \antonymref{未来}{wei4lai2}
  \antonymref{先前}{xian1qian2}
  \antonymref{以前}{yi3qian2}
\end{EntryWithPhonetic}

\begin{EntryWithPhonetic}{现抓}{xian4zhua1}{8,7}{⾒,⼿}
  \definition{v.}{improvisar}
\end{EntryWithPhonetic}

\begin{EntryWithPhonetic}{现状}{xian4zhuang4}{8,7}{⾒,⽝}[HSK 5]
  \definition{s.}{situação atual}
\end{EntryWithPhonetic}

\begin{EntryWithPhonetic}{现做}{xian4zuo4}{8,11}{⾒,⼈}
  \definition{adj.}{preparar (comida) na hora; fazer sob encomenda}
  \definition{v.}{fazer (comida) no local}
\end{EntryWithPhonetic}

%%%%%%%%%% 线 %%%%%%%%%%
\subsection*{线}\addcontentsline{loh}{figure}{线 \dpy{xian4}}

\begin{EntryWithPhonetic}{线}{xian4}{8}{⽷}[HSK 3]
  \definition{clas.}{usado para coisas abstratas, o número é limitado a ``一'' (1)}
  \definition[根,个]{s.}{fio; corda; arame; objetos finos e longos feitos de seda, algodão, metal, etc. | linha; figura formada pelo movimento arbitrário de um ponto| feito de fio de algodão | algo em forma de linha, fio, etc. | rota de transporte; linha | linha de demarcação; limite; zona de fronteira; zona de transição | beira; borda | linha ideológica e política | pista; fio}
\end{EntryWithPhonetic}

\begin{EntryWithPhonetic}{线路}{xian4lu4}{8,13}{⽷,⾜}[HSK 6]
  \definition[条]{s.}{linha; rota; as rotas que os veículos de transporte percorrem, etc., que as pessoas podem usar para chegar aos seus destinos | Eletricidade: linha; circuito; a rota da corrente elétrica}
\end{EntryWithPhonetic}

\begin{EntryWithPhonetic}{线索}{xian4suo3}{8,10}{⽷,⽷}[HSK 5]
  \definition[条,个]{s.}{pista; fio; metáfora para o desenvolvimento das coisas ou a maneira de explorar um problema | fio; linha; refere"-se ao contexto de desenvolvimento do enredo em obras literárias}
\end{EntryWithPhonetic}

\begin{EntryWithPhonetic}{线条}{xian4tiao2}{8,7}{⽷,⽊}[HSK 7-9]
  \definition[根,道,笔]{s.}{linha (no desenho); linhas desenhadas com caneta ou outras ferramentas | as linhas ou contornos de um objeto tridimensional; a linha mais externa de uma pessoa ou objeto, geralmente referindo"-se à sua forma ou aparência}
  \synonymref{曲线}{qu1xian4}
\end{EntryWithPhonetic}

\begin{EntryWithPhonetic}{线香}{xian4xiang1}{8,9}{⽷,⾹}
  \definition{s.}{bastão fino de incenso}
\end{EntryWithPhonetic}

%%%%%%%%%% 限 %%%%%%%%%%
\subsection*{限}\addcontentsline{loh}{figure}{限 \dpy{xian4}}

\begin{EntryWithPhonetic}{限}{xian4}{8}{⾩}[HSK 7-9]
  \definition{s.}{limite; o intervalo especificado não deve exceder | intervalo especificado; limite}
  \definition{v.}{estabelecer um limite; limitar; restringir}
\end{EntryWithPhonetic}

\begin{EntryWithPhonetic}{限定}{xian4ding4}{8,8}{⾩,⼧}[HSK 7-9]
  \definition{v.}{limitar; restringir; prescrever (ou estabelecer) um limite para; o âmbito do tempo, espaço ou ação especificados}
  \synonymref{规定}{gui1ding4}
  \synonymref{截至}{jie2zhi4}
  \synonymref{界定}{jie4ding4}
  \synonymref{戒指}{jie4zhi5}
  \synonymref{局限}{ju2xian4}
  \synonymref{绝版}{jue2ban3}
  \synonymref{限度}{xian4du4}
  \synonymref{限制}{xian4zhi4}
  \antonymref{放松}{fang4song1}
\end{EntryWithPhonetic}

\begin{EntryWithPhonetic}{限度}{xian4du4}{8,9}{⾩,⼴}[HSK 7-9]
  \definition{s.}{faixa; limite; medida; limitação; a faixa ou grau especificado; a faixa ou grau mais alto ou mais baixo}
  \synonymref{范围}{fan4wei2}
  \synonymref{截至}{jie2zhi4}
  \synonymref{戒指}{jie4zhi5}
  \synonymref{局部}{ju2bu4}
  \synonymref{局限}{ju2xian4}
  \synonymref{限定}{xian4ding4}
  \synonymref{限制}{xian4zhi4}
\end{EntryWithPhonetic}

\begin{EntryWithPhonetic}{限于}{xian4yu2}{8,3}{⾩,⼆}[HSK 7-9]
  \definition{v.}{ser limitado a; ser sujeito a certas condições ou circunstâncias | ser limitado a; ser limitado a um determinado intervalo}
\end{EntryWithPhonetic}

\begin{EntryWithPhonetic}{限制}{xian4zhi4}{8,8}{⾩,⼑}[HSK 4]
  \definition[些]{s.}{limite; restrição; confinamento}
  \definition{v.}{limitar; adstringir; restringir; confinar; fechar em (sobre)}
\end{EntryWithPhonetic}

%%%%%%%%%% 宪 %%%%%%%%%%
\subsection*{宪}\addcontentsline{loh}{figure}{宪 \dpy{xian4}}

\begin{EntryWithPhonetic}{宪}{xian4}{9}{⼧}
  \definition*{s.}{Sobrenome: Xian}
  \definition{s.}{estatuto; decreto | constituição}
\end{EntryWithPhonetic}

\begin{EntryWithPhonetic}{宪法}{xian4fa3}{9,8}{⼧,⽔}[HSK 7-9]
  \definition[部]{s.}{constituição; a lei fundamental de um Estado possui a mais alta autoridade legal e serve como base para toda a legislação; ela geralmente estipula o sistema social, o sistema estatal, as instituições estatais e os direitos e obrigações básicos de um país}
\end{EntryWithPhonetic}

\begin{EntryWithPhonetic}{宪法法院}{xian4fa3fa3yuan4}{9,8,8,9}{⼧,⽔,⽔,⾩}
  \definition*{s.}{Tribunal Constitucional; Suprema Corte}
\end{EntryWithPhonetic}

\begin{EntryWithPhonetic}{宪政}{xian4zheng4}{9,9}{⼧,⽁}
  \definition{s.}{governo constitucional; constitucionalismo}
\end{EntryWithPhonetic}

\begin{EntryWithPhonetic}{宪制}{xian4zhi4}{9,8}{⼧,⼑}
  \definition*{s.}{Constituição}
  \definition{adj.}{constitucional}
\end{EntryWithPhonetic}

%%%%%%%%%% 陷 %%%%%%%%%%
\subsection*{陷}\addcontentsline{loh}{figure}{陷 \dpy{xian4}}

\begin{EntryWithPhonetic}{陷}{xian4}{10}{⾩}[HSK 7-9]
  \definition[个]{s.}{armadilha; cilada | defeito | deficiência; desvantagem}
  \definition{v.}{ficar preso (ou atolado); enredar | afundar; desabar | acusar falsamente; incriminar; armar | (uma cidade, etc.) ser capturado; cair | ser enquadrado; ser capturado}
\end{EntryWithPhonetic}

\begin{EntryWithPhonetic}{陷阱}{xian4jing3}{10,6}{⾩,⾩}[HSK 7-9]
  \definition[个,处]{s.}{fosso; armadilha; laço; uma cova cavada para capturar animais selvagens ou inimigos | armadilha; fosso; uma metáfora para uma armadilha ou plano para prejudicar alguém}
  \synonymref{机关}{ji1guan1}
  \synonymref{圈套}{quan1tao4}
  \synonymref{组织}{zu3zhi1}
  \antonymref{机遇}{ji1yu4}
\end{EntryWithPhonetic}

\begin{EntryWithPhonetic}{陷入}{xian4ru4}{10,2}{⾩,⼊}[HSK 6]
  \definition{v.}{afundar em; cair em; cair em uma situação desfavorável | estar perdido em; estar profundamente em; estar imerso em; metaforicamente, estar profundamente imerso em (uma situação ou pensamento) | estar atolado (lama fofa, areia, etc.)}
\end{EntryWithPhonetic}

%%%%%%%%%% 馅 %%%%%%%%%%
\subsection*{馅}\addcontentsline{loh}{figure}{馅 \dpy{xian4}}

\begin{EntryWithPhonetic}{馅}{xian4}{11}{⾷}
  \definition[种]{s.}{(comida) recheio | Coloquial: informação privilegiada; segredo | enchimento; massas, doces, massas com pasta de feijão vermelho ou vegetais, e outras sobremesas}
  \seealsoref{馅儿}{xian4r5}
\end{EntryWithPhonetic}

\begin{EntryWithPhonetic}{馅儿}{xian4r5}{11,2}{⾷,⼉}[HSK 7-9]
  \definition{s.}{recheio (de comida); carne, legumes, açúcar, etc., envoltos em massa ou salgadinhos}
\end{EntryWithPhonetic}

%%%%%%%%%% 羡 %%%%%%%%%%
\subsection*{羡}\addcontentsline{loh}{figure}{羡 \dpy{xian4}}

\begin{EntryWithPhonetic}{羡}{xian4}{12}{⽺}
  \definition{v.}{admirar; invejar}
\end{EntryWithPhonetic}

\begin{EntryWithPhonetic}{羡慕}{xian4mu4}{12,14}{⽺,⼼}[HSK 7-9]
  \definition{v.}{invejar; admirar; ver os outros terem certos pontos fortes ou vantagens e desejar tê"-los também}
  \seealsoref{忌妒}{ji4du4}
  \synonymref{爱戴}{ai4dai4}
  \synonymref{向往}{xiang4wang3}
  \antonymref{鄙视}{bi3shi4}
  \antonymref{妒忌}{du4ji4}
  \antonymref{嫉妒}{ji2du4}
  \antonymref{忌妒}{ji4du4}
  \antonymref{讨厌}{tao3/yan4}
  \antonymref{嫌弃}{xian2qi4}
  \antonymref{眼红}{yan3hong2}
\end{EntryWithPhonetic}

%%%%%%%%%% 献 %%%%%%%%%%
\subsection*{献}\addcontentsline{loh}{figure}{献 \dpy{xian4}}

\begin{EntryWithPhonetic}{献}{xian4}{13}{⽝}[HSK 5]
  \definition{v.}{oferecer; apresentar; dedicar; doar | mostrar; apresentar; exibir | exibir"-se; mostrar"-se para que os outros vejam}
\end{EntryWithPhonetic}

\begin{EntryWithPhonetic}{献血}{xian4xie3}{13,6}{⽝,⾎}[HSK 7-9]
  \definition{v.}{doar sangue; doar sangue para instituições médicas especializadas e ajude quem precisa}
\end{EntryWithPhonetic}

%%%%%%%%%% 腺 %%%%%%%%%%
\subsection*{腺}\addcontentsline{loh}{figure}{腺 \dpy{xian4}}

\begin{EntryWithPhonetic}{腺}{xian4}{13}{⾁}[HSK 7-9]
  \definition[个]{s.}{glândula; tecidos em um organismo que podem secretar certas substâncias químicas}
\end{EntryWithPhonetic}

%%%%%%%%%% 乡 %%%%%%%%%%
\subsection*{乡}\addcontentsline{loh}{figure}{乡 \dpy{xiang1}}

\begin{EntryWithPhonetic}{乡}{xiang1}{3}{⼄}[HSK 5]
  \definition[个,座,片]{s.}{país; campo; vilarejo; área rural | local de origem; vila ou cidade natal | município (uma unidade administrativa rural subordinada ao condado) | vila natal; cidade natal | terra ou local famoso por produzir algo}
  \antonymref{城}{cheng2}
\end{EntryWithPhonetic}

\begin{EntryWithPhonetic}{乡巴佬}{xiang1ba1lao3}{3,4,8}{⼄,⼰,⼈}
  \definition{s.}{Pejorativo: aldeão; caipira}
\end{EntryWithPhonetic}

\begin{EntryWithPhonetic}{乡村}{xiang1cun1}{3,7}{⼄,⽊}[HSK 5]
  \definition{adj.}{rural | rústico}
  \definition[个]{s.}{vila; campo; área rural; principalmente envolvido na agricultura; áreas com distribuição populacional mais dispersa em relação às cidades}
  \synonymref{农村}{nong2cun1}
  \antonymref{城市}{cheng2shi4}
\end{EntryWithPhonetic}

\begin{EntryWithPhonetic}{乡亲}{xiang1qin1}{3,9}{⼄,⼇}[HSK 7-9]
  \definition{s.}{pessoa da mesma aldeia ou cidade; conterrâneo ou morador da mesma cidade | população local; aldeões; gente; povo local | termo de tratamento direto para pessoas locais ou moradores de vilarejos}
  \synonymref{故乡}{gu4xiang1}
  \synonymref{家园}{jia1yuan2}
\end{EntryWithPhonetic}

\begin{EntryWithPhonetic}{乡下}{xiang1xia4}{3,3}{⼄,⼀}[HSK 7-9]
  \definition[个]{s.}{país; campo; interior; área rural}
  \synonymref{村庄}{cun1zhuang1}
  \synonymref{农村}{nong2cun1}
  \synonymref{乡村}{xiang1cun1}
  \antonymref{城市}{cheng2shi4}
  \antonymref{都市}{du1shi4}
\end{EntryWithPhonetic}

%%%%%%%%%% 相 %%%%%%%%%%
\subsection*{相}\addcontentsline{loh}{figure}{相 \dpy{xiang1}}

\begin{EntryWithPhonetic}{相}{xiang1}{9}{⽬}
  \definition*{s.}{Sobrenome: Xiang}
  \definition{adv.}{uns aos outros; mutuamente | (para uma ação realizada por uma pessoa em relação a outra) | indica a ação de uma parte em relação à outra parte}
  \definition{s.}{qualidade; substância}
  \definition{v.}{ver por si mesmo (se algo ou algo é do seu agrado)}
  \seeref{xiang4}
\end{EntryWithPhonetic}

\begin{EntryWithPhonetic}{相伴}{xiang1ban4}{9,7}{⽬,⼈}[HSK 7-9]
  \definition{s.}{estar juntos; acompanhar um ao outro; viver ou realizar outras atividades com alguém ou algo}
  \synonymref{陪伴}{pei2ban4}
\end{EntryWithPhonetic}

\begin{EntryWithPhonetic}{相比}{xiang1bi3}{9,4}{⽬,⽐}[HSK 3]
  \definition{v.}{combinar; comparar com | comparar mutuamente, usar uma coisa como padrão, perceber as características de outra coisa ou obter uma opinião}
\end{EntryWithPhonetic}

\begin{EntryWithPhonetic}{相比之下}{xiang1bi3 zhi1xia4}{9,4,3,3}{⽬,⽐,⼂,⼀}[HSK 7-9]
  \definition{expr.}{por comparação; em comparação; destacar características diferentes através da comparação com outra coisa}
\end{EntryWithPhonetic}

\begin{EntryWithPhonetic}{相差}{xiang1cha4}{9,9}{⽬,⼯}[HSK 7-9]
  \definition{adj.}{diferente}
  \synonymref{出入}{chu1ru4}
  \synonymref{进出}{jin4chu1}
  \synonymref{收支}{shou1zhi1}
\end{EntryWithPhonetic}

\begin{EntryWithPhonetic}{相处}{xiang1chu3}{9,5}{⽬,⼡}[HSK 4]
  \definition{v.}{dar"-se bem; viver juntos; dar"-se bem (uns com os outros); viver uns com os outros; entrar em contato uns com os outros, tratar uns aos outros}
\end{EntryWithPhonetic}

\begin{EntryWithPhonetic}{相传}{xiang1chuan2}{9,6}{⽬,⼈}[HSK 7-9]
  \definition{v.}{diz a lenda que\dots; segundo a lenda; a tradição conta que\dots; tem sido transmitido ou circulado por muito tempo (referindo"-se a uma afirmação sobre algo que não se baseia em evidências reais, mas é meramente um boato) | passar a mão; repassar; transmitir}
  \synonymref{传递}{chuan2di4}
  \antonymref{危险}{wei1xian3}
\end{EntryWithPhonetic}

\begin{EntryWithPhonetic}{相当}{xiang1dang1}{9,6}{⽬,⼹}[HSK 3]
  \definition{adj.}{adequado; apropriado}
  \definition{adv.}{bastante; razoavelmente; consideravelmente; indica um grau relativamente alto e profundo}
  \definition{v.}{combinar; equilibrar; corresponder a; ser aproximadamente igual a; ser proporcional a}
\end{EntryWithPhonetic}

\begin{EntryWithPhonetic}{相当于}{xiang1dang1yu2}{9,6,3}{⽬,⼹,⼆}[HSK 7-9]
  \definition{v.}{ser equivalente a; a situação de uma das partes é igual ou semelhante à da outra parte}
\end{EntryWithPhonetic}

\begin{EntryWithPhonetic}{相等}{xiang1deng3}{9,12}{⽬,⽵}[HSK 5]
  \definition{v.}{ser igual a; possuir a mesma quantidade, peso, tamanho e grau}
\end{EntryWithPhonetic}

\begin{EntryWithPhonetic}{相对}{xiang1dui4}{9,5}{⽬,⼨}[HSK 7-9]
  \definition{adj.}{relativo (a determinadas condições); existente, dependendo de certas condições; variável de acordo com certas condições | comparativo; relativo; uma situação não é independente nem absoluta, mas varia dependendo das condições; ela é derivada por meio da comparação}
  \definition{v.}{ser oposto; são opostos na natureza, como o bem e o mal, o alto e o baixo, o bonito e o feio}
  \synonymref{对应}{dui4ying4}
  \synonymref{相反}{xiang1fan3}
  \synonymref{相应}{xiang1ying5}
  \antonymref{绝对}{jue2dui4}
\end{EntryWithPhonetic}

\begin{EntryWithPhonetic}{相对而言}{xiang1dui4-er2yan2}{9,5,6,7}{⽬,⼨,⽽,⾔}[HSK 7-9]
  \definition{expr.}{em termos relativos; comparativamente falando; relativamente falando}[相对而言，他更自信。===Comparativamente falando, ele está mais confiante.]
\end{EntryWithPhonetic}

\begin{EntryWithPhonetic}{相反}{xiang1fan3}{9,4}{⽬,⼜}[HSK 4]
  \definition{adj.}{oposto; contrário; dois aspectos das coisas são contraditórios e mutuamente exclusivos}
  \definition{conj.}{pelo contrário; usado no início ou no meio de uma frase para indicar uma contradição de significado com o que foi dito anteriormente.}
\end{EntryWithPhonetic}

\begin{EntryWithPhonetic}{相辅相成}{xiang1fu3-xiang1cheng2}{9,11,9,6}{⽬,⾞,⽬,⼽}[HSK 7-9]
  \definition{expr.}{complementam"-se mutuamente; complementar"-se um ao outro; cooperem entre si}
\end{EntryWithPhonetic}

\begin{EntryWithPhonetic}{相关}{xiang1guan1}{9,6}{⽬,⼋}[HSK 3]
  \definition{v.}{estar mutuamente relacionado; estar intimamente relacionado; estar inter"-relacionado}
\end{EntryWithPhonetic}

\begin{EntryWithPhonetic}{相互}{xiang1hu4}{9,4}{⽬,⼆}[HSK 3]
  \definition{adj.}{mútuo; recíproco; entre duas pessoas ou coisas}
  \definition{adv.}{mutuamente; um ao outro; tratamento recíproco}
\end{EntryWithPhonetic}

\begin{EntryWithPhonetic}{相继}{xiang1ji4}{9,10}{⽬,⽷}[HSK 7-9]
  \definition{adv.}{em sucessão; um após o outro}
\end{EntryWithPhonetic}

\begin{EntryWithPhonetic}{相聚}{xiang1ju4}{9,14}{⽬,⽿}
  \definition{v.}{reunir; juntar; congregar}
  \synonymref{汇集}{hui4ji2}
  \synonymref{汇聚}{hui4ju4}
  \synonymref{集中}{ji2zhong1}
  \synonymref{聚会}{ju4hui4}
  \synonymref{联合}{lian2he2}
  \synonymref{相约}{xiang1yue1}
  \antonymref{分开}{fen1/kai1}
  \antonymref{分手}{fen1/shou3}
  \antonymref{分离}{fen1li2}
  \antonymref{告别}{gao4/bie2}
  \antonymref{孤单}{gu1dan1}
\end{EntryWithPhonetic}

\begin{EntryWithPhonetic}{相连}{xiang1lian2}{9,7}{⽬,⾡}[HSK 7-9]
  \definition{v.}{estar unido; estar ligado; estar conectado}
  \synonymref{接连}{jie1lian2}
  \synonymref{连接}{lian2jie1}
  \synonymref{连续}{lian2xu4}
  \synonymref{衔接}{xian2jie1}
\end{EntryWithPhonetic}

\begin{EntryWithPhonetic}{相亲}{xiang1qin1}{9,9}{⽬,⼇}
  \definition{v.}{ir em um encontro às cegas; avaliar um possível parceiro em um encontro combinado | antes do noivado, o homem e a mulher, juntamente com seus pais, se encontram para se conhecerem e visitarem a pessoa que lhes foi apresentada}
  \antonymref{分离}{fen1li2}
\end{EntryWithPhonetic}

\begin{EntryWithPhonetic}{相识}{xiang1shi2}{9,7}{⽬,⾔}[HSK 7-9]
  \definition[个]{s.}{conhecido; pessoas que se conhecem}
  \definition{v.}{conhecer"-se mutuamente; conhecer"-se melhor}
  \synonymref{重逢}{chong2feng2}
  \synonymref{遇到}{yu4/dao4}
  \synonymref{遇见}{yu4/jian5}
  \antonymref{陌生}{mo4sheng1}
\end{EntryWithPhonetic}

\begin{EntryWithPhonetic}{相思病}{xiang1si1bing4}{9,9,10}{⽬,⼼,⽧}
  \definition{s.}{mal de amor; desilusão amorosa; doença de amor}
\end{EntryWithPhonetic}

\begin{EntryWithPhonetic}{相似}{xiang1si4}{9,6}{⽬,⼈}[HSK 3]
  \definition{v.}{assemelhar"-se; ser semelhante; ser parecido}
\end{EntryWithPhonetic}

\begin{EntryWithPhonetic}{相提并论}{xiang1ti2-bing4lun4}{9,12,6,6}{⽬,⼿,⼲,⾔}[HSK 7-9]
  \definition{expr.}{colocar em uma barra; mencionar na mesma frase; comparar; contrastar; juntar; ser paralelo; discutir; tratar; discutir ou observar duas pessoas ou duas coisas juntas}
  \synonymref{一概而论}{yi2gai4'er2lun4}
\end{EntryWithPhonetic}

\begin{EntryWithPhonetic}{相通}{xiang1tong1}{9,10}{⽬,⾡}[HSK 7-9]
  \definition{v.}{estar interligado; estar conectado; as coisas estão interligadas e comunicadas}
  \synonymref{沟通}{gou1tong1}
  \synonymref{雷同}{lei2tong2}
  \synonymref{相似}{xiang1si4}
  \synonymref{相同}{xiang1tong2}
  \synonymref{一样}{yi2yang4}
\end{EntryWithPhonetic}

\begin{EntryWithPhonetic}{相同}{xiang1tong2}{9,6}{⽬,⼝}[HSK 2]
  \definition{adj.}{semelhante; similar; igual; idêntico; o mesmo; consistentes entre si, sem diferença}
  \seealsoref{相通}{xiang1tong1}
  \seealsoref{一样}{yi2yang4}
  \synonymref{雷同}{lei2tong2}
  \synonymref{同样}{tong2yang4}
  \synonymref{相似}{xiang1si4}
  \synonymref{相通}{xiang1tong1}
  \synonymref{一律}{yi2lv4}
  \synonymref{一样}{yi2yang4}
  \synonymref{一致}{yi2zhi4}
  \antonymref{不同}{bu4tong2}
  \antonymref{差异}{cha1yi4}
  \antonymref{分歧}{fen1qi2}
  \antonymref{区别}{qu1bie2}
  \antonymref{相反}{xiang1fan3}
\end{EntryWithPhonetic}

\begin{EntryWithPhonetic}{相信}{xiang1xin4}{9,9}{⽬,⼈}[HSK 2]
  \definition{v.}{acreditar em; estar convencido de; ter fé em; acreditar que algo é certo ou verdadeiro sem dúvida}
  \seealsoref{信任}{xin4ren4}
  \synonymref{坚信}{jian1xin4}
  \synonymref{肯定}{ken3ding4}
  \synonymref{确信}{que4xin4}
  \synonymref{信赖}{xin4lai4}
  \synonymref{信任}{xin4ren4}
  \antonymref{猜忌}{cai1ji4}
  \antonymref{猜疑}{cai1yi2}
  \antonymref{怀疑}{huai2yi2}
  \antonymref{疑惑}{yi2huo4}
  \antonymref{质疑}{zhi4yi2}
\end{EntryWithPhonetic}

\begin{EntryWithPhonetic}{相依为命}{xiang1yi1-wei2ming4}{9,8,4,8}{⽬,⼈,⼂,⼝}[HSK 7-9]
  \definition{expr.}{viver dependendo uns dos outros; depender uns dos outros para sobreviver; nenhum de nós consegue viver sem o outro}
\end{EntryWithPhonetic}

\begin{EntryWithPhonetic}{相宜}{xiang1yi2}{9,8}{⽬,⼧}
  \definition{adj.}{adequado; apropriado; conveniente}
  \synonymref{合适}{he2shi4}
  \synonymref{适合}{shi4he2}
  \synonymref{事宜}{shi4yi2}
  \synonymref{适宜}{shi4yi2}
\end{EntryWithPhonetic}

\begin{EntryWithPhonetic}{相应}{xiang1ying5}{9,7}{⽬,⼴}[HSK 5]
  \definition{adj.}{Dialeto: barato}
  \definition{v.}{corresponder}
\end{EntryWithPhonetic}

\begin{EntryWithPhonetic}{相遇}{xiang1yu4}{9,12}{⽬,⾡}[HSK 7-9]
  \definition{v.}{encontrar (reunião, encontro, etc.); deparar"-se com; se encontrar}
  \synonymref{碰见}{peng4jian4}
  \antonymref{错过}{cuo4guo4}
  \antonymref{分离}{fen1li2}
  \antonymref{告别}{gao4/bie2}
\end{EntryWithPhonetic}

\begin{EntryWithPhonetic}{相约}{xiang1yue1}{9,6}{⽬,⽷}[HSK 7-9]
  \definition{v.}{concordar (com o local da reunião, data, etc.); chegar a um acordo; marcar uma reunião ou consulta | combinar (um local de encontro, data, etc.)}
  \synonymref{相聚}{xiang1ju4}
\end{EntryWithPhonetic}

%%%%%%%%%% 香 %%%%%%%%%%
\subsection*{香}\addcontentsline{loh}{figure}{香 \dpy{xiang1}}

\begin{EntryWithPhonetic}{香}{xiang1}{9}{⾹}[HSK 3][Kangxi 186]
  \definition*{s.}{Sobrenome: Xiang}
  \definition{adj.}{aromático; perfumado; fragrante; cheiroso | saboroso; saboroso; delicioso; apetitoso | com gosto; com bom apetite | (sono) profundo; dormir confortavelmente e tranquilamente | popular; valorizado; apreciado}
  \definition[根,炷]{s.}{especiaria; perfume; fragrância; aromatizante; substância com aroma intenso | incenso; bastão de incenso; tiras finas feitas de serragem e especiarias, queimadas em rituais para honrar os antepassados ou deuses e budas, e também usadas para afastar odores desagradáveis ou mosquitos| antigamente, referia"-se a coisas relacionadas com mulheres}
  \antonymref{臭}{chou4}
\end{EntryWithPhonetic}

\begin{EntryWithPhonetic}{香槟酒}{xiang1bin1jiu3}{9,14,10}{⾹,⽊,⾣}
  \definition[瓶,杯,口]{s.}{Empréstimo Linguístico: \emph{champagne}; espumante}
  \synonymref{葡萄酒}{pu2tao2jiu3}
\end{EntryWithPhonetic}

\begin{EntryWithPhonetic}{香波}{xiang1bo1}{9,8}{⾹,⽔}
  \definition{s.}{Empréstimo Linguístico: xampu}
  \seealsoref{洗发皂}{xi3fa4 zao4}
\end{EntryWithPhonetic}

\begin{EntryWithPhonetic}{香肠}{xiang1chang2}{9,7}{⾹,⾁}[HSK 5]
  \definition[根]{s.}{salsicha; linguiça; alimento feito com intestino de porco, recheado com carne picada e temperos}
\end{EntryWithPhonetic}

\begin{EntryWithPhonetic}{香港}{xiang1gang3}{9,12}{⾹,⽔}
  \definition*{s.}{Hong Kong}
  \seealsoref{香港岛}{xiang1gang3dao3}
\end{EntryWithPhonetic}

\begin{EntryWithPhonetic}{香港岛}{xiang1gang3dao3}{9,12,7}{⾹,⽔,⼭}
  \definition*{s.}{Ilha de Hong Kong}
  \seealsoref{香港}{xiang1gang3}
\end{EntryWithPhonetic}

\begin{EntryWithPhonetic}{香蕉}{xiang1jiao1}{9,15}{⾹,⾋}[HSK 3]
  \definition[枝,根,个,把,串,束,弓]{s.}{banana}
\end{EntryWithPhonetic}

\begin{EntryWithPhonetic}{香料}{xiang1liao4}{9,10}{⾹,⽃}[HSK 7-9]
  \definition[种]{s.}{perfume; especiaria; condimento; substâncias orgânicas que emitem fragrância à temperatura ambiente}[这种香料可以让菜更美味。===Este tempero pode deixar os pratos mais saborosos.]
\end{EntryWithPhonetic}

\begin{EntryWithPhonetic}{香炉}{xiang1lu2}{9,8}{⾹,⽕}
  \definition{s.}{um incensário (para queimar incenso); incensário | queimador de incenso | turíbulo}
\end{EntryWithPhonetic}

\begin{EntryWithPhonetic}{香气}{xiang1qi4}{9,4}{⾹,⽓}
  \definition{s.}{cheiro doce; fragrância; aroma | fragrância; os produtos cosméticos são feitos de fragrâncias, álcool e água destilada, entre outras coisas}
  \synonymref{香味}{xiang1wei4}
  \synonymref{直到}{zhi2dao4}
  \antonymref{臭气}{chou4qi4}
\end{EntryWithPhonetic}

\begin{EntryWithPhonetic}{香水}{xiang1shui3}{9,4}{⾹,⽔}[HSK 7-9]
  \definition[瓶,款,种,滴]{s.}{aroma; perfume; cosméticos feitos com fragrâncias, álcool, água destilada, etc.}
\end{EntryWithPhonetic}

\begin{EntryWithPhonetic}{香味}{xiang1wei4}{9,8}{⾹,⼝}[HSK 7-9]
  \definition[股,阵]{s.}{aroma; odor; cheiro doce; fragrância; o termo se refere ao conjunto de sentidos do olfato e do paladar que evoca uma sensação agradável e confortável; é percebido através dos sentidos do olfato e do paladar}
  \synonymref{香气}{xiang1qi4}
\end{EntryWithPhonetic}

\begin{EntryWithPhonetic}{香蕈}{xiang1xun4}{9,15}{⾹,⾋}
  \definition{s.}{cogumelo perfumado (\emph{Lentinus edodes}); fungo shiitake, um cogumelo comestível}
\end{EntryWithPhonetic}

\begin{EntryWithPhonetic}{香烟}{xiang1yan1}{9,10}{⾹,⽕}[HSK 7-9]
  \definition[根,支,盒,包]{s.}{a fumaça da queima de incenso | cigarro}
  \synonymref{烟草}{yan1cao3}
  \synonymref{纸烟}{zhi3yan1}
\end{EntryWithPhonetic}

\begin{EntryWithPhonetic}{香艳}{xiang1yan4}{9,10}{⾹,⾊}
  \definition{adj.}{sexy; erótico; pornográfico | florido; de boudoir (isto é, de meninas ou mulheres); com linguagem rebuscada ou conteúdo relacionado aos aposentos femininos}
  \synonymref{艳丽}{yan4li4}
  \antonymref{平凡}{ping2fan2}
  \antonymref{朴素}{pu3su4}
\end{EntryWithPhonetic}

\begin{EntryWithPhonetic}{香油}{xiang1you2}{9,8}{⾹,⽔}[HSK 7-9]
  \definition{s.}{óleo perfumado | óleo de gergelim}
\end{EntryWithPhonetic}

\begin{EntryWithPhonetic}{香皂}{xiang1zao4}{9,7}{⾹,⽩}
  \definition[块,盒]{s.}{sabonete de toalete; sabonete perfumado; sabonete aromatizado; o sabonete feito com a adição de fragrância a matérias"-primas refinadas é usado principalmente para lavar o rosto}
  \synonymref{肥皂}{fei2zao4}
\end{EntryWithPhonetic}

%%%%%%%%%% 湘 %%%%%%%%%%
\subsection*{湘}\addcontentsline{loh}{figure}{湘 \dpy{xiang1}}

\begin{EntryWithPhonetic}{湘}{xiang1}{12}{⽔}
  \definition*{s.}{Abreviação de rio Xiangjiang (湘江) na província de Hunan (湖南) | Abreviação de Hunan (província de Hunan, 湖南)}
  \seealsoref{湖南}{hu2nan2}
  \seealsoref{湘江}{xiang1jiang1}
\end{EntryWithPhonetic}

\begin{EntryWithPhonetic}{湘江}{xiang1jiang1}{12,6}{⽔,⽔}
  \definition{s.}{Rio Xiangjiang, nasce em Guangxi (广西) e deságua no Lago Dongting (洞庭)}
  \seealsoref{洞庭}{dong4ting2}
  \seealsoref{广西}{guang3xi1}
\end{EntryWithPhonetic}

\begin{EntryWithPhonetic}{湘西}{xiang1xi1}{12,6}{⽔,⾑}
  \definition*{s.}{Prefeitura Autônoma de Xiangxi Tujia (土家) e Miao (苗) | oeste de Hunan (湖南)}
  \seealsoref{湖南}{hu2nan2}
  \seealsoref{苗}{miao2}
  \seealsoref{土家}{tu3jia1}
\end{EntryWithPhonetic}

%%%%%%%%%% 箱 %%%%%%%%%%
\subsection*{箱}\addcontentsline{loh}{figure}{箱 \dpy{xiang1}}

\begin{EntryWithPhonetic}{箱}{xiang1}{15}{⾋}[HSK 4]
  \definition{s.}{caixa; estojo; baú | qualquer coisa no formato de caixa}
\end{EntryWithPhonetic}

\begin{EntryWithPhonetic}{箱子}{xiang1zi5}{15,3}{⾋,⼦}[HSK 4]
  \definition[个,只]{s.}{baú; caixa; estojo; maleta; pasta executiva}
\end{EntryWithPhonetic}

%%%%%%%%%% 镶 %%%%%%%%%%
\subsection*{镶}\addcontentsline{loh}{figure}{镶 \dpy{xiang1}}

\begin{EntryWithPhonetic}{镶}{xiang1}{22}{⾦}[HSK 7-9]
  \definition{v.}{para incrustar; cravar; montar | marginar; orlar; rodear; decorar; colocar uma moldura | inserir; integrar}
\end{EntryWithPhonetic}

\begin{EntryWithPhonetic}{镶嵌}{xiang1qian4}{22,12}{⾦,⼭}[HSK 7-9]
  \definition{v.}{incrustar; cravar; montar; incorporar um objeto em outro}
\end{EntryWithPhonetic}

%%%%%%%%%% 详 %%%%%%%%%%
\subsection*{详}\addcontentsline{loh}{figure}{详 \dpy{xiang2}}

\begin{EntryWithPhonetic}{详}{xiang2}{8}{⾔}
  \definition{adj.}{conhecido; reconhecido; saber claramente | detalhado; minucioso; pormenorizado}
  \definition{s.}{detalhes; particularidades}
  \definition{v.}{contar; explicar; elaborar | saber claramente}
  \antonymref{略}{lve4}
\end{EntryWithPhonetic}

\begin{EntryWithPhonetic}{详尽}{xiang2jin4}{8,6}{⾔,⼫}[HSK 7-9]
  \definition{adj.}{detalhado e completo; exaustivo; minucioso | detalhado; amplo; preciso}
  \synonymref{精确}{jing1que4}
  \synonymref{精细}{jing1xi4}
  \synonymref{具体}{ju4ti3}
  \synonymref{细致}{xi4zhi4}
  \synonymref{详细}{xiang2xi4}
  \synonymref{周密}{zhou1mi4}
  \synonymref{仔细}{zi3xi4}
  \antonymref{粗略}{cu1lve4}
  \antonymref{概括}{gai4kuo4}
  \antonymref{简要}{jian3yao4}
  \antonymref{省略}{sheng3lve4}
\end{EntryWithPhonetic}

\begin{EntryWithPhonetic}{详细}{xiang2xi4}{8,8}{⾔,⽷}[HSK 5]
  \definition{adj.}{explícito; detalhado; minucioso; circunstancial; meticuloso}
\end{EntryWithPhonetic}

%%%%%%%%%% 祥 %%%%%%%%%%
\subsection*{祥}\addcontentsline{loh}{figure}{祥 \dpy{xiang2}}

\begin{EntryWithPhonetic}{祥}{xiang2}{10}{⽰}
  \definition*{s.}{Sobrenome: Xiang}
  \definition{adj.}{auspicioso; propício; sortudo}
\end{EntryWithPhonetic}

\begin{EntryWithPhonetic}{祥和}{xiang2he2}{10,8}{⽰,⼝}[HSK 7-9]
  \definition{adj.}{(atmosfera, humor, etc.) feliz e auspicioso | (aparência, sentimento, etc.) gentil; benigno | auspicioso e pacífico}
  \synonymref{和蔼}{he2'ai3}
  \synonymref{和谐}{he2xie2}
  \synonymref{平和}{ping2he2}
\end{EntryWithPhonetic}

%%%%%%%%%% 享 %%%%%%%%%%
\subsection*{享}\addcontentsline{loh}{figure}{享 \dpy{xiang3}}

\begin{EntryWithPhonetic}{享}{xiang3}{8}{⼇}[HSK 7-9]
  \definition{v.}{aproveitar}
\end{EntryWithPhonetic}

\begin{EntryWithPhonetic}{享受}{xiang3shou4}{8,8}{⼇,⼜}[HSK 5]
  \definition{v.}{aproveitar; desfrutar; estar satisfeito material ou espiritualmente}
  \antonymref{吃苦}{chi1/ku3}
  \antonymref{负担}{fu4dan1}
\end{EntryWithPhonetic}

\begin{EntryWithPhonetic}{享有}{xiang3you3}{8,6}{⼇,⽉}[HSK 7-9]
  \definition{v.}{desfrutar (direitos, prestígio, etc.); obter (direitos, reputação, prestígio, etc.) na sociedade}
  \synonymref{获得}{huo4de2}
  \synonymref{具有}{ju4you3}
  \synonymref{占有}{zhan4you3}
  \antonymref{剥夺}{bo1duo2}
\end{EntryWithPhonetic}

%%%%%%%%%% 响 %%%%%%%%%%
\subsection*{响}\addcontentsline{loh}{figure}{响 \dpy{xiang3}}

\begin{EntryWithPhonetic}{响}{xiang3}{9}{⼝}[HSK 2]
  \definition{adj.}{barulhento; ressonante}
  \definition[声,阵]{s.}{som; ruído; barulho | eco}
  \definition{v.}{tocar; soar; ressoar; fazer um som | soar; fazer algo emitir um som}
  \synonymref{声}{sheng1}
\end{EntryWithPhonetic}

\begin{EntryWithPhonetic}{响亮}{xiang3liang4}{9,9}{⼝,⼇}[HSK 7-9]
  \definition{adj.}{vibrante; ressonante; sonoro; ressonante; alto e claro}
  \synonymref{洪亮}{hong2liang4}
  \synonymref{嘹亮}{liao2liang4}
  \synonymref{清脆}{qing1cui4}
\end{EntryWithPhonetic}

\begin{EntryWithPhonetic}{响起}{xiang3qi3}{9,10}{⼝,⾛}[HSK 7-9]
  \definition{v.}{(som) surgir; (fonte sonora) soar; emitir som; explodir}
\end{EntryWithPhonetic}

\begin{EntryWithPhonetic}{响声}{xiang3sheng1}{9,7}{⼝,⼠}[HSK 6]
  \definition{s.}{som; ruído}
\end{EntryWithPhonetic}

\begin{EntryWithPhonetic}{响应}{xiang3ying4}{9,7}{⼝,⼴}[HSK 7-9]
  \definition{v.}{responder; dar uma resposta}
  \synonymref{答应}{da1ying5}
  \synonymref{反响}{fan3xiang3}
  \synonymref{反映}{fan3ying4}
  \synonymref{反应}{fan3ying4}
  \synonymref{呼应}{hu1ying4}
  \synonymref{回应}{hui2ying4}
  \synonymref{接受}{jie1shou4}
  \synonymref{提倡}{ti2chang4}
  \synonymref{同意}{tong2yi4}
  \synonymref{相应}{xiang1ying5}
  \antonymref{反对}{fan3dui4}
\end{EntryWithPhonetic}

%%%%%%%%%% 想 %%%%%%%%%%
\subsection*{想}\addcontentsline{loh}{figure}{想 \dpy{xiang3}}

\begin{EntryWithPhonetic}{想}{xiang3}{13}{⼼}[HSK 1]
  \definition{v.}{pensar; ponderar; refletir | supor; contar; considerar; pensar; estimar | querer; gostaria de; sentir vontade (de fazer algo) | lembrar com saudade; sentir falta}
  \seealsoref{想念}{xiang3nian4}
  \seealsoref{要}{yao4}
  \synonymref{记}{ji4}
  \synonymref{念}{nian4}
  \synonymref{思}{si1}
  \synonymref{希}{xi1}
  \synonymref{想念}{xiang3nian4}
  \synonymref{要}{yao4}
\end{EntryWithPhonetic}

\begin{EntryWithPhonetic}{想不到}{xiang3bu2dao4}{13,4,8}{⼼,⼀,⼑}[HSK 6]
  \definition{adj.}{inesperado; imprevisto}
\end{EntryWithPhonetic}

\begin{EntryWithPhonetic}{想到}{xiang3dao4}{13,8}{⼼,⼑}[HSK 2]
  \definition{v.}{pensar em; trazer à mente; ter no coração; ter uma ideia (na mente); ter uma ideia (no coração)}
  \seealsoref{料到}{liao4dao4}
  \synonymref{料到}{liao4dao4}
\end{EntryWithPhonetic}

\begin{EntryWithPhonetic}{想法}{xiang3fa5}{13,8}{⼼,⽔}[HSK 2]
  \definition[种]{s.}{ideia; opinião; pensamento; noção; o que alguém tem em mente; visões e opiniões sobre alguém ou algo obtidas através do pensamento}
  \definition{s.}{maneira de pensar | opinião | noção}
  \definition{v.}{tentar; pensar em uma maneira (de fazer algo); fazer o que puder; encontrar um jeito}
  \seealsoref{观念}{guan1nian4}
  \seealsoref{看法}{kan4fa5}
  \synonymref{点子}{dian3zi5}
  \synonymref{观点}{guan1dian3}
  \synonymref{观念}{guan1nian4}
  \synonymref{见解}{jian4jie3}
  \synonymref{念头}{nian4tou5}
  \synonymref{设法}{she4fa3}
  \synonymref{思想}{si1xiang3}
  \synonymref{主张}{zhu3zhang1}
  \synonymref{主意}{zhu3yi5}
\end{EntryWithPhonetic}

\begin{EntryWithPhonetic}{想方设法}{xiang3fang1-she4fa3}{13,4,6,8}{⼼,⽅,⾔,⽔}[HSK 7-9]
  \definition{expr.}{fazer tudo o que for possível; tentar todos os meios; tentar por bem ou por mal}
  \synonymref{千方百计}{qian1fang1-bai3ji4}
\end{EntryWithPhonetic}

\begin{EntryWithPhonetic}{想念}{xiang3nian4}{13,8}{⼼,⼼}[HSK 4]
  \definition{v.}{sentir falta; pensar em; lembrar com carinho; ficar doente por; desejar ver novamente; lembrar com saudade}
\end{EntryWithPhonetic}

\begin{EntryWithPhonetic}{想起}{xiang3qi3}{13,10}{⼼,⾛}[HSK 2]
  \definition{v.}{recordar; lembrar; pensar em; trazer à mente; cruzar pelos pensamentos de alguém; passar pelo pensamento de alguém}
\end{EntryWithPhonetic}

\begin{EntryWithPhonetic}{想想看}{xiang3xiang3kan4}{13,13,9}{⼼,⼼,⽬}
  \definition{v.}{pensar sobre isso}
\end{EntryWithPhonetic}

\begin{EntryWithPhonetic}{想象}{xiang3xiang4}{13,11}{⼼,⾗}[HSK 4]
  \definition[个,种,面]{s.}{imaginação; refere"-se ao processo mental de processamento e transformação de representações armazenadas na mente para formar novas imagens}
  \definition{v.}{imaginar; vislumbrar; visualizar; refere"-se a ter uma imagem concreta de algo que não está na frente dos olhos}
\end{EntryWithPhonetic}

%%%%%%%%%% 向 %%%%%%%%%%
\subsection*{向}\addcontentsline{loh}{figure}{向 \dpy{xiang4}}

\begin{EntryWithPhonetic}{向}{xiang4}{6}{⼝}[HSK 2]
  \definition*{s.}{Sobrenome: Xiang}
  \definition{adv.}{sempre; o tempo todo}
  \definition{prep.}{em direção a; para}
  \definition{s.}{direção | a janela voltada para o norte}
  \definition{v.}{encarar; virar"-se para | estar do lado de; ser parcial com; tomar o partido de alguém}
  \seealsoref{朝}{chao2}
  \seealsoref{对}{dui4}
  \seealsoref{往}{wang3}
  \synonymref{朝}{chao2}
  \synonymref{从}{cong2}
  \synonymref{对}{dui4}
  \synonymref{方}{fang1}
  \synonymref{近}{jin4}
  \synonymref{偏}{pian1}
  \synonymref{世}{shi4}
  \synonymref{位}{wei4}
  \synonymref{压}{ya1}
  \antonymref{背}{bei4}
  \antonymref{公}{gong1}
  \antonymref{今}{jin1}
  \antonymref{现}{xian4}
  \antonymref{正}{zheng4}
\end{EntryWithPhonetic}

\begin{EntryWithPhonetic}{向导}{xiang4dao3}{6,6}{⼝,⼨}[HSK 5]
  \definition[位]{s.}{guia; a pessoa que lidera todos e lhes indica a direção ao caminhar}
  \definition{v.}{agir como um guia; mostrar a alguém o caminho; levar alguém a algum lugar}
\end{EntryWithPhonetic}

\begin{EntryWithPhonetic}{向来}{xiang4lai2}{6,7}{⼝,⽊}[HSK 7-9]
  \definition{adv.}{sempre; o tempo todo}
  \synonymref{本来}{ben3lai2}
  \synonymref{从来}{cong2lai2}
  \synonymref{历来}{li4lai2}
  \synonymref{一贯}{yi2guan4}
  \synonymref{一向}{yi2xiang4}
  \synonymref{一直}{yi4zhi2}
  \synonymref{原来}{yuan2lai2}
  \synonymref{原先}{yuan2xian1}
\end{EntryWithPhonetic}

\begin{EntryWithPhonetic}{向前}{xiang4qian2}{6,9}{⼝,⼑}[HSK 5]
  \definition{adv.}{para frente; adiante}
  \definition{v.}{avançar; ir em direção à frente; mover"-se para frente; avançar um pouco mais}
\end{EntryWithPhonetic}

\begin{EntryWithPhonetic}{向上}{xiang4shang4}{6,3}{⼝,⼀}[HSK 5]
  \definition{adv.}{o superior; acima}
  \definition{v.}{mover"-se; subir; ir para um lugar mais alto; ir para um lugar mais alto em relação a um determinado ponto; ir para um desenvolvimento mais alto que o atual | avançar; continuar se aperfeiçoar; subir na vida; desenvolver"-se em direção ao progresso}
\end{EntryWithPhonetic}

\begin{EntryWithPhonetic}{向往}{xiang4wang3}{6,8}{⼝,⼻}[HSK 7-9]
  \definition{s.}{anseio; desejo; saudade; refere"-se a um estado psicológico de desejo ou saudade}
  \definition{v.}{ansiar por; aguardar com expectativa; desejar ou alcançar algo por amor ou admiração}
\end{EntryWithPhonetic}

\begin{EntryWithPhonetic}{向着}{xiang4zhe5}{6,11}{⼝,⽬}[HSK 7-9]
  \definition{v.}{virar"-se para; encarar | Coloquial: tomar partido de alguém; ficar do lado de; ser parcial em relação a | ser o oposto de | assumir o papel de}
  \synonymref{朝着}{chao2zhe5}
  \synonymref{靠拢}{kao4long3}
\end{EntryWithPhonetic}

%%%%%%%%%% 相 %%%%%%%%%%
\subsection*{相}\addcontentsline{loh}{figure}{相 \dpy{xiang4}}

\begin{EntryWithPhonetic}{相}{xiang4}{9}{⽬}
  \definition*{s.}{Sobrenome: Xiang}
  \definition{s.}{aparência | postura; porte; postura sentada, em pé, etc. | Física: fase; refere"-se a uma parte homogênea de uma substância com a mesma composição e as mesmas propriedades físicas e químicas | fotografia | primeiro"-ministro (na China antiga) | ministro; títulos oficiais de certos países | fácies marinha (carvão) | elefante, uma das peças do xadrez chinês | recepcionista (pessoa que ajuda o anfitrião a receber o hóspede); antigamente, referia"-se a alguém que ajudava o anfitrião a receber convidados}
  \definition{v.}{olhar e avaliar; observe a aparência das coisas; julgar sua qualidade | assistir; ajudar; auxiliar}
  \seeref{xiang1}
\end{EntryWithPhonetic}

\begin{EntryWithPhonetic}{相机}{xiang4ji1}{9,6}{⽬,⽊}[HSK 2]
  \definition[台,部,架,个]{s.}{câmera; máquina fotográfica}
  \definition{v.}{ficar atento a uma oportunidade; procurar oportunidades}
\end{EntryWithPhonetic}

\begin{EntryWithPhonetic}{相片}{xiang4pian4}{9,4}{⽬,⽚}[HSK 4]
  \definition[张]{s.}{foto; fotografia; uma imagem de uma pessoa ou objeto feita pela exposição de papel fotográfico a um negativo fotográfico e, em seguida, revelando e fixando a imagem.}
\end{EntryWithPhonetic}

\begin{EntryWithPhonetic}{相声}{xiang4sheng5}{9,7}{⽬,⼠}[HSK 5]
  \definition[个,段]{s.}{conversa cruzada; diálogo cômico; forma de performance humorística, em que os atores usam piadas, canções e imitações para satirizar e elogiar}
\end{EntryWithPhonetic}

%%%%%%%%%% 项 %%%%%%%%%%
\subsection*{项}\addcontentsline{loh}{figure}{项 \dpy{xiang4}}

\begin{EntryWithPhonetic}{项}{xiang4}{9}{⾴}[HSK 4]
  \definition*{s.}{Sobrenome: Xiang}
  \definition{clas.}{usado para itens discriminados; taxonomia}
  \definition{s.}{nuca (do pescoço); a parte de trás do pescoço | soma (de dinheiro); fundos para fins especiais | termo; em álgebra, significa uma única fórmula que não é unida por um sinal de mais ou de menos | item}
\end{EntryWithPhonetic}

\begin{EntryWithPhonetic}{项链}{xiang4lian4}{9,12}{⾴,⾦}[HSK 7-9]
  \definition[条,串,款]{s.}{colar; joias em formato de corrente, usadas ao redor do pescoço e penduradas no peito, geralmente feitas de ouro, prata ou pérolas}
\end{EntryWithPhonetic}

\begin{EntryWithPhonetic}{项目}{xiang4mu4}{9,5}{⾴,⽬}[HSK 4]
  \definition{s.}{evento; categorias em que as coisas são divididas | item; projeto; trabalhos de engenharia, acadêmicos, etc., de conteúdo específico}
\end{EntryWithPhonetic}

\begin{EntryWithPhonetic}{项羽}{xiang4 yu3}{9,6}{⾴,⽻}
  \definition*{s.}{Xiang Yu, o Conquistador (232--202 a.C.), senhor da guerra derrotado pelo primeiro imperador Han}
\end{EntryWithPhonetic}

%%%%%%%%%% 象 %%%%%%%%%%
\subsection*{象}\addcontentsline{loh}{figure}{象 \dpy{xiang4}}

\begin{EntryWithPhonetic}{象}{xiang4}{11}{⾗}
  \definition*{s.}{Sobrenome: Xiang}
  \definition[头,群,个]{s.}{elefante | elefante, uma das peças do xadrez chinês | aparência; forma; imagem}
  \definition{v.}{imitar | latir}
\end{EntryWithPhonetic}

\begin{EntryWithPhonetic}{象棋}{xiang4qi2}{11,12}{⾗,⽊}
  \definition[副]{s.}{xadrez chinês; um tipo de jogo de xadrez em que dois jogadores têm dezesseis peças cada: um general, dois soldados, dois elefantes, duas carruagens, dois cavalos, dois canhões e cinco soldados ; cada jogador joga de acordo com as regras e o vencedor é aquele que der o xeque no general do adversário}
\end{EntryWithPhonetic}

\begin{EntryWithPhonetic}{象征}{xiang4zheng1}{11,8}{⾗,⼻}[HSK 5]
  \definition[种]{s.}{símbolo; emblema; insígnia; \emph{token}; objeto concreto que simboliza um significado especial}
  \definition{v.}{simbolizar; significar; representar; expressar um significado especial através de algo concreto}
\end{EntryWithPhonetic}

%%%%%%%%%% 像 %%%%%%%%%%
\subsection*{像}\addcontentsline{loh}{figure}{像 \dpy{xiang4}}

\begin{EntryWithPhonetic}{像}{xiang4}{13}{⼈}[HSK 2,7-9]
  \definition{adv.}{como; parecer; dar a impressão de que}
  \definition[幅]{s.}{imagem; retrato; semelhança a alguém; obras baseadas no personagem | imagem}
  \definition{v.}{assemelhar"-se; ser como; parecer"-se com | ser como; ser tal como}
  \seealsoref{好像}{hao3xiang4}
  \seealsoref{似}{shi4}
  \seealsoref{似}{si4}
  \synonymref{好像}{hao3xiang4}
  \synonymref{似}{shi4}
  \synonymref{似}{si4}
\end{EntryWithPhonetic}

\begin{EntryWithPhonetic}{像样}{xiang4yang4}{13,10}{⼈,⽊}[HSK 7-9]
  \definition{adj.}{decente; apresentável; impecável; suficiente até certo ponto; atende a certos padrões}
\end{EntryWithPhonetic}

%%%%%%%%%% 橡 %%%%%%%%%%
\subsection*{橡}\addcontentsline{loh}{figure}{橡 \dpy{xiang4}}

\begin{EntryWithPhonetic}{橡}{xiang4}{15}{⽊}
  \definition{s.}{carvalho | seringueira}
\end{EntryWithPhonetic}

\begin{EntryWithPhonetic}{橡胶}{xiang4jiao1}{15,10}{⽊,⾁}[HSK 7-9]
  \definition[块,片]{s.}{borracha; vulcanizado; os produtos de borracha, feitos de látex obtido de seringueiras, capim"-borracha e outras plantas, são materiais elásticos, isolantes, impermeáveis e herméticos, amplamente utilizados em diversos aspectos da indústria e da vida cotidiana}
\end{EntryWithPhonetic}

\begin{EntryWithPhonetic}{橡皮}{xiang4pi2}{15,5}{⽊,⽪}[HSK 7-9]
  \definition[块,个,盒]{s.}{borracha; borracha maleável; um material escolar que apaga escritas a lápis | borracha; borracha"-da"-índia; o termo geral para borracha vulcanizada}
\end{EntryWithPhonetic}

%%%%%%%%%% 肖 %%%%%%%%%%
\subsection*{肖}\addcontentsline{loh}{figure}{肖 \dpy{xiao1}}

\begin{EntryWithPhonetic}{肖}{xiao1}{7}{⾁}
  \definition*{s.}{Sobrenome: Xiao}
  \seeref{xiao4}
\end{EntryWithPhonetic}

%%%%%%%%%% 削 %%%%%%%%%%
\subsection*{削}\addcontentsline{loh}{figure}{削 \dpy{xiao1}}

\begin{EntryWithPhonetic}{削}{xiao1}{9}{⼑}[HSK 7-9]
  \definition{v.}{descascar com uma faca; usar uma faca para remover a camada superficial do objeto em um ângulo | cortar (golpe de efeito); um estilo de jogar tênis de mesa}
  \seeref{xue1}
\end{EntryWithPhonetic}

%%%%%%%%%% 消 %%%%%%%%%%
\subsection*{消}\addcontentsline{loh}{figure}{消 \dpy{xiao1}}

\begin{EntryWithPhonetic}{消}{xiao1}{10}{⽔}[HSK 7-9]
  \definition{v.}{desaparecer | dissipar; remover; eliminar; fazer desaparecer | passar o tempo de forma descontraída (recreação) | precisar; tomar (necessidade, geralmente precedido por 不, 几, 何)}
  \seealsoref{不}{bu4}
  \seealsoref{何}{he2}
  \seealsoref{几}{ji3}
  \antonymref{长}{zhang3}
\end{EntryWithPhonetic}

\begin{EntryWithPhonetic}{消沉}{xiao1chen2}{10,7}{⽔,⽔}[HSK 7-9]
  \definition{adj.}{abatido; deprimido; desanimado; cabisbaixo}
  \synonymref{悲观}{bei1guan1}
  \synonymref{低落}{di1luo4}
  \synonymref{灰心}{hui1/xin1}
  \synonymref{降低}{jiang4di1}
  \synonymref{沮丧}{ju3sang4}
  \synonymref{气馁}{qi4nei3}
  \synonymref{失望}{shi1wang4}
  \synonymref{颓废}{tui2fei4}
  \synonymref{下降}{xia4jiang4}
  \synonymref{消极}{xiao1ji2}
  \antonymref{高昂}{gao1'ang2}
  \antonymref{高涨}{gao1zhang3}
  \antonymref{鼓舞}{gu3wu3}
  \antonymref{气魄}{qi4po4}
  \antonymref{兴奋}{xing1fen4}
  \antonymref{振奋}{zhen4fen4}
\end{EntryWithPhonetic}

\begin{EntryWithPhonetic}{消除}{xiao1chu2}{10,9}{⽔,⾩}[HSK 5]
  \definition{v.}{dissipar; eliminar; limpar; tornar algo inexistente; remover (algo desfavorável)}
\end{EntryWithPhonetic}

\begin{EntryWithPhonetic}{消毒}{xiao1du2}{10,9}{⽔,⽏}[HSK 5]
  \definition{v.}{desinfetar; esterilizar; matar os microrganismos causadores de doenças por meios físicos ou químicos}
\end{EntryWithPhonetic}

\begin{EntryWithPhonetic}{消防}{xiao1fang2}{10,6}{⽔,⾩}[HSK 5]
  \definition{s.}{combate a incêncios; controle de incêndios}
\end{EntryWithPhonetic}

\begin{EntryWithPhonetic}{消防员}{xiao1fang2yuan2}{10,6,7}{⽔,⾩,⼝}
  \definition[名,位,个]{s.}{bombeiro; pessoal especializado em combate a incêndios e trabalhos de resgate}
\end{EntryWithPhonetic}

\begin{EntryWithPhonetic}{消费}{xiao1fei4}{10,9}{⽔,⾙}[HSK 3]
  \definition{v.}{gastar; consumir; consumir materiais para satisfazer as necessidades de produção ou de vida (geralmente refere"-se ao consumo doméstico) | consumir (recursos naturais)}
\end{EntryWithPhonetic}

\begin{EntryWithPhonetic}{消费者}{xiao1fei4zhe3}{10,9,8}{⽔,⾙,⽼}[HSK 5]
  \definition{s.}{consumidor; cliente; consumo; indivíduos membros da sociedade que compram e utilizam bens e serviços para consumo pessoal}
\end{EntryWithPhonetic}

\begin{EntryWithPhonetic}{消耗}{xiao1hao4}{10,10}{⽔,⽾}[HSK 6]
  \definition{v.}{gastar; esgotar; consumir; usar; (espírito, força, coisas, etc.) diminuir gradualmente devido ao uso ou perda}
\end{EntryWithPhonetic}

\begin{EntryWithPhonetic}{消化}{xiao1hua4}{10,4}{⽔,⼔}[HSK 4]
  \definition{v.}{digerir (alimentos) | digerir (conhecimento); pensar e absorver; uma metáfora para a compreensão total de novos conhecimentos ou informações e a capacidade de transformá"-los em algo que possa ser usado}
\end{EntryWithPhonetic}

\begin{EntryWithPhonetic}{消极}{xiao1ji2}{10,7}{⽔,⽊}[HSK 5]
  \definition{adj.}{negativo; oposto; adverso | passivo; inativo; sem ambição; sem iniciativa; desanimado; apático}
\end{EntryWithPhonetic}

\begin{EntryWithPhonetic}{消灭}{xiao1mie4}{10,5}{⽔,⽕}[HSK 6]
  \definition{v.}{perecer; morrer; falecer; desaparecer | abolir; erradicar; eliminar; aniquilar; exterminar; acabar com; fazer com que não exista}
\end{EntryWithPhonetic}

\begin{EntryWithPhonetic}{消遣}{xiao1qian3}{10,13}{⽔,⾡}[HSK 7-9]
  \definition[个,种]{s.}{diversão; passatempo; \emph{hobby}; jogos casuais}
  \definition{v.}{divertir"-se; passar o tempo; preencher seu tempo livre ou aliviar o tédio fazendo coisas que te façam sentir relaxado e feliz}
  \synonymref{娱乐}{yu2le4}
\end{EntryWithPhonetic}

\begin{EntryWithPhonetic}{消失}{xiao1shi1}{10,5}{⽔,⼤}[HSK 3]
  \definition{v.}{desaparecer; desvanecer; dissolver; dissipar; evaporar; sumir}
\end{EntryWithPhonetic}

\begin{EntryWithPhonetic}{消息}{xiao1xi5}{10,10}{⽔,⼼}[HSK 3]
  \definition[个,条,篇,些]{s.}{notícias; informação; reportagem sobre pessoas ou situações | notícias; novidades}
\end{EntryWithPhonetic}

%%%%%%%%%% 萧 %%%%%%%%%%
\subsection*{萧}\addcontentsline{loh}{figure}{萧 \dpy{xiao1}}

\begin{EntryWithPhonetic}{萧}{xiao1}{11}{⾋}
  \definition*{s.}{Sobrenome: Xiao}
  \definition{adj.}{desolado; sombrio}
  \definition{s.}{artemísia}
\end{EntryWithPhonetic}

\begin{EntryWithPhonetic}{萧条}{xiao1tiao2}{11,7}{⾋,⽊}[HSK 7-9]
  \definition{adj.}{sombrio; triste; desolado; desolado e sem vida | depressão; a recessão econômica, ou o período imediatamente posterior a uma crise econômica cíclica na sociedade capitalista, é caracterizada pela estagnação da produção industrial, preços baixos e retração do comércio}
  \synonymref{低迷}{di1mi2}
  \synonymref{荒凉}{huang1liang2}
  \synonymref{荒芜}{huang1wu2}
  \synonymref{冷落}{leng3luo4}
  \synonymref{凄凉}{qi1liang2}
  \antonymref{繁华}{fan2hua2}
  \antonymref{繁荣}{fan2rong2}
  \antonymref{复苏}{fu4su1}
  \antonymref{兴旺}{xing1wang4}
\end{EntryWithPhonetic}

%%%%%%%%%% 销 %%%%%%%%%%
\subsection*{销}\addcontentsline{loh}{figure}{销 \dpy{xiao1}}

\begin{EntryWithPhonetic}{销}{xiao1}{12}{⾦}[HSK 7-9]
  \definition*{s.}{Sobrenome: Xiao}
  \definition{s.}{gasto; despesa | pino}
  \definition{v.}{derreter (metal) | cancelar; anular | vender; comercializar | aferrolhar; fixar; prender; pregar | fixar com um parafuso; parafusar | gastar (consumo) | inserir um pino}
  \synonymref{卖}{mai4}
  \antonymref{购}{gou4}
  \antonymref{售}{shou4}
\end{EntryWithPhonetic}

\begin{EntryWithPhonetic}{销毁}{xiao1hui3}{12,13}{⾦,⽎}[HSK 7-9]
  \definition{v.}{destruir por fusão; destruir por combustão; derreter e destruir; queimar}
  \synonymref{毁灭}{hui3mie4}
  \synonymref{歼灭}{jian1mie4}
  \synonymref{烧毁}{shao1hui3}
  \synonymref{消灭}{xiao1mie4}
  \antonymref{保存}{bao3cun2}
  \antonymref{保留}{bao3liu2}
\end{EntryWithPhonetic}

\begin{EntryWithPhonetic}{销量}{xiao1liang4}{12,12}{⾦,⾥}[HSK 7-9]
  \definition{s.}{vendas (volume); quantidade de mercadorias vendidas}
\end{EntryWithPhonetic}

\begin{EntryWithPhonetic}{销售}{xiao1shou4}{12,11}{⾦,⼝}[HSK 4]
  \definition{v.}{vender; comercializar}
\end{EntryWithPhonetic}

%%%%%%%%%% 潇 %%%%%%%%%%
\subsection*{潇}\addcontentsline{loh}{figure}{潇 \dpy{xiao1}}

\begin{EntryWithPhonetic}{潇}{xiao1}{14}{⽔}
  \definition{adj.}{Literário: (água) profunda e límpida}
  \definition{s.}{Rio Xiaoshui, um afluente do rio Xiangjiang (湘江)}
  \definition{v.}{(chuva fraca) cair | (vento e da chuva) uivar e bater forte}
  \seealsoref{湘江}{xiang1jiang1}
\end{EntryWithPhonetic}

\begin{EntryWithPhonetic}{潇洒}{xiao1sa3}{14,9}{⽔,⽔}[HSK 7-9]
  \definition{adj.}{natural e sem restrições; descreve uma pessoa como alguém que aparenta ser natural, relaxada e sem restrições | arrojado; elegante e não convencional; descreve alguém cuja personalidade, comportamento e opiniões não são restritos e que possui relativa liberdade}
  \synonymref{大方}{da4fang5}
  \synonymref{轻松}{qing1song1}
  \synonymref{自然}{zi4ran5}
\end{EntryWithPhonetic}

%%%%%%%%%% 嚣 %%%%%%%%%%
\subsection*{嚣}\addcontentsline{loh}{figure}{嚣 \dpy{xiao1}}

\begin{EntryWithPhonetic}{嚣}{xiao1}{18}{⼝}
  \definition*{s.}{Sobrenome: Xiao}
  \definition{adj.}{lazer}
  \definition{v.}{clamar; fazer barulho}
\end{EntryWithPhonetic}

\begin{EntryWithPhonetic}{嚣张}{xiao1zhang1}{18,7}{⼝,⼸}
  \definition{adj.}{desenfreado; arrogante; descontrolado; agressivo}
  \synonymref{猖狂}{chang1kuang2}
  \synonymref{疯狂}{feng1kuang2}
  \synonymref{强势}{qiang2shi4}
  \synonymref{张狂}{zhang1kuang2}
  \antonymref{沉默}{chen2mo4}
\end{EntryWithPhonetic}

%%%%%%%%%% 小 %%%%%%%%%%
\subsection*{小}\addcontentsline{loh}{figure}{小 \dpy{xiao3}}

\begin{EntryWithPhonetic}{小}{xiao3}{3}{⼩}[HSK 1,2][Kangxi 42]
  \definition*{s.}{Sobrenome: Xiao}
  \definition{adj.}{menor; pequeno; insignificante; pouco; volume, área, quantidade, intensidade, etc. não são grandes | jovem | expressões humildes, referindo"-se a si mesmo ou a pessoas ou coisas relacionadas a si mesmo | por um tempo; por um curto período; por um curto período de tempo | o mais novo; o último na ordem de antiguidade; em último lugar na classificação}
  \definition{pref.}{usado antes do sobrenome, nome, posição na família, etc.}
  \definition{s.}{os jovens; pessoas mais jovens | concubina}
  \synonymref{略}{lve4}
  \synonymref{浅}{qian3}
  \synonymref{稍}{shao1}
  \synonymref{微}{wei1}
  \synonymref{细}{xi4}
  \synonymref{纤}{xian1}
  \synonymref{幼}{you4}
  \synonymref{暂}{zan4}
  \antonymref{大}{da4}
  \antonymref{巨}{ju4}
  \antonymref{老}{lao3}
\end{EntryWithPhonetic}

\begin{EntryWithPhonetic}{小白菜}{xiao3bai2cai4}{3,5,11}{⼩,⽩,⾋}
  \definition[棵,些]{s.}{repolho chinês (\emph{Brassica chinensis}) | uma variedade de repolho chinês; \emph{pak choi} | \emph{bok choy}}
  \seealsoref{青菜}{qing1cai4}
  \seealsoref{小白菜儿}{xiao3bai2cai4r5}
\end{EntryWithPhonetic}

\begin{EntryWithPhonetic}{小白菜儿}{xiao3bai2cai4r5}{3,5,11,2}{⼩,⽩,⾋,⼉}
  \definition{s.}{uma variedade de repolho chinês; \emph{pak choi}}
\end{EntryWithPhonetic}

\begin{EntryWithPhonetic}{小吃}{xiao3chi1}{3,6}{⼩,⼝}[HSK 4]
  \definition[家]{s.}{lanche; petiscos; comida com especialidades locais, não muito para uma porção | prato frio; prato feito; cortes de frios na culinária ocidental | pratos pequenos e baratos; pratos simples em restaurantes com porções pequenas e preços baixos}
\end{EntryWithPhonetic}

\begin{EntryWithPhonetic}{小丑}{xiao3chou3}{3,4}{⼩,⼀}[HSK 7-9]
  \definition[个,名,位,帮]{s.}{palhaço; bufão; patife; na ópera tradicional chinesa, o palhaço é uma metáfora para alguém que se comporta de maneira indigna e é hábil em bajular os outros | um miserável desprezível; referindo"-se a uma pessoa insignificante}
\end{EntryWithPhonetic}

\begin{EntryWithPhonetic}{小贩}{xiao3fan4}{3,8}{⼩,⾙}[HSK 7-9]
  \definition[个,名,位,帮]{s.}{vendedor ambulante; vendedor; mascate; comerciantes com muito pouco capital}
\end{EntryWithPhonetic}

\begin{EntryWithPhonetic}{小费}{xiao3fei4}{3,9}{⼩,⾙}[HSK 6]
  \definition[笔]{s.}{gorjeta; gratificação; dinheiro extra pago por clientes e viajantes a funcionários de serviços em setores de serviços, como hotéis e pousadas}
\end{EntryWithPhonetic}

\begin{EntryWithPhonetic}{小狗}{xiao3gou3}{3,8}{⼩,⽝}
  \definition{s.}{filhote de cachorro; cães de pequeno porte, geralmente filhotes ou raças miniatura, são mantidos principalmente como animais de estimação}
\end{EntryWithPhonetic}

\begin{EntryWithPhonetic}{小孩儿}{xiao3hai2r5}{3,9,2}{⼩,⼦,⼉}[HSK 1]
  \definition[个,名,位]{s.}{criança; bebê}
  \synonymref{儿童}{er2tong2}
  \synonymref{儿子}{er2zi5}
  \synonymref{孩子}{hai2zi5}
  \synonymref{女儿}{nv3'er2}
  \synonymref{小朋友}{xiao3peng2you3}
  \synonymref{子女}{zi3nv3}
  \antonymref{成人}{cheng2ren2}
  \antonymref{大人}{da4ren2}
\end{EntryWithPhonetic}

\begin{EntryWithPhonetic}{小伙子}{xiao3huo3zi5}{3,6,3}{⼩,⼈,⼦}[HSK 4]
  \definition[位]{s.}{rapaz jovem; jovem colega}
\end{EntryWithPhonetic}

\begin{EntryWithPhonetic}{小姐}{xiao3jie3}{3,8}{⼩,⼥}[HSK 1]
  \definition[位,名,个,些]{s.}{jovem senhora; anteriormente, era assim que se referiam às filhas de famílias ricas. | senhorita; título honorífico para mulheres jovens | (gíria) prostituta}
  \seealsoref{姑娘}{gu1niang5}
  \synonymref{姑娘}{gu1niang5}
  \synonymref{女士}{nv3shi4}
\end{EntryWithPhonetic}

\begin{EntryWithPhonetic}{小看}{xiao3kan4}{3,9}{⼩,⽬}[HSK 7-9]
  \definition{v.}{subestimar; desprezar; menosprezar; desdenhar}
  \synonymref{鄙视}{bi3shi4}
  \synonymref{忽视}{hu1shi4}
  \synonymref{看不起}{kan4bu5qi3}
  \synonymref{歧视}{qi2shi4}
  \synonymref{无视}{wu2shi4}
  \antonymref{重视}{zhong4shi4}
\end{EntryWithPhonetic}

\begin{EntryWithPhonetic}{小康}{xiao3kang1}{3,11}{⼩,⼴}[HSK 7-9]
  \definition{adj.}{bem"-sucedido; (em relação ao padrão de vida) bem de vida; refere"-se à situação econômica de uma família que consegue manter um padrão de vida de classe média}
  \synonymref{富裕}{fu4yu4}
\end{EntryWithPhonetic}

\begin{EntryWithPhonetic}{小路}{xiao3lu4}{3,13}{⼩,⾜}[HSK 7-9]
  \definition[条]{s.}{caminho; trilha | atalho | faixa | estrada secundária | via}
  \antonymref{大道}{da4dao4}
\end{EntryWithPhonetic}

\begin{EntryWithPhonetic}{小麦}{xiao3mai4}{3,7}{⼩,⿆}[HSK 6]
  \definition[粒,公斤,吨,棵]{s.}{trigo}
\end{EntryWithPhonetic}

\begin{EntryWithPhonetic}{小朋友}{xiao3peng2you3}{3,8,4}{⼩,⽉,⼜}[HSK 1]
  \definition[个,名,位]{s.}{criança; crianças; refere"-se a crianças e adolescentes | (termo usado por um adulto para se dirigir a uma criança) amiguinho; menino (ou menina); termo carinhoso para se referir a crianças e adolescentes}
  \synonymref{少儿}{shao4'er2}
  \synonymref{小孩儿}{xiao3hai2r5}
  \antonymref{成人}{cheng2ren2}
  \antonymref{大人}{da4ren2}
\end{EntryWithPhonetic}

\begin{EntryWithPhonetic}{小品}{xiao3pin3}{3,9}{⼩,⼝}[HSK 7-9]
  \definition[个,部,篇]{s.}{criação literária ou artística curta e simples; ensaio; esboço; esquete | esboço curto; ato curto; uma criação literária ou artística breve e simples; originalmente referindo"-se a versões abreviadas das escrituras budistas, agora se refere a ensaios curtos ou outras formas concisas de expressão}
\end{EntryWithPhonetic}

\begin{EntryWithPhonetic}{小气鬼}{xiao3qi4gui3}{3,4,9}{⼩,⽓,⿁}
  \definition{adj.}{avarento; mesquinho; mão"-de"-vaca; miserável; pão"-duro}
\end{EntryWithPhonetic}

\begin{EntryWithPhonetic}{小气}{xiao3qi5}{3,4}{⼩,⽓}[HSK 7-9]
  \definition{adj.}{mesquinho; apertado; avarento; parcimonioso; valorizar excessivamente os próprios bens | mesquinho; pequeno; de mente fechada}
  \antonymref{大量}{da4liang4}
  \antonymref{大气}{da4qi4}
  \antonymref{大方}{da4fang5}
  \antonymref{慷慨}{kang1kai3}
\end{EntryWithPhonetic}

\begin{EntryWithPhonetic}{小区}{xiao3qu1}{3,4}{⼩,⼖}[HSK 7-9]
  \definition[个]{s.}{bairro; conjunto habitacional; comunidade residencial; áreas residenciais relativamente independentes na cidade, com instalações de serviços bem equipadas}
  \synonymref{社区}{she4qu1}
\end{EntryWithPhonetic}

\begin{EntryWithPhonetic}{小曲}{xiao3qu1}{3,6}{⼩,⽈}[HSK 7-9]
  \definition[首,支,段]{s.}{levedura para fabricação de vinho de arroz ou fermentação de arroz glutinoso | balada; melodia | melodia folclórica | música popular}
\end{EntryWithPhonetic}

\begin{EntryWithPhonetic}{小人}{xiao3ren2}{3,2}{⼩,⼈}[HSK 7-9]
  \definition[个]{s.}{pessoa mesquinha; homem pequeno; homem ordinário; uma pessoa de caráter desprezível | ``seu humilde servo''; pessoa de baixa posição social; eu (usado por um homem de posição social inferior ao falar com seus superiores) | na antiguidade, referia"-se a pessoas de baixa condição social; era frequentemente usado como autorreferência | vilão; pessoa desprezível; personagem vil; pessoas imorais}
  \antonymref{君子}{jun1zi3}
\end{EntryWithPhonetic}

\begin{EntryWithPhonetic}{小声}{xiao3sheng1}{3,7}{⼩,⼠}[HSK 2]
  \definition{v.}{falar em voz baixa; falar baixinho; sussurar}
\end{EntryWithPhonetic}

\begin{EntryWithPhonetic}{小时}{xiao3shi2}{3,7}{⼩,⽇}[HSK 1]
  \definition{clas.}{hora; unidade de medida legal do tempo, 1 hora equivale a 60 minutos, é 1/24 de um dia}
  \definition[个,半,些]{s.}{hora; refere"-se a um período de uma hora}
  \seealsoref{点钟}{dian3zhong1}
  \synonymref{点钟}{dian3zhong1}
  \synonymref{时刻}{shi2ke4}
  \synonymref{钟头}{zhong1tou2}
\end{EntryWithPhonetic}

\begin{EntryWithPhonetic}{小时候}{xiao3shi2hou5}{3,7,10}{⼩,⽇,⼈}[HSK 2]
  \definition{s.}{na infância; quando alguém era jovem; refere"-se à infância}
  \synonymref{童年}{tong2nian2}
  \antonymref{老年}{lao3nian2}
  \antonymref{中年}{zhong1nian2}
\end{EntryWithPhonetic}

\begin{EntryWithPhonetic}{小树}{xiao3shu4}{3,9}{⼩,⽊}
  \definition[棵]{s.}{muda; árvore jovem; arbusto}
\end{EntryWithPhonetic}

\begin{EntryWithPhonetic}{小说}{xiao3shuo1}{3,9}{⼩,⾔}[HSK 2]
  \definition[本,部,篇,章]{s.}{história; romance; ficção; uma forma literária que reflete a vida social por meio da descrição de personagens, ambiente e enredo}
  \synonymref{故事}{gu4shi5}
  \antonymref{纪实}{ji4shi2}
\end{EntryWithPhonetic}

\begin{EntryWithPhonetic}{小提琴}{xiao3ti2qin2}{3,12,12}{⼩,⼿,⽟}[HSK 7-9]
  \definition[个,些,把]{s.}{violino; um tipo de violino, o menor em tamanho e com o som mais agudo}
\end{EntryWithPhonetic}

\begin{EntryWithPhonetic}{小偷儿}{xiao3tou1er5}{3,11,2}{⼩,⼈,⼉}[HSK 5]
  \definition{s.}{ladrão insignificante (ou furtivo); ladrãozinho | ladrão}
\end{EntryWithPhonetic}

\begin{EntryWithPhonetic}{小腿}{xiao3tui3}{3,13}{⼩,⾁}
  \definition{s.}{perna; parte inferior da perna (do joelho ao tornozelo)}
\end{EntryWithPhonetic}

\begin{EntryWithPhonetic}{小屋}{xiao3wu1}{3,9}{⼩,⼫}
  \definition[间]{s.}{chalé; cabana; pousada | casita; cabana; casa geminada | cabine}
\end{EntryWithPhonetic}

\begin{EntryWithPhonetic}{小溪}{xiao3xi1}{3,13}{⼩,⽔}[HSK 7-9]
  \definition[条]{s.}{riacho; córrego; um pequeno riacho raso, geralmente que flui por um vale ou floresta}
\end{EntryWithPhonetic}

\begin{EntryWithPhonetic}{小小}{xiao3xiao3}{3,3}{⼩,⼩}
  \definition{adj.}{muito pouco | muito pequeno}
\end{EntryWithPhonetic}

\begin{EntryWithPhonetic}{小心翼翼}{xiao3xin1-yi4yi4}{3,4,17,17}{⼩,⼼,⽻,⽻}[HSK 7-9]
  \definition{expr.}{cautelosamente; cuidadosamente; com muita cautela; com grande cautela; originalmente significando solene e reverente, agora é usado para descrever ações extremamente cautelosas, sem espaço para negligência}
  \antonymref{粗心大意}{cu1xin1-da4yi4}
\end{EntryWithPhonetic}

\begin{EntryWithPhonetic}{小心}{xiao3xin5}{3,4}{⼩,⼼}[HSK 2]
  \definition{adj.}{cuidadoso; atento; com cautela}
  \definition{v.}{ter cuidado; ser cauteloso; estar atento; tomar cuidado; prestar atenção}
  \seealsoref{当心}{dang1xin1}
  \seealsoref{仔细}{zi3xi4}
  \synonymref{当心}{dang1xin1}
  \synonymref{谨慎}{jin3shen4}
  \synonymref{警惕}{jing3ti4}
  \synonymref{留神}{liu2/shen2}
  \synonymref{留心}{liu2/xin1}
  \synonymref{留意}{liu2/yi4}
  \synonymref{慎重}{shen4zhong4}
  \synonymref{细心}{xi4xin1}
  \synonymref{注意}{zhu4/yi4}
  \synonymref{仔细}{zi3xi4}
  \antonymref{粗心}{cu1xin1}
  \antonymref{大意}{da4yi5}
  \antonymref{鲁莽}{lu3mang3}
\end{EntryWithPhonetic}

\begin{EntryWithPhonetic}{小型}{xiao3xing2}{3,9}{⼩,⼟}[HSK 4]
  \definition{adj.}{de tamanho pequeno; em pequena escala; miniatura; tipo pequeno; tamanho de bolso; tipo compacto}
  \definition{s.}{Mediterrâneo: escunas, pequenos veleiros de pesca ou turismo | pequeno \emph{rover} lunar (duas pessoas)}
\end{EntryWithPhonetic}

\begin{EntryWithPhonetic}{小学}{xiao3xue2}{3,8}{⼩,⼦}[HSK 1]
  \definition[所]{s.}{escola primária (ou fundamental); escolas que oferecem ensino fundamental básico | estudos filológicos; antigamente, referia"-se ao estudo da escrita, da fonética e da exegese}
  \antonymref{大学}{da4xue2}
\end{EntryWithPhonetic}

\begin{EntryWithPhonetic}{小学生}{xiao3xue2sheng5}{3,8,5}{⼩,⼦,⽣}[HSK 1]
  \definition[个,位,名,些]{s.}{aluno; estudante; estudante do sexo masculino (男); estudante do sexo feminino (女) | um aluno mais novo (do que os outros da sua turma) | (dialeto) um menino pequeno}
  \seealsoref{男}{nan2}
  \seealsoref{女}{nv3}
\end{EntryWithPhonetic}

\begin{EntryWithPhonetic}{小洋白菜}{xiao3 yang2bai2cai4}{3,9,5,11}{⼩,⽔,⽩,⾋}
  \definition{s.}{couve de bruxelas}
\end{EntryWithPhonetic}

\begin{EntryWithPhonetic}{小于}{xiao3yu2}{3,3}{⼩,⼆}[HSK 6]
  \definition{prep.}{menor que; menos que; indica que um número ou quantidade é menor que outro}
\end{EntryWithPhonetic}

\begin{EntryWithPhonetic}{小众}{xiao3 zhong4}{3,6}{⼩,⼈}
  \definition{s.}{uma elite; uma minoria; um pequeno círculo; um grupo relativamente pequeno de pessoas; um pequeno grupo (no contexto de 大众)}
  \seealsoref{大众}{da4zhong4}
\end{EntryWithPhonetic}

\begin{EntryWithPhonetic}{小卒}{xiao3zu2}{3,8}{⼩,⼗}[HSK 7-9]
  \definition{s.}{soldado de infantaria |Figurativo: ninguém | Xadrez: peão | particular | homem comum; um homem sem importância; pessoas comuns | um mero peão}
  \seealsoref{小卒子}{xiao3zu2 zi3}
  \antonymref{大王}{da4wang2}
  \antonymref{大王}{dai4wang5}
\end{EntryWithPhonetic}

\begin{EntryWithPhonetic}{小卒子}{xiao3zu2 zi3}{3,8,3}{⼩,⼗,⼦}
  \definition{s.}{peão}
\end{EntryWithPhonetic}

\begin{EntryWithPhonetic}{小组}{xiao3zu3}{3,8}{⼩,⽷}[HSK 2]
  \definition[个,名,位]{s.}{grupo; um pequeno grupo de pessoas}
  \synonymref{集体}{ji2ti3}
  \synonymref{团队}{tuan2dui4}
  \antonymref{个人}{ge4ren2}
\end{EntryWithPhonetic}

%%%%%%%%%% 晓 %%%%%%%%%%
\subsection*{晓}\addcontentsline{loh}{figure}{晓 \dpy{xiao3}}

\begin{EntryWithPhonetic}{晓}{xiao3}{10}{⽇}
  \definition{s.}{amanhecer; alvorada}
  \definition{v.}{(um dia) amanhecer; romper | saber; deixar alguém saber; dizer}
\end{EntryWithPhonetic}

\begin{EntryWithPhonetic}{晓得}{xiao3de2}{10,11}{⽇,⼻}[HSK 6]
  \definition{v.}{saber; entender}[我不晓得他在哪里。===Não sei onde ele está.]
\end{EntryWithPhonetic}

%%%%%%%%%% 孝 %%%%%%%%%%
\subsection*{孝}\addcontentsline{loh}{figure}{孝 \dpy{xiao4}}

\begin{EntryWithPhonetic}{孝}{xiao4}{7}{⼦}
  \definition*{s.}{Sobrenome: Xiao}
  \definition{adj.}{filial}
  \definition{s.}{piedade filial (ou obediência); costumes tradicionais observados por um certo período após a morte de um ancião | roupas de luto}
\end{EntryWithPhonetic}

\begin{EntryWithPhonetic}{孝敬}{xiao4jing4}{7,12}{⼦,⽁}[HSK 7-9]
  \definition{v.}{presentear; dar dinheiro ou presentes aos idosos é uma forma de demonstrar piedade filial ou respeito | demonstrar respeito filial; demonstrar respeito aos mais velhos}
  \synonymref{孝顺}{xiao4shun4}
  \antonymref{虐待}{nve4dai4}
\end{EntryWithPhonetic}

\begin{EntryWithPhonetic}{孝顺}{xiao4shun4}{7,9}{⼦,⾴}[HSK 7-9]
  \definition{adj.}{obediente; seja bom com os mais velhos da sua família, ouça"-os e cuide deles}
  \definition{v.}{demonstrar obediência; cumprir o dever de servir os pais ou outros idosos de todo o coração e atender aos seus desejos}
  \synonymref{孝敬}{xiao4jing4}
\end{EntryWithPhonetic}

%%%%%%%%%% 肖 %%%%%%%%%%
\subsection*{肖}\addcontentsline{loh}{figure}{肖 \dpy{xiao4}}

\begin{EntryWithPhonetic}{肖}{xiao4}{7}{⾁}
  \definition{v.}{assemelhar"-se a; ser semelhante a}
  \seeref{xiao1}
\end{EntryWithPhonetic}

\begin{EntryWithPhonetic}{肖像}{xiao4xiang4}{7,13}{⾁,⼈}[HSK 7-9]
  \definition[幅]{s.}{retrato; um retrato ou fotografia de uma pessoa geralmente não possui fundo; a pessoa é o assunto principal}
\end{EntryWithPhonetic}

%%%%%%%%%% 哮 %%%%%%%%%%
\subsection*{哮}\addcontentsline{loh}{figure}{哮 \dpy{xiao4}}

\begin{EntryWithPhonetic}{哮}{xiao4}{10}{⼝}
  \definition{s.}{respiração pesada; chiado}
  \definition{v.}{rugir; uivar}
\end{EntryWithPhonetic}

\begin{EntryWithPhonetic}{哮喘}{xiao4chuan3}{10,12}{⼝,⼝}
  \definition{s.}{asma; sintomas de dispneia: os pacientes sentem que a respiração está muito difícil; pneumonia, insuficiência cardíaca, bronquite crônica e outras doenças causadas por espasmo da musculatura lisa respiratória frequentemente apresentam esse sintoma}
  \definition{v.}{sofrer de asma}
\end{EntryWithPhonetic}

%%%%%%%%%% 效 %%%%%%%%%%
\subsection*{效}\addcontentsline{loh}{figure}{效 \dpy{xiao4}}

\begin{EntryWithPhonetic}{效}{xiao4}{10}{⽁}
  \definition{s.}{efeito; função | eficiência; resultado}
  \definition{v.}{imitar; seguir o exemplo de | dedicar (a energia ou a vida de alguém) a; prestar (um serviço)}
\end{EntryWithPhonetic}

\begin{EntryWithPhonetic}{效仿}{xiao4fang3}{10,6}{⽁,⼈}[HSK 7-9]
  \definition{v.}{imitar; seguir | seguir o exemplo de}
  \synonymref{模仿}{mo2fang3}
  \antonymref{创新}{chuang4xin1}
\end{EntryWithPhonetic}

\begin{EntryWithPhonetic}{效果}{xiao4guo3}{10,8}{⽁,⽊}[HSK 3]
  \definition[种,个]{s.}{efeito; resultado | efeitos sonoros; vários sons ou fenômenos naturais criados para combinar com o enredo em dramas e filmes, como vento e chuva, tiros, fogo, neve, etc.}
\end{EntryWithPhonetic}

\begin{EntryWithPhonetic}{效力}{xiao4li4}{10,2}{⽁,⼒}[HSK 7-9]
  \definition{s.}{eficácia; efeito; os efeitos positivos de uma ação ou método geralmente se referem a leis, políticas, regulamentos, etc.}
  \definition{v.}{servir; prestar um serviço a; servir ou fazer coisas pelo país, organização, etc.}
  \synonymref{成效}{cheng2xiao4}
  \synonymref{功能}{gong1neng2}
  \synonymref{功效}{gong1xiao4}
  \synonymref{听从}{ting1cong2}
  \synonymref{效率}{xiao4lv4}
  \synonymref{遵守}{zun1shou3}
  \synonymref{遵循}{zun1xun2}
  \synonymref{作用}{zuo4yong4}
\end{EntryWithPhonetic}

\begin{EntryWithPhonetic}{效率}{xiao4lv4}{10,11}{⽁,⽞}[HSK 4]
  \definition[种]{s.}{eficiência; produtividade; a quantidade de trabalho concluído por unidade de tempo}
\end{EntryWithPhonetic}

\begin{EntryWithPhonetic}{效益}{xiao4yi4}{10,10}{⽁,⽫}[HSK 7-9]
  \definition{s.}{benefício; resultado benéfico; efeitos e benefícios}
  \synonymref{效果}{xiao4guo3}
\end{EntryWithPhonetic}

\begin{EntryWithPhonetic}{效应}{xiao4ying4}{10,7}{⽁,⼴}[HSK 7-9]
  \definition[种]{s.}{efeito; resultado; impacto; a reação ou efeito social causado por uma pessoa ou coisa | efeito físico ou químico; um determinado efeito produzido por ação física ou química}
  \synonymref{反应}{fan3ying4}
  \synonymref{效果}{xiao4guo3}
  \antonymref{失效}{shi1/xiao4}
\end{EntryWithPhonetic}

%%%%%%%%%% 校 %%%%%%%%%%
\subsection*{校}\addcontentsline{loh}{figure}{校 \dpy{xiao4}}

\begin{EntryWithPhonetic}{校}{xiao4}{10}{⽊}
  \definition*{s.}{Sobrenome: Xiao}
  \definition[所]{s.}{escola | oficial de campo}
  \seeref{jiao4}
\end{EntryWithPhonetic}

\begin{EntryWithPhonetic}{校服}{xiao4fu2}{10,8}{⽊,⽉}
  \definition[套,件,身]{s.}{uniformes escolares confeccionados com tecido e estilo padronizados}
  \synonymref{制服}{zhi4fu2}
\end{EntryWithPhonetic}

\begin{EntryWithPhonetic}{校规}{xiao4gui1}{10,8}{⽊,⾒}
  \definition{s.}{regras e regulamentos escolares; regras escolares que os alunos devem seguir}
\end{EntryWithPhonetic}

\begin{EntryWithPhonetic}{校监}{xiao4jian1}{10,10}{⽊,⽫}
  \definition{s.}{diretor; superintendente escolar | supervisor (da escola)}
\end{EntryWithPhonetic}

\begin{EntryWithPhonetic}{校园}{xiao4yuan2}{10,7}{⽊,⼞}[HSK 2]
  \definition[个]{s.}{campus; pátio da escola; refere"-se a todos os terrenos e edifícios dentro da área escolar}
  \synonymref{学校}{xue2xiao4}
\end{EntryWithPhonetic}

\begin{EntryWithPhonetic}{校长}{xiao4zhang3}{10,4}{⽊,⾧}[HSK 2]
  \definition[个,位,名]{s.}{diretor; presidente; reitor; o mais alto líder administrativo e empresarial de uma escola}
\end{EntryWithPhonetic}

%%%%%%%%%% 笑 %%%%%%%%%%
\subsection*{笑}\addcontentsline{loh}{figure}{笑 \dpy{xiao4}}

\begin{EntryWithPhonetic}{笑}{xiao4}{10}{⽵}[HSK 1]
  \definition{adj.}{ridículo; engraçado; risível; hilário}
  \definition{v.}{sorrir; rir; mostrar expressão de alegria; emitir sons de alegria | ridicularizar; rir de; zombar}
  \synonymref{嘲}{chao2}
  \synonymref{接}{jie1}
  \synonymref{乐}{le4}
  \synonymref{受}{shou4}
  \antonymref{哭}{ku1}
\end{EntryWithPhonetic}

\begin{EntryWithPhonetic}{笑话儿}{xiao4hua4r5}{10,8,2}{⽵,⾔,⼉}[HSK 2]
  \definition{s.}{piada; brincadeira; gracejo}
\end{EntryWithPhonetic}

\begin{EntryWithPhonetic}{笑话}{xiao4hua5}{10,8}{⽵,⾔}[HSK 2]
  \definition[个]{s.}{piada; brincadeira; uma conversa ou história que faz as pessoas rirem; algo que as pessoas usam como piada}
  \definition{v.}{ridicularizar; zombar; rir de}
  \synonymref{嘲笑}{chao2xiao4}
  \synonymref{耻笑}{chi3xiao4}
  \synonymref{讥笑}{ji1xiao4}
  \synonymref{取笑}{qu3xiao4}
\end{EntryWithPhonetic}

\begin{EntryWithPhonetic}{笑脸}{xiao4lian3}{10,11}{⽵,⾁}[HSK 6]
  \definition{s.}{\emph{smiley}; rosto sorridente (emoji)}
\end{EntryWithPhonetic}

\begin{EntryWithPhonetic}{笑容}{xiao4rong2}{10,10}{⽵,⼧}[HSK 6]
  \definition[丝,抹,个]{s.}{sorriso; expressão sorridente; o olhar no rosto de alguém ao sorrir}
\end{EntryWithPhonetic}

\begin{EntryWithPhonetic}{笑声}{xiao4sheng1}{10,7}{⽵,⼠}[HSK 6]
  \definition{s.}{riso; risada}
\end{EntryWithPhonetic}

%%%%%%%%%% 些 %%%%%%%%%%
\subsection*{些}\addcontentsline{loh}{figure}{些 \dpy{xie1}}

\begin{EntryWithPhonetic}{些}{xie1}{8}{⼆}[HSK 4]
  \definition{adv.}{um pouco; um pouco mais; usado após um adjetivo ou parte de um verbo para indicar uma pequena quantidade, equivalente a 一点儿}
  \definition{clas.}{alguns; um pouco; denota uma quantidade indefinida}
  \seealsoref{一点儿}{yi4dian3r5}
\end{EntryWithPhonetic}

\begin{EntryWithPhonetic}{些许}{xie1xu3}{8,6}{⼆,⾔}
  \definition{adj.}{um pouco; alguns}
  \definition{num.}{um pouco; alguns; uma pequena quantidade}
  \synonymref{适量}{shi4liang4}
  \antonymref{长足}{chang2zu2}
\end{EntryWithPhonetic}

%%%%%%%%%% 楔 %%%%%%%%%%
\subsection*{楔}\addcontentsline{loh}{figure}{楔 \dpy{xie1}}

\begin{EntryWithPhonetic}{楔}{xie1}{13}{⽊}
  \definition[个]{s.}{cunha | pino; pregos de madeira; pregos de bambu}
  \definition{v.}{cunhar}
\end{EntryWithPhonetic}

\begin{EntryWithPhonetic}{楔子}{xie1zi5}{13,3}{⽊,⼦}
  \definition{s.}{cunha | pino | prólogo ou interlúdio no drama da Dinastia Yuan | prólogo em alguns romances modernos; introduções a óperas e romances | calço; chuteira; lascas de madeira inseridas nas juntas de encaixe e espiga, etc. | estaca de madeira; estaca de bambu; pregos de madeira; pregos de bambu}
\end{EntryWithPhonetic}

%%%%%%%%%% 歇 %%%%%%%%%%
\subsection*{歇}\addcontentsline{loh}{figure}{歇 \dpy{xie1}}

\begin{EntryWithPhonetic}{歇}{xie1}{13}{⽋}[HSK 5]
  \definition*{s.}{Sobrenome: Xie}
  \definition{s.}{um pouco de tempo}
  \definition{v.}{descansar; fazer uma pausa | parar (o trabalho); encerrar o expediente | dormir; ir para a cama}
\end{EntryWithPhonetic}

\begin{EntryWithPhonetic}{歇息}{xie1xi5}{13,10}{⽋,⼼}
  \definition{v.}{descansar | ir para a cama; deitar"-se para dormir; dormir}
  \seealsoref{休息}{xiu1xi5}
  \synonymref{睡觉}{shui4/jiao4}
  \synonymref{停歇}{ting2xie1}
  \synonymref{休憩}{xiu1qi4}
  \synonymref{休息}{xiu1xi5}
  \antonymref{干活}{gan4/huo2}
  \antonymref{工作}{gong1zuo4}
\end{EntryWithPhonetic}

%%%%%%%%%% 协 %%%%%%%%%%
\subsection*{协}\addcontentsline{loh}{figure}{协 \dpy{xie2}}

\begin{EntryWithPhonetic}{协}{xie2}{6}{⼗}
  \definition*{s.}{Sobrenome: Xie}
  \definition{adv.}{conjuntamente; coordenadamente; juntos}
  \definition{s.}{harmonioso}
  \definition{v.}{auxiliar; assistir; ajudar}
\end{EntryWithPhonetic}

\begin{EntryWithPhonetic}{协定}{xie2ding4}{6,8}{⼗,⼧}[HSK 7-9]
  \definition[份]{s.}{acordo; concordância; os termos serão mutuamente observados após a negociação}
  \definition{v.}{concordar com; decidir por meio de negociação}
  \synonymref{合同}{he2tong5}
  \synonymref{契约}{qi4yue1}
  \synonymref{协议}{xie2yi4}
\end{EntryWithPhonetic}

\begin{EntryWithPhonetic}{协会}{xie2hui4}{6,6}{⼗,⼈}[HSK 6]
  \definition[个]{s.}{sociedade; instituto; associação; uma organização de massa formada para promover uma causa comum}
\end{EntryWithPhonetic}

\begin{EntryWithPhonetic}{协商}{xie2shang1}{6,11}{⼗,⼝}[HSK 6]
  \definition{v.}{discutir; consultar; negociar; várias partes discutiram e decidiram em conjunto para chegar à mesma visão}
\end{EntryWithPhonetic}

\begin{EntryWithPhonetic}{协调}{xie2tiao2}{6,10}{⼗,⾔}[HSK 6]
  \definition{adj.}{coordenado; harmonioso; em sintonia}
  \definition{v.}{coordenar; concertar; integrar; harmonizar; fazer a harmonia apropriada}
\end{EntryWithPhonetic}

\begin{EntryWithPhonetic}{协同}{xie2tong2}{6,6}{⼗,⼝}[HSK 7-9]
  \definition{v.}{cooperar; trabalhar em conjunto ou em coordenação.}
  \synonymref{搭配}{da1pei4}
  \synonymref{共同}{gong4tong2}
  \synonymref{合伙}{he2huo3}
  \synonymref{互动}{hu4dong4}
  \synonymref{联合}{lian2he2}
  \synonymref{配合}{pei4he5}
\end{EntryWithPhonetic}

\begin{EntryWithPhonetic}{协议}{xie2yi4}{6,5}{⼗,⾔}[HSK 5]
  \definition[份,项]{s.}{acordo; tratado; decisão conjunta alcançada através de negociação e consulta}
  \definition{v.}{concordar em}
\end{EntryWithPhonetic}

\begin{EntryWithPhonetic}{协议书}{xie2yi4shu1}{6,5,4}{⼗,⾔,⼄}[HSK 5]
  \definition{s.}{contrato | protocolo}
\end{EntryWithPhonetic}

\begin{EntryWithPhonetic}{协助}{xie2zhu4}{6,7}{⼗,⼒}[HSK 6]
  \definition{v.}{ajudar; auxiliar; dar assistência; fornecer ajuda}
\end{EntryWithPhonetic}

\begin{EntryWithPhonetic}{协作}{xie2zuo4}{6,7}{⼗,⼈}[HSK 7-9]
  \definition{v.}{cooperar com; cooperar uns com os outros; trabalhar em conjunto para concluir (uma tarefa)}
  \synonymref{合作}{he2zuo4}
  \synonymref{互动}{hu4dong4}
  \synonymref{互助}{hu4zhu4}
  \synonymref{配合}{pei4he5}
  \synonymref{团结}{tuan2jie2}
  \synonymref{协调}{xie2tiao2}
  \antonymref{比赛}{bi3sai4}
  \antonymref{分工}{fen1gong1}
\end{EntryWithPhonetic}

%%%%%%%%%% 邪 %%%%%%%%%%
\subsection*{邪}\addcontentsline{loh}{figure}{邪 \dpy{xie2}}

\begin{EntryWithPhonetic}{邪}{xie2}{6}{⾢}[HSK 7-9]
  \definition{adj.}{maligno; herético; estranho; irregular; impróprio | amaldiçoado; enfeitiçado; perverso; pessoas supersticiosas se referem às calamidades como sendo enviadas por fantasmas e deuses}
  \definition{s.}{fator patogênico; na medicina tradicional chinesa, as doenças são abordadas como causas de fatores ambientais}
  \antonymref{歪}{wai1}
  \antonymref{正}{zheng4}
\end{EntryWithPhonetic}

\begin{EntryWithPhonetic}{邪恶}{xie2'e4}{6,10}{⾢,⼼}[HSK 7-9]
  \definition{adj.}{mal; doente; perverso; vicioso}
  \definition{s.}{mal; doente; perverso; vicioso; refere"-se a pessoas ou forças malignas}
  \synonymref{凶恶}{xiong1'e4}
  \synonymref{凶狠}{xiong1hen3}
  \synonymref{罪恶}{zui4'e4}
  \antonymref{纯洁}{chun2jie2}
  \antonymref{善良}{shan4liang2}
  \antonymref{神圣}{shen2sheng4}
  \antonymref{正义}{zheng4yi4}
  \antonymref{正直}{zheng4zhi2}
\end{EntryWithPhonetic}

%%%%%%%%%% 挟 %%%%%%%%%%
\subsection*{挟}\addcontentsline{loh}{figure}{挟 \dpy{xie2}}

\begin{EntryWithPhonetic}{挟}{xie2}{9}{⼿}
  \definition*{s.}{Sobrenome: Xie}
  \definition{v.}{segurar algo debaixo do braço; segure com o braço | coagir; forçar alguém a submeter"-se à própria vontade | abrigar (ressentimento, etc.) | contar com; confiar em | depender de}
\end{EntryWithPhonetic}

\begin{EntryWithPhonetic}{挟持}{xie2chi2}{9,9}{⼿,⼿}[HSK 7-9]
  \definition{v.}{agarrar alguém pelos braços, segurando"-o por ambos os lados; agarrar pelos dois lados | manter alguém sob coação; forçar alguém à submissão; obrigar a outra parte a obedecer}
  \synonymref{要挟}{yao1xie2}
\end{EntryWithPhonetic}

%%%%%%%%%% 斜 %%%%%%%%%%
\subsection*{斜}\addcontentsline{loh}{figure}{斜 \dpy{xie2}}

\begin{EntryWithPhonetic}{斜}{xie2}{11}{⽃}[HSK 5]
  \definition{adj.}{oblíquo; inclinado | enviesado; chanfrado; diagonal; torto; nem paralelo nem perpendicular a um plano ou linha}
  \definition{v.}{virar de lado; inclinar}
\end{EntryWithPhonetic}

\begin{EntryWithPhonetic}{斜阳}{xie2yang2}{11,6}{⽃,⾩}
  \definition{s.}{sol poente}
  \synonymref{落日}{luo4ri4}
  \synonymref{夕阳}{xi1yang2}
  \antonymref{旭日}{xu4ri4}
\end{EntryWithPhonetic}

%%%%%%%%%% 谐 %%%%%%%%%%
\subsection*{谐}\addcontentsline{loh}{figure}{谐 \dpy{xie2}}

\begin{EntryWithPhonetic}{谐}{xie2}{11}{⾔}
  \definition{adj.}{humorístico | em sintonia; em harmonia; em consonância; harmonioso}
  \definition{v.}{estar em harmonia; estar em acordo | Literário: chegar a um acordo; resolver; concordar; concluir}
\end{EntryWithPhonetic}

%%%%%%%%%% 携 %%%%%%%%%%
\subsection*{携}\addcontentsline{loh}{figure}{携 \dpy{xie2}}

\begin{EntryWithPhonetic}{携}{xie2}{13}{⼿}
  \definition{v.}{carregar; levar consigo | pegar (segurar) alguém pela mão}
\end{EntryWithPhonetic}

\begin{EntryWithPhonetic}{携带}{xie2dai4}{13,9}{⼿,⼱}[HSK 7-9]
  \definition{v.}{carregar; levar consigo}
  \seealsoref{带}{dai4}
  \synonymref{附带}{fu4dai4}
  \synonymref{领导}{ling3dao3}
  \antonymref{遗落}{yi2luo4}
\end{EntryWithPhonetic}

\begin{EntryWithPhonetic}{携手}{xie2shou3}{13,4}{⼿,⼿}[HSK 7-9]
  \definition{v.}{dar as mãos; andar de mãos dadas; juntar as mãos | cooperar; trabalhar em conjunto}
\end{EntryWithPhonetic}

%%%%%%%%%% 鞋 %%%%%%%%%%
\subsection*{鞋}\addcontentsline{loh}{figure}{鞋 \dpy{xie2}}

\begin{EntryWithPhonetic}{鞋}{xie2}{15}{⾰}[HSK 2]
  \definition[双,只]{s.}{sapatos; usado nos pés; algo que toca o chão ao caminhar; sem cano alto}
\end{EntryWithPhonetic}

%%%%%%%%%% 写 %%%%%%%%%%
\subsection*{写}\addcontentsline{loh}{figure}{写 \dpy{xie3}}

\begin{EntryWithPhonetic}{写}{xie3}{5}{⼍}[HSK 1]
  \definition{v.}{escrever | compor; escrever (como autor, repórter, etc.) | descrever; retratar | pintar; desenhar | expressar a imagem das coisas através da linguagem e da escrita | desenhar (pintura)}
  \synonymref{编写}{bian1xie3}
  \synonymref{创作}{chuang4zuo4}
  \synonymref{画}{hua4}
  \synonymref{书写}{shu1xie3}
  \synonymref{写作}{xie3zuo4}
\end{EntryWithPhonetic}

\begin{EntryWithPhonetic}{写意}{xie3yi4}{5,13}{⼍,⼼}
  \definition{s.}{estilo de pintura chinesa com pinceladas livres, caracterizado por expressões vívidas e contornos marcantes}
  \definition{v.}{sugerir (em vez de descrever em detalhes)}
  \seeref{xie4yi4}
\end{EntryWithPhonetic}

\begin{EntryWithPhonetic}{写照}{xie3zhao4}{5,13}{⼍,⽕}[HSK 7-9]
  \definition{s.}{retrato}
  \definition{v.}{retratar (uma pessoa ou personagem)}
  \synonymref{缩影}{suo1ying3}
  \synonymref{写真}{xie3zhen1}
\end{EntryWithPhonetic}

\begin{EntryWithPhonetic}{写真}{xie3zhen1}{5,10}{⼍,⼗}
  \definition{s.}{retrato; pintura de retrato | representação fiel à realidade; representação fiel}
  \definition{v.}{retratar uma pessoa; desenhar um retrato | descrever algo como é}
\end{EntryWithPhonetic}

\begin{EntryWithPhonetic}{写字}{xie3zi4}{5,6}{⼍,⼦}
  \definition{v.}{escrever (à mão); praticar caligrafia}
\end{EntryWithPhonetic}

\begin{EntryWithPhonetic}{写字匠}{xie3zi4 jiang4}{5,6,6}{⼍,⼦,⼕}
  \definition{s.}{escriba; escritor; calígrafo}
\end{EntryWithPhonetic}

\begin{EntryWithPhonetic}{写字楼}{xie3zi4lou2}{5,6,13}{⼍,⼦,⽊}[HSK 6]
  \definition{s.}{prédio de escritórios}
\end{EntryWithPhonetic}

\begin{EntryWithPhonetic}{写字台}{xie3zi4tai2}{5,6,5}{⼍,⼦,⼝}[HSK 6]
  \definition[个,张]{s.}{escrivaninha; secretária; escrivaninha de escrever; uma mesa retangular usada para escrever e trabalhar, com gavetas e algumas com pequenos armários}
\end{EntryWithPhonetic}

\begin{EntryWithPhonetic}{写作}{xie3zuo4}{5,7}{⼍,⼈}[HSK 3]
  \definition{v.}{escrever artigos; escrever livros, etc.; também se refere especificamente à criação de obras literárias}
\end{EntryWithPhonetic}

%%%%%%%%%% 血 %%%%%%%%%%
\subsection*{血}\addcontentsline{loh}{figure}{血 \dpy{xie3}}

\begin{EntryWithPhonetic}{血}{xie3}{6}{⾎}[Kangxi 143]
  \seeref{xue4}
\end{EntryWithPhonetic}

%%%%%%%%%% 写 %%%%%%%%%%
\subsection*{写}\addcontentsline{loh}{figure}{写 \dpy{xie4}}

\begin{EntryWithPhonetic}{写意}{xie4yi4}{5,13}{⼍,⼼}
  \definition{adj.}{confortável; agradável}
  \seeref{xie3yi4}
\end{EntryWithPhonetic}

%%%%%%%%%% 泄 %%%%%%%%%%
\subsection*{泄}\addcontentsline{loh}{figure}{泄 \dpy{xie4}}

\begin{EntryWithPhonetic}{泄}{xie4}{8}{⽔}[HSK 7-9]
  \definition*{s.}{Sobrenome: Xie}
  \definition{v.}{deixar sair (um fluido ou gás); descarregar; liberar | revelar (um segredo); vazar (notícias, segredos, etc.) | dar vazão a; desabafar}
  \antonymref{鼓}{gu3}
\end{EntryWithPhonetic}

\begin{EntryWithPhonetic}{泄底}{xie4di3}{8,8}{⽔,⼴}
  \definition{v.}{revelar ou expor o que está no fundo de algo | divulgar a história interna; vazar segredos}
\end{EntryWithPhonetic}

\begin{EntryWithPhonetic}{泄愤}{xie4fen4}{8,12}{⽔,⼼}
  \definition{v.}{dar vazão à raiva}
\end{EntryWithPhonetic}

\begin{EntryWithPhonetic}{泄洪}{xie4hong2}{8,9}{⽔,⽔}
  \definition{v.}{liberar água da enchente (descarga de inundação)}
\end{EntryWithPhonetic}

\begin{EntryWithPhonetic}{泄漏}{xie4lou4}{8,14}{⽔,⽔}[HSK 7-9]
  \definition{v.}{vazar; vazamentos de líquidos, gases, etc., através de buracos, rachaduras, etc. | divulgar; essa metáfora descreve o vazamento acidental de (informações secretas, etc.)}
  \synonymref{暴露}{bao4lu4}
  \synonymref{揭发}{jie1fa1}
  \synonymref{流露}{liu2lu4}
  \synonymref{透露}{tou4lu4}
  \synonymref{线路}{xian4lu4}
  \synonymref{泄露}{xie4lou4}
  \synonymref{宣泄}{xuan1xie4}
  \antonymref{保密}{bao3/mi4}
  \antonymref{保守}{bao3shou3}
\end{EntryWithPhonetic}

\begin{EntryWithPhonetic}{泄露}{xie4lou4}{8,21}{⽔,⾬}[HSK 7-9]
  \definition{v.}{vazar; deixar escapar; divulgar; revelar (um segredo, etc.) | vazar; escapar; descarregar (um fluido ou gás)}
  \synonymref{暴露}{bao4lu4}
  \synonymref{揭发}{jie1fa1}
  \synonymref{揭露}{jie1lu4}
  \synonymref{流露}{liu2lu4}
  \synonymref{透露}{tou4lu4}
  \synonymref{泄漏}{xie4lou4}
  \synonymref{宣泄}{xuan1xie4}
  \antonymref{保密}{bao3/mi4}
  \antonymref{保守}{bao3shou3}
\end{EntryWithPhonetic}

\begin{EntryWithPhonetic}{泄密}{xie4/mi4}{8,11}{⽔,⼧}[HSK 7-9]
  \definition{v.+compl.}{divulgar um segredo; trair assuntos confidenciais | revelar (contar um segredo) | revelar um segredo}
  \antonymref{保密}{bao3/mi4}
\end{EntryWithPhonetic}

\begin{EntryWithPhonetic}{泄气}{xie4/qi4}{8,4}{⽔,⽓}[HSK 7-9]
  \definition{adj.}{decepcionante | frustrante | patético}
  \definition{v.+compl.}{esvaziar; descarga de gás | desanimar; sentir"-se desanimado; ficar desalentado; perder a confiança e a coragem}
  \synonymref{灰心}{hui1/xin1}
  \synonymref{懒散}{lan3san3}
  \synonymref{气馁}{qi4nei3}
  \antonymref{奋斗}{fen4dou4}
  \antonymref{起劲}{qi3jin4}
  \antonymref{振奋}{zhen4fen4}
\end{EntryWithPhonetic}

%%%%%%%%%% 泻 %%%%%%%%%%
\subsection*{泻}\addcontentsline{loh}{figure}{泻 \dpy{xie4}}

\begin{EntryWithPhonetic}{泻}{xie4}{8}{⽔}[HSK 7-9]
  \definition{v.}{fluir rapidamente; descer impetuosamente; derramar"-se | ter fezes soltas; ter diarreia}
\end{EntryWithPhonetic}

%%%%%%%%%% 卸 %%%%%%%%%%
\subsection*{卸}\addcontentsline{loh}{figure}{卸 \dpy{xie4}}

\begin{EntryWithPhonetic}{卸}{xie4}{9}{⼙}[HSK 7-9]
  \definition{v.}{descarregar; deitar | remover; despir | livrar"-se de; esquivar"-se (de sua tarefa, responsabilidade) | remover carga ou frete; descarregar a mercadoria transportada do veículo de transporte | desmontar; desmantelar; remover; despojar; desamarrar os arreios do gado | esquivar"-se; evitar; libertar"-se; esquivar"-se da responsabilidade}
  \antonymref{装}{zhuang1}
\end{EntryWithPhonetic}

%%%%%%%%%% 契 %%%%%%%%%%
\subsection*{契}\addcontentsline{loh}{figure}{契 \dpy{xie4}}

\begin{EntryWithPhonetic}{契}{xie4}{9}{⼤}
  \definition*{s.}{Xie (ancestral da dinastia Yin, considerado ministro do Imperador Shun)}
  \seeref{qi4}
\end{EntryWithPhonetic}

%%%%%%%%%% 谢 %%%%%%%%%%
\subsection*{谢}\addcontentsline{loh}{figure}{谢 \dpy{xie4}}

\begin{EntryWithPhonetic}{谢}{xie4}{12}{⾔}
  \definition*{s.}{Sobrenome: Xie}
  \definition{v.}{agradecer | desculpar"-se; pedir desculpas; admitir a própria culpa | recusar; declinar; renunciar | murchar; perder de flores ou folhas}
\end{EntryWithPhonetic}

\begin{EntryWithPhonetic}{谢病}{xie4bing4}{12,10}{⾔,⽧}
  \definition{v.}{desculpar"-se por causa de doença | Literário: recusar"-se por motivos de doença}
\end{EntryWithPhonetic}

\begin{EntryWithPhonetic}{谢恩}{xie4'en1}{12,10}{⾔,⼼}
  \definition{v.}{(geralmente de um ministro a serviço do imperador) expressar gratidão por um favor | agradecer a alguém por um favor (especialmente do imperador ou de um oficial superior)}
\end{EntryWithPhonetic}

\begin{EntryWithPhonetic}{谢媒}{xie4mei2}{12,12}{⾔,⼥}
  \definition{v.}{(recém"-casados) expressar agradecimentos ao casamenteiro | agradecer ao casamenteiro}
\end{EntryWithPhonetic}

\begin{EntryWithPhonetic}{谢世}{xie4shi4}{12,5}{⾔,⼀}
  \definition{v.}{falecer; morrer; ir para um mundo melhor}
  \synonymref{去世}{qu4shi4}
  \synonymref{逝世}{shi4shi4}
  \antonymref{出生}{chu1sheng1}
  \antonymref{诞生}{dan4sheng1}
\end{EntryWithPhonetic}

\begin{EntryWithPhonetic}{谢天谢地}{xie4tian1xie4di4}{12,4,12,6}{⾔,⼤,⾔,⼟}
  \definition{expr.}{agradecer às estrelas (da sorte); expressar gratidão ao Céu e à Terra; ``Graças a Deus!''; ``Graças aos Céus!''; ``Graças ao Céu e à Terra!''}
\end{EntryWithPhonetic}

\begin{EntryWithPhonetic}{谢谢}{xie4xie5}{12,12}{⾔,⾔}[HSK 1]
  \definition{interj.}{``Obrigado!''}
  \definition{v.}{agradecer; agradecer a gentileza dos outros}
  \seealsoref{感激}{gan3ji1}
  \seealsoref{感谢}{gan3xie4}
  \synonymref{感激}{gan3ji1}
  \synonymref{感谢}{gan3xie4}
  \antonymref{抱怨}{bao4yuan5}
  \antonymref{埋怨}{man2yuan4}
\end{EntryWithPhonetic}

\begin{EntryWithPhonetic}{谢意}{xie4yi4}{12,13}{⾔,⼼}
  \definition{s.}{gratidão; agradecimento; um sincero agradecimento}
  \synonymref{感谢}{gan3xie4}
  \antonymref{歉意}{qian4yi4}
\end{EntryWithPhonetic}

%%%%%%%%%% 心 %%%%%%%%%%
\subsection*{心}\addcontentsline{loh}{figure}{心 \dpy{xin1}}

\begin{EntryWithPhonetic}{心}{xin1}{4}{⼼}[HSK 3][Kangxi 61]
  \definition*{s.}{Xin, uma das mansões lunares; uma das vinte e oito constelações}
  \definition[颗,个]{s.}{o coraçã; órgão que impulsiona a circulação sanguínea no corpo humano e nos vertebrados | coração; mente; sentimento; intenção; refere"-se aos órgãos do pensamento e ao pensamento, sentimentos, etc. | centro; núcleo; parte central}
\end{EntryWithPhonetic}

\begin{EntryWithPhonetic}{心爱}{xin1'ai4}{4,10}{⼼,⽖}[HSK 7-9]
  \definition{v.}{amar; valorizar; ter afeição sincera; ter um relacionamento próximo; ter afeto profundo}
  \synonymref{可爱}{ke3'ai4}
  \synonymref{亲爱}{qin1'ai4}
  \synonymref{热爱}{re4'ai4}
  \synonymref{喜爱}{xi3'ai4}
  \synonymref{喜欢}{xi3huan5}
\end{EntryWithPhonetic}

\begin{EntryWithPhonetic}{心安}{xin1'an1}{4,6}{⼼,⼧}
  \definition{adj.}{à vontade; tranquilizado}
  \seealsoref{安心}{an1xin1}
\end{EntryWithPhonetic}

\begin{EntryWithPhonetic}{心安理得}{xin1'an1-li3de2}{4,6,11,11}{⼼,⼧,⽟,⼻}[HSK 7-9]
  \definition{expr.}{sentir"-se à vontade e justificado; ter a consciência tranquila; com a mente em paz e a consciência limpa; ter a consciência limpa; não ter receio algum de algo}
  \synonymref{理直气壮}{li3zhi2-qi4zhuang4}
  \antonymref{提心吊胆}{ti2xin1-diao4dan3}
\end{EntryWithPhonetic}

\begin{EntryWithPhonetic}{心病}{xin1bing4}{4,10}{⼼,⽧}[HSK 7-9]
  \definition{s.}{preocupação; ansiedade; humor deprimido | ponto sensível; problema secreto; histórias ocultas ou dor oculta | doença cardíaca}
\end{EntryWithPhonetic}

\begin{EntryWithPhonetic}{心肠}{xin1chang2}{4,7}{⼼,⾁}[HSK 7-9]
  \definition{s.}{coração; intenção; estado emocional | preocupações; coisas que pesam na mente; pensamentos; assuntos do coração;  refere"-se a sentimentos pessoais, inquietações ou coisas que perturbam a mente | humor}
\end{EntryWithPhonetic}

\begin{EntryWithPhonetic}{心得}{xin1de2}{4,11}{⼼,⼻}[HSK 7-9]
  \definition[篇,点,项,个]{s.}{percepção; experiência; conhecimento adquirido; os princípios que se aprendem e se vivenciam no trabalho e nos estudos}
\end{EntryWithPhonetic}

\begin{EntryWithPhonetic}{心慌}{xin1/huang1}{4,12}{⼼,⼼}[HSK 7-9]
  \definition{adj.}{perturbado; nervoso; alarmado; ansioso}
  \definition{v.+compl.}{ter palpitações; ter aumento ou irregularidade na frequência cardíaca}
  \synonymref{慌乱}{huang1luan4}
  \synonymref{慌张}{huang1zhang1}
\end{EntryWithPhonetic}

\begin{EntryWithPhonetic}{心机}{xin1ji1}{4,6}{⼼,⽊}
  \definition{s.}{ideias; estratégias; pensamentos; planos; tramas}
  \synonymref{心思}{xin1si5}
  \antonymref{心绪}{xin1xu4}
\end{EntryWithPhonetic}

\begin{EntryWithPhonetic}{心急如焚}{xin1ji2-ru2fen2}{4,9,6,12}{⼼,⼼,⼥,⽕}[HSK 7-9]
  \definition{s.}{ansioso; ardendo de impaciência}
  \seealsoref{心急如火}{xin1ji2-ru2huo3}
\end{EntryWithPhonetic}

\begin{EntryWithPhonetic}{心急如火}{xin1ji2-ru2huo3}{4,9,6,4}{⼼,⼼,⼥,⽕}
  \definition{expr.}{ansioso e impaciente; ardendo de impaciência}
\end{EntryWithPhonetic}

\begin{EntryWithPhonetic}{心里}{xin1li3}{4,7}{⼼,⾥}[HSK 2]
  \definition[个]{s.}{no coração; no coração de alguém | no coração; na mente; na cabeça e no peito}
  \synonymref{内心}{nei4xin1}
  \synonymref{心胸}{xin1xiong1}
  \synonymref{心中}{xin1zhong1}
  \synonymref{胸膛}{xiong1tang2}
  \antonymref{表面}{biao3mian4}
  \antonymref{外表}{wai4biao3}
\end{EntryWithPhonetic}

\begin{EntryWithPhonetic}{心里话}{xin1li3hua4}{4,7,8}{⼼,⾥,⾔}[HSK 7-9]
  \definition[句]{s.}{os pensamentos e sentimentos mais íntimos de alguém}
\end{EntryWithPhonetic}

\begin{EntryWithPhonetic}{心理}{xin1li3}{4,11}{⼼,⽟}[HSK 4]
  \definition[个]{s.}{mentalidade; refere"-se à reflexão da mente humana sobre coisas objetivas, incluindo sensação, percepção, memória, pensamento e emoções | psicologia}
\end{EntryWithPhonetic}

\begin{EntryWithPhonetic}{心灵}{xin1ling2}{4,7}{⼼,⽕}[HSK 6]
  \definition[个,颗]{s.}{alma; coração; espírito; refere"-se ao coração, espírito, pensamentos, etc.}
\end{EntryWithPhonetic}

\begin{EntryWithPhonetic}{心灵手巧}{xin1ling2-shou3qiao3}{4,7,4,5}{⼼,⽕,⼿,⼯}[HSK 7-9]
  \definition{expr.}{inteligente e habilidoso; capaz; hábil}
\end{EntryWithPhonetic}

\begin{EntryWithPhonetic}{心目}{xin1mu4}{4,5}{⼼,⽬}[HSK 7-9]
  \definition{s.}{humor; estado de espírito; refere"-se a sentimentos na mente ou nos sentidos visuais | mente; visão mental; refere"-se a pensamentos e opiniões | memória}
  \synonymref{内心}{nei4xin1}
\end{EntryWithPhonetic}

\begin{EntryWithPhonetic}{心情}{xin1qing2}{4,11}{⼼,⼼}[HSK 2]
  \definition{s.}{humor; tom de sentimento; estado de espírito; estado emocional interior}
  \seealsoref{情绪}{qing2xu4}
  \seealsoref{心绪}{xin1xu4}
  \synonymref{感情}{gan3qing2}
  \synonymref{情感}{qing2gan3}
  \synonymref{情绪}{qing2xu4}
  \synonymref{心理}{xin1li3}
  \synonymref{心绪}{xin1xu4}
  \synonymref{心思}{xin1si5}
\end{EntryWithPhonetic}

\begin{EntryWithPhonetic}{心声}{xin1sheng1}{4,7}{⼼,⼠}[HSK 7-9]
  \definition{s.}{desejos sinceros; aspirações; pensamentos}
  \synonymref{渴望}{ke3wang4}
  \synonymref{心事}{xin1shi4}
\end{EntryWithPhonetic}

\begin{EntryWithPhonetic}{心事}{xin1shi4}{4,8}{⼼,⼅}[HSK 7-9]
  \definition[桩,件]{s.}{mágoa; preocupação; algo que pesa na mente; um fardo para a mente; as coisas em que estou pensando (geralmente me referindo a coisas que me deixam desconfortável ou constrangido)}
  \synonymref{心思}{xin1si5}
\end{EntryWithPhonetic}

\begin{EntryWithPhonetic}{心思}{xin1si5}{4,9}{⼼,⼼}[HSK 7-9]
  \definition[番]{s.}{pensamento; ideia | pensamento; refere"-se a habilidades como raciocínio e memória | humor; estado de espírito; a sensação de querer fazer algo}
  \seealsoref{心事}{xin1shi4}
  \synonymref{脑筋}{nao3jin1}
  \synonymref{情绪}{qing2xu4}
  \synonymref{思想}{si1xiang3}
  \synonymref{头脑}{tou2nao3}
  \synonymref{心情}{xin1qing2}
  \synonymref{心绪}{xin1xu4}
\end{EntryWithPhonetic}

\begin{EntryWithPhonetic}{心酸}{xin1suan1}{4,14}{⼼,⾣}[HSK 7-9]
  \definition{adj.}{de partir o coração; triste; dor de coração causada pelo luto}
  \synonymref{悲哀}{bei1'ai1}
  \synonymref{悲伤}{bei1shang1}
  \synonymref{辛酸}{xin1suan1}
  \antonymref{甜蜜}{tian2mi4}
\end{EntryWithPhonetic}

\begin{EntryWithPhonetic}{心态}{xin1tai4}{4,8}{⼼,⼼}[HSK 5]
  \definition[种,个]{s.}{mentalidade; psicologia; estado mental}
\end{EntryWithPhonetic}

\begin{EntryWithPhonetic}{心疼}{xin1teng2}{4,10}{⼼,⽧}[HSK 5]
  \definition{v.}{amar profundamente; sentir pena porque coisas valiosas foram destruídas ou perdidas; não querer se separar delas | sentir pena; ficar angustiado; preocupar"-se e sofrer pelo sofrimento dos outros; estar disposto a cuidar mais por causa da preocupação}
\end{EntryWithPhonetic}

\begin{EntryWithPhonetic}{心想事成}{xin1xiang3-shi4cheng2}{4,13,8,6}{⼼,⼼,⼅,⼽}[HSK 7-9]
  \definition{expr.}{``Que todos os seus desejos se realizem.''; ter os próprios desejos realizados; independentemente do que você imaginar, você terá sucesso}
  \synonymref{如愿以偿}{ru2yuan4yi3chang2}
  \synonymref{一帆风顺}{yi4fan1-feng1shun4}
\end{EntryWithPhonetic}

\begin{EntryWithPhonetic}{心胸}{xin1xiong1}{4,10}{⼼,⾁}[HSK 7-9]
  \definition{s.}{mentalidade; tolerância; mente aberta; amplitude de espírito | aspiração; ambição}
  \seealsoref{胸怀}{xiong1huai2}
\end{EntryWithPhonetic}

\begin{EntryWithPhonetic}{心绪}{xin1xu4}{4,11}{⼼,⽷}
  \definition{s.}{estado de espírito | sentimentos; mente; humor (geralmente descrito como estável ou desordenado)}
  \seealsoref{心情}{xin1qing2}
  \synonymref{情绪}{qing2xu4}
  \synonymref{心理}{xin1li3}
  \synonymref{心情}{xin1qing2}
  \synonymref{心思}{xin1si5}
\end{EntryWithPhonetic}

\begin{EntryWithPhonetic}{心血}{xin1xue4}{4,6}{⼼,⾎}[HSK 7-9]
  \definition{s.}{esforço meticuloso; esforço empregado para alcançar o sucesso}
  \synonymref{血汗}{xue4han4}
\end{EntryWithPhonetic}

\begin{EntryWithPhonetic}{心眼儿}{xin1yan3r5}{4,11,2}{⼼,⽬,⼉}[HSK 7-9]
  \definition{s.}{coração; mente; intenção; a mente de uma pessoa | inteligência; astúcia; engenhosidade | dúvidas infundadas; receios desnecessários; preocupações e considerações desnecessárias sobre as pessoas | tolerância}
\end{EntryWithPhonetic}

\begin{EntryWithPhonetic}{心意}{xin1yi4}{4,13}{⼼,⼼}[HSK 7-9]
  \definition{s.}{consideração; propósito; sentimentos amáveis; os sentimentos positivos em relação aos outros geralmente são expressos por meio de ajuda, agradecimento e votos de felicidades | mente; o que eu quero expressar; meus pensamentos mais íntimos}
  \seealsoref{心愿}{xin1yuan4}
\end{EntryWithPhonetic}

\begin{EntryWithPhonetic}{心愿}{xin1yuan4}{4,14}{⼼,⽕}[HSK 6]
  \definition[桩]{s.}{desejo acalentado; aspiração; desejo; sonho | o desejo do coração}
\end{EntryWithPhonetic}

\begin{EntryWithPhonetic}{心脏}{xin1zang4}{4,10}{⼼,⾁}[HSK 6]
  \definition[颗,个]{s.}{coração; um órgão importante no corpo de humanos ou animais superiores que faz o sangue circular | coração; o centro ou a parte mais importante de uma metáfora}
\end{EntryWithPhonetic}

\begin{EntryWithPhonetic}{心脏病}{xin1zang4bing4}{4,10,10}{⼼,⾁,⽧}[HSK 6]
  \definition{s.}{doença cardíaca; cardiopatia um termo geral para anormalidades ou doenças na estrutura e função do coração humano}
\end{EntryWithPhonetic}

\begin{EntryWithPhonetic}{心中}{xin1zhong1}{4,4}{⼼,⼁}[HSK 2]
  \definition{s.}{no coração; na mente}
  \synonymref{内心}{nei4xin1}
  \synonymref{心里}{xin1li3}
\end{EntryWithPhonetic}

%%%%%%%%%% 芯 %%%%%%%%%%
\subsection*{芯}\addcontentsline{loh}{figure}{芯 \dpy{xin1}}

\begin{EntryWithPhonetic}{芯}{xin1}{7}{⾋}
  \definition{s.}{medula de junco | pavio}
  \seeref{xin4}
\end{EntryWithPhonetic}

\begin{EntryWithPhonetic}{芯片}{xin1pian4}{7,4}{⾋,⽚}[HSK 7-9]
  \definition{s.}{\emph{chip} de computador; \emph{microchip}; um substrato (geralmente uma pastilha de silício) que contém um circuito integrado completo}
\end{EntryWithPhonetic}

%%%%%%%%%% 辛 %%%%%%%%%%
\subsection*{辛}\addcontentsline{loh}{figure}{辛 \dpy{xin1}}

\begin{EntryWithPhonetic}{辛}{xin1}{7}{⾟}[Kangxi 160]
  \definition*{s.}{Sobrenome: Xin}
  \definition{adj.}{quente (no sabor, etc.); pungente | difícil; trabalhoso | ponto da bússola chinesa antiga: 285° | oitavo na ordem}
  \definition{pref.}{octa-}
  \definition{s.}{sofrimento}
  \definition{s.}{oitavo dos dez caules celestiais}
\end{EntryWithPhonetic}

\begin{EntryWithPhonetic}{辛苦}{xin1ku3}{7,8}{⾟,⾋}[HSK 5]
  \definition{adj.}{difícil; trabalhoso; árduo; descreve muito trabalho, alta intensidade e pouco descanso}
  \definition{s.}{dificuldades}
  \definition{v.}{trabalhar duro; passar por grandes dificuldades; passar por dificuldades}
\end{EntryWithPhonetic}

\begin{EntryWithPhonetic}{辛勤}{xin1qin2}{7,13}{⾟,⼒}[HSK 7-9]
  \definition{adj.}{trabalhador; esforçado; diligente; que não tem medo de se esforçar}
  \synonymref{吃力}{chi1li4}
  \synonymref{艰苦}{jian1ku3}
  \synonymref{劳累}{lao2lei4}
  \synonymref{努力}{nu3li4}
  \synonymref{勤奋}{qin2fen4}
  \synonymref{勤劳}{qin2lao2}
  \synonymref{勤快}{qin2kuai5}
  \synonymref{辛苦}{xin1ku3}
  \synonymref{用功}{yong4gong1}
  \antonymref{懒惰}{lan3duo4}
\end{EntryWithPhonetic}

\begin{EntryWithPhonetic}{辛酸}{xin1suan1}{7,14}{⾟,⾣}[HSK 7-9]
  \definition{adj.}{triste; amargo; miserável}
  \definition{s.}{Química: ácido caprílico; ácido octóico; ácido octanoico}
  \synonymref{悲哀}{bei1'ai1}
  \synonymref{悲伤}{bei1shang1}
  \synonymref{心酸}{xin1suan1}
  \antonymref{甜美}{tian2mei3}
  \antonymref{甜蜜}{tian2mi4}
  \antonymref{喜悦}{xi3yue4}
\end{EntryWithPhonetic}

%%%%%%%%%% 欣 %%%%%%%%%%
\subsection*{欣}\addcontentsline{loh}{figure}{欣 \dpy{xin1}}

\begin{EntryWithPhonetic}{欣}{xin1}{8}{⽋}
  \definition*{s.}{Sobrenome: Xin}
  \definition{adj.}{alegre; feliz; contente}
\end{EntryWithPhonetic}

\begin{EntryWithPhonetic}{欣赏}{xin1shang3}{8,12}{⽋,⾙}[HSK 5]
  \definition{v.}{apreciar; admirar; valorizar; apreciar as coisas boas e descubrir o prazer que elas proporcionam | apreciar; gostar; considerar bom}
\end{EntryWithPhonetic}

\begin{EntryWithPhonetic}{欣慰}{xin1wei4}{8,15}{⽋,⼼}[HSK 7-9]
  \definition{adj.}{aliviado; satisfeito; à vontade e satisfeito; alegre e bem"-disposto; confortável e despreocupado}
  \synonymref{安抚}{an1fu3}
  \synonymref{安慰}{an1wei4}
  \synonymref{惭愧}{can2kui4}
  \synonymref{慰问}{wei4wen4}
  \synonymref{欣喜}{xin1xi3}
  \antonymref{伤感}{shang1gan3}
  \antonymref{痛心}{tong4xin1}
\end{EntryWithPhonetic}

\begin{EntryWithPhonetic}{欣喜}{xin1xi3}{8,12}{⽋,⼝}[HSK 7-9]
  \definition{adj.}{contente; feliz; alegre}
  \synonymref{沸腾}{fei4teng2}
  \synonymref{高兴}{gao1xing4}
  \synonymref{欢喜}{huan1xi3}
  \synonymref{开心}{kai1/xin1}
  \synonymref{喜悦}{xi3yue4}
  \synonymref{欣慰}{xin1wei4}
  \synonymref{愉快}{yu2kuai4}
  \antonymref{悲哀}{bei1'ai1}
  \antonymref{忧愁}{you1chou2}
\end{EntryWithPhonetic}

\begin{EntryWithPhonetic}{欣欣向荣}{xin1xin1-xiang4rong2}{8,8,6,9}{⽋,⽋,⼝,⾋}[HSK 7-9]
  \definition{expr.}{próspero; florescente; descrevendo uma vegetação exuberante e verdejante; metaforicamente, significa um empreendimento próspero e bem"-sucedido}
  \synonymref{万古长青}{wan4gu3-chang2qing1}
  \synonymref{朝气蓬勃}{zhao1qi4-peng2bo2}
\end{EntryWithPhonetic}

%%%%%%%%%% 新 %%%%%%%%%%
\subsection*{新}\addcontentsline{loh}{figure}{新 \dpy{xin1}}

\begin{EntryWithPhonetic}{新}{xin1}{13}{⽄}[HSK 1]
  \definition*{s.}{Xinjiang, abreviação de 新疆 | Singapura, abreviação de 新加坡 | Sobrenome: Xin}
  \definition{adj.}{novo; fresco; inovador; atualizado; aparecer ou ser experimentado pela primeira vez | nunca utilizado; novo; não foi usado ou foi usado por pouco tempo | recém"-casado}
  \definition{adv.}{recém; recentemente; há pouco tempo}
  \definition{pref.}{Química: meso-}
  \definition{v.}{atualizar; renovar}
  \seealsoref{新加坡}{xin1jia1po1}
  \seealsoref{新疆}{xin1jiang1}
  \synonymref{初}{chu1}
  \synonymref{刚}{gang1}
  \synonymref{始}{shi3}
  \synonymref{鲜}{xian1}
  \synonymref{属}{zhu3}
  \antonymref{陈}{chen2}
  \antonymref{故}{gu4}
  \antonymref{旧}{jiu4}
  \antonymref{老}{lao3}
\end{EntryWithPhonetic}

\begin{EntryWithPhonetic}{新潮}{xin1chao2}{13,15}{⽄,⽔}[HSK 7-9]
  \definition{adj.}{na moda; atualizado; \emph{à la mode} | nova maré; nova onda; moda; moderno}
  \definition{s.}{nova tendência; nova moda}
  \synonymref{潮流}{chao2liu2}
  \synonymref{时髦}{shi2mao2}
  \synonymref{新颖}{xin1ying3}
  \antonymref{传统}{chuan2tong3}
\end{EntryWithPhonetic}

\begin{EntryWithPhonetic}{新陈代谢}{xin1chen2-dai4xie4}{13,7,5,12}{⽄,⾩,⼈,⾔}[HSK 7-9]
  \definition{expr.}{o novo tomando o lugar do antigo; refere"-se ao processo pelo qual os organismos vivos substituem constantemente substâncias antigas por novas; de forma mais ampla, refere"-se ao surgimento e desenvolvimento contínuos de novas coisas, que substituem as antigas, que se decompõem constantemente | Biologia: metabolismo}
\end{EntryWithPhonetic}

\begin{EntryWithPhonetic}{新房}{xin1fang2}{13,8}{⽄,⼾}[HSK 7-9]
  \definition[间,套]{s.}{quarto nupcial | casa recém"-construída ou comprada; casa nova | casa novinha em folha}
  \seealsoref{新房子}{xin1fang2zi3}
\end{EntryWithPhonetic}

\begin{EntryWithPhonetic}{新房子}{xin1fang2zi3}{13,8,3}{⽄,⼾,⼦}
  \definition{s.}{casa nova}
\end{EntryWithPhonetic}

\begin{EntryWithPhonetic}{新加坡}{xin1jia1po1}{13,5,8}{⽄,⼒,⼟}
  \definition*{s.}{Singapura}
\end{EntryWithPhonetic}

\begin{EntryWithPhonetic}{新疆}{xin1jiang1}{13,19}{⽄,⼸}
  \definition*{s.}{Região Autônoma Uigur de Xinjiang}
\end{EntryWithPhonetic}

\begin{EntryWithPhonetic}{新疆维吾尔自治区}{xin1jiang1wei2wu2'er3zi4zhi4qu1}{13,19,11,7,5,6,8,4}{⽄,⼸,⽷,⼝,⼩,⾃,⽔,⼖}
  \definition*{s.}{Região Autônoma Uigur de Xinjiang, abreviada como 新, capital 乌鲁木齐}
  \seealsoref{乌鲁木齐}{wu1lu3mu4qi2}
  \seealsoref{新}{xin1}
\end{EntryWithPhonetic}

\begin{EntryWithPhonetic}{新款}{xin1kuan3}{13,12}{⽄,⽋}[HSK 7-9]
  \definition{s.}{novo estilo; nova moda; novo modelo}
\end{EntryWithPhonetic}

\begin{EntryWithPhonetic}{新郎}{xin1lang2}{13,8}{⽄,⾢}[HSK 4]
  \definition[位,名,个,些]{s.}{noivo; homens no momento do casamento}
\end{EntryWithPhonetic}

\begin{EntryWithPhonetic}{新年}{xin1nian2}{13,6}{⽄,⼲}[HSK 1]
  \definition*[个]{s.}{Ano Novo}
  \synonymref{过年}{guo4/nian2}
\end{EntryWithPhonetic}

\begin{EntryWithPhonetic}{新娘}{xin1niang2}{13,10}{⽄,⼥}[HSK 4]
  \definition[位,个]{s.}{noiva; a mulher no momento do casamento}
  \seealsoref{新娘子}{xin1niang2zi5}
\end{EntryWithPhonetic}

\begin{EntryWithPhonetic}{新娘服装}{xin1niang2 fu2zhuang1}{13,10,8,12}{⽄,⼥,⽉,⾐}
  \definition{s.}{vestido de noiva}
\end{EntryWithPhonetic}

\begin{EntryWithPhonetic}{新娘子}{xin1niang2zi5}{13,10,3}{⽄,⼥,⼦}
  \definition{s.}{noiva}
  \seealsoref{新娘}{xin1niang2}
\end{EntryWithPhonetic}

\begin{EntryWithPhonetic}{新奇}{xin1qi2}{13,8}{⽄,⼤}[HSK 7-9]
  \definition{adj.}{novo; inédito; estranho}
  \synonymref{别致}{bie2zhi4}
  \synonymref{海鲜}{hai3xian1}
  \synonymref{好奇}{hao4qi2}
  \synonymref{稀奇}{xi1qi2}
  \synonymref{新鲜}{xin1xian1}
  \synonymref{新颖}{xin1ying3}
  \antonymref{陈旧}{chen2jiu4}
\end{EntryWithPhonetic}

\begin{EntryWithPhonetic}{新人}{xin1ren2}{13,2}{⽄,⼈}[HSK 6]
  \definition[位]{s.}{pessoas de um novo tipo; nova pessoa;  pessoa que virou uma nova página | nova personalidade; novo talento | recém"-chegado; novo membro | noiva ou noivo; recém"-casado | \emph{neoanthropus}; \emph{homo sapiens}}
\end{EntryWithPhonetic}

\begin{EntryWithPhonetic}{新生}{xin1sheng1}{13,5}{⽄,⽣}[HSK 7-9]
  \definition{adj.}{recém"-nascido}
  \definition[个,名,位]{s.}{renascimento; regeneração; nova vida; uma nova oportunidade de vida}
\end{EntryWithPhonetic}

\begin{EntryWithPhonetic}{新式}{xin1shi4}{13,6}{⽄,⼷}[HSK 7-9]
  \definition{adj.}{novo (tipo); mais recente (tipo); novo (estilo) | novo estilo; moderno; neotérico; atualizado | novo modelo}
  \synonymref{新型}{xin1xing2}
\end{EntryWithPhonetic}

\begin{EntryWithPhonetic}{新手}{xin1shou3}{13,4}{⽄,⼿}[HSK 7-9]
  \definition[双]{s.}{novato; recruta inexperiente; estreante; pessoas que são novas em um determinado trabalho}
\end{EntryWithPhonetic}

\begin{EntryWithPhonetic}{新闻}{xin1wen2}{13,9}{⽄,⾨}[HSK 2]
  \definition[个,条,则,版]{s.}{notícias; notícias nacionais e internacionais reportadas em jornais, estações de rádio, etc. | notícias; refere"-se a coisas importantes ou novas que aconteceram recentemente na sociedade}
  \seealsoref{消息}{xiao1xi5}
  \synonymref{时事}{shi2shi4}
  \synonymref{消息}{xiao1xi5}
  \synonymref{信息}{xin4xi1}
  \synonymref{资讯}{zi1xun4}
\end{EntryWithPhonetic}

\begin{EntryWithPhonetic}{新鲜}{xin1xian1}{13,14}{⽄,⿂}[HSK 4]
  \definition{adj.}{fresco (experiência, alimento, etc.) | novo; inédito; estranho; raro; incomum | fresco; (o ar) é frequentemente circulado e livre de impurezas}
  \definition{s.}{frescor}
  \seealsoref{鲜}{xian1}
  \seealsoref{新颖}{xin1ying3}
  \synonymref{别致}{bie2zhi4}
  \synonymref{清新}{qing1xin1}
  \synonymref{稀奇}{xi1qi2}
  \synonymref{稀罕}{xi1han5}
  \synonymref{鲜活}{xian1huo2}
  \synonymref{鲜美}{xian1mei3}
  \synonymref{新奇}{xin1qi2}
  \synonymref{新颖}{xin1ying3}
  \antonymref{陈旧}{chen2jiu4}
  \antonymref{大众}{da4zhong4}
  \antonymref{腐败}{fu3bai4}
  \antonymref{腐烂}{fu3lan4}
  \antonymref{腐朽}{fu3xiu3}
  \antonymref{普通}{pu3tong1}
\end{EntryWithPhonetic}

\begin{EntryWithPhonetic}{新兴}{xin1xing1}{13,6}{⽄,⼋}[HSK 6]
  \definition[个]{adj.}{recém"-desenvolvido; crescente; florescente; emergente; descreve algo que está apenas começando a se tornar popular ou se desenvolver}
\end{EntryWithPhonetic}

\begin{EntryWithPhonetic}{新型}{xin1xing2}{13,9}{⽄,⼟}[HSK 4]
  \definition[种]{s.}{ultimo modelo; novo tipo; novo padrão; novo estilo}
\end{EntryWithPhonetic}

\begin{EntryWithPhonetic}{新颖}{xin1ying3}{13,13}{⽄,⾴}[HSK 7-9]
  \definition{adj.}{novo; original; único}
  \synonymref{别致}{bie2zhi4}
  \synonymref{清新}{qing1xin1}
  \synonymref{稀奇}{xi1qi2}
  \synonymref{现代}{xian4dai4}
  \synonymref{新潮}{xin1chao2}
  \synonymref{新奇}{xin1qi2}
  \synonymref{新鲜}{xin1xian1}
  \antonymref{陈旧}{chen2jiu4}
  \antonymref{古老}{gu3lao3}
  \antonymref{普通}{pu3tong1}
  \antonymref{一般}{yi4ban1}
\end{EntryWithPhonetic}

%%%%%%%%%% 薪 %%%%%%%%%%
\subsection*{薪}\addcontentsline{loh}{figure}{薪 \dpy{xin1}}

\begin{EntryWithPhonetic}{薪}{xin1}{16}{⾋}
  \definition{s.}{lenha; combustível | salário; ordenado; pagamento}
\end{EntryWithPhonetic}

\begin{EntryWithPhonetic}{薪水}{xin1shui5}{16,4}{⾋,⽔}[HSK 6]
  \definition[份,笔]{s.}{pagamento; salário; ordenados; dinheiro ou bens pagos regularmente aos trabalhadores como compensação pelo seu trabalho}
\end{EntryWithPhonetic}

%%%%%%%%%% 寻 %%%%%%%%%%
\subsection*{寻}\addcontentsline{loh}{figure}{寻 \dpy{xin2}}

\begin{EntryWithPhonetic}{寻}{xin2}{6}{⼨}
  \definition{v.}{pesquisar}
  \seeref{xun2}
  \synonymref{找}{zhao3}
\end{EntryWithPhonetic}

%%%%%%%%%% 芯 %%%%%%%%%%
\subsection*{芯}\addcontentsline{loh}{figure}{芯 \dpy{xin4}}

\begin{EntryWithPhonetic}{芯}{xin4}{7}{⾋}
  \definition{s.}{núcleo; a parte central de um objeto | língua de cobra}
  \seeref{xin1}
\end{EntryWithPhonetic}

%%%%%%%%%% 信 %%%%%%%%%%
\subsection*{信}\addcontentsline{loh}{figure}{信 \dpy{xin4}}

\begin{EntryWithPhonetic}{信}{xin4}{9}{⼈}[HSK 2,3]
  \definition*{s.}{Sobrenome: Xin}
  \definition{adj.}{verdade}
  \definition[封,个,张]{s.}{carta; correio | mensagem; notícia; informação | sinal; evidência | confiança; fé; crédito | detonador (de bombas, etc.) | arsênico}
  \definition{v.}{acreditar; fazer um balanço; dar crédito | deixar à vontade; deixar à mercê; deixar ao acaso | professar fé em; acreditar em}
  \seealsoref{信件}{xin4jian4}
  \synonymref{函}{han2}
  \synonymref{据}{ju4}
  \synonymref{任}{ren4}
  \synonymref{随}{sui2}
  \synonymref{息}{xi1}
  \synonymref{信件}{xin4jian4}
  \synonymref{信仰}{xin4yang3}
  \synonymref{证}{zheng4}
  \antonymref{束}{shu4}
  \antonymref{限}{xian4}
  \antonymref{疑}{yi2}
\end{EntryWithPhonetic}

\begin{EntryWithPhonetic}{信贷}{xin4dai4}{9,9}{⼈,⾙}[HSK 7-9]
  \definition{s.}{crédito; termo geral para atividades de crédito, como depósitos bancários e empréstimos, geralmente se refere a empréstimos bancários}
\end{EntryWithPhonetic}

\begin{EntryWithPhonetic}{信访}{xin4fang3}{9,6}{⼈,⾔}
  \definition{s.}{cartas e visitas (ou seja, cartas de reclamação das pessoas e as visitas que elas fazem para apresentar reclamações) | carta de reclamação | carta de petição}
  \seealsoref{人民来信来访}{ren2min2 lai2xin4 lai2fang3}
  \synonymref{上访}{shang4fang3}
  \synonymref{投诉}{tou2su4}
\end{EntryWithPhonetic}

\begin{EntryWithPhonetic}{信封}{xin4feng1}{9,9}{⼈,⼨}[HSK 3]
  \definition[个,封]{s.}{envelope para cartas}
\end{EntryWithPhonetic}

\begin{EntryWithPhonetic}{信号}{xin4hao4}{9,5}{⼈,⼝}[HSK 2]
  \definition[个,道]{s.}{sinal; luz, ondas de rádio, som, movimento, etc. usados para transmitir mensagens ou comandos | ponte de sinalização; marcação para chamar a atenção, ajudar na identificação e na memória}
  \synonymref{灯号}{deng1hao4}
  \synonymref{记号}{ji4hao5}
\end{EntryWithPhonetic}

\begin{EntryWithPhonetic}{信号灯}{xin4hao4deng1}{9,5,6}{⼈,⼝,⽕}
  \definition[盏]{s.}{lâmpada de sinalização; semáforo | luz de sinalização}
\end{EntryWithPhonetic}

\begin{EntryWithPhonetic}{信件}{xin4jian4}{9,6}{⼈,⼈}[HSK 7-9]
  \definition[封,摞,批,沓]{s.}{cartas, documentos, material impresso, etc. (enviados por correio ou por mensageiro) | carta; correspondência; missiva; cartas e documentos entregues, material impresso}
  \seealsoref{信}{xin4}
\end{EntryWithPhonetic}

\begin{EntryWithPhonetic}{信经}{xin4jing1}{9,8}{⼈,⽷}
  \definition[个]{s.}{crença | credo (seção da missa católica)}
\end{EntryWithPhonetic}

\begin{EntryWithPhonetic}{信赖}{xin4lai4}{9,13}{⼈,⾙}[HSK 7-9]
  \definition{v.}{confiar; contar com; ter fé em}
  \synonymref{相信}{xiang1xin4}
  \synonymref{信任}{xin4ren4}
  \antonymref{猜疑}{cai1yi2}
  \antonymref{怀疑}{huai2yi2}
\end{EntryWithPhonetic}

\begin{EntryWithPhonetic}{信念}{xin4nian4}{9,8}{⼈,⼼}[HSK 5]
  \definition[个,种]{s.}{fé; crença; convicção; concepções consideradas corretas e acreditadas com convicção}
\end{EntryWithPhonetic}

\begin{EntryWithPhonetic}{信任}{xin4ren4}{9,6}{⼈,⼈}[HSK 3]
  \definition{s.}{confiança; um estado mental positivo e conexão emocional}
  \definition{v.}{confiar; ter confiança em; acreditar e ousar confiar}
\end{EntryWithPhonetic}

\begin{EntryWithPhonetic}{信息}{xin4xi1}{9,10}{⼈,⼼}[HSK 2]
  \definition[个,条,段,些]{s.}{notícias; informações; as últimas notícias sobre alguém ou alguma coisa | mensagem; informação; na teoria da informação, uma mensagem transmitida usando símbolos, cujo conteúdo é desconhecido pelo receptor}
  \seealsoref{消息}{xiao1xi5}
  \synonymref{消息}{xiao1xi5}
  \synonymref{新闻}{xin1wen2}
  \synonymref{资讯}{zi1xun4}
\end{EntryWithPhonetic}

\begin{EntryWithPhonetic}{信箱}{xin4xiang1}{9,15}{⼈,⾋}[HSK 5]
  \definition{s.}{caixa de correio; caixa postal instalada pelos correios para que as pessoas possam depositar cartas | caixa postal; caixas com números, localizadas nos correios, que podem ser alugadas para receber correspondência; chamadas de caixas postais exclusivas}
\end{EntryWithPhonetic}

\begin{EntryWithPhonetic}{信心}{xin4xin1}{9,4}{⼈,⼼}[HSK 2]
  \definition[个]{s.}{confiança; fé (em alguém ou algo); a crença de que os desejos se tornarão realidade}
  \synonymref{底气}{di3qi4}
  \synonymref{决心}{jue2xin1}
  \synonymref{信念}{xin4nian4}
  \synonymref{自信}{zi4xin4}
  \antonymref{胆怯}{dan3qie4}
\end{EntryWithPhonetic}

\begin{EntryWithPhonetic}{信仰}{xin4yang3}{9,6}{⼈,⼈}[HSK 6]
  \definition[种]{s.}{crença; religião; refere"-se à ideia de acreditar, adorar e tomar algo como padrão e guia para palavras e ações}
  \definition{v.}{acreditar; crer em; acreditar e adorar uma determinada religião ou doutrina e tomá"-la como guia para palavras e ações}
\end{EntryWithPhonetic}

\begin{EntryWithPhonetic}{信用}{xin4yong4}{9,5}{⼈,⽤}[HSK 6]
  \definition{adj.}{crédito; não é necessária nenhuma garantia material e o dinheiro pode ser reembolsado no prazo}
  \definition[些]{s.}{crédito; confiabilidade; a confiança que você ganha ao fazer o que prometeu | crédito; uma relação de empréstimo"-pagamento ou situação em que o empréstimo é condicionado ao pagamento; uma situação em que um banco empresta dinheiro temporariamente a um cliente e este posteriormente devolve o dinheiro ao banco}
\end{EntryWithPhonetic}

\begin{EntryWithPhonetic}{信用卡}{xin4yong4ka3}{9,5,5}{⼈,⽤,⼘}[HSK 2]
  \definition[张]{s.}{cartão de crédito; moeda eletrônica emitida por um banco ou outra instituição especializada para consumidores; os titulares do cartão podem usá"-lo para sacar dinheiro ou fazer compras de acordo com os regulamentos}
\end{EntryWithPhonetic}

\begin{EntryWithPhonetic}{信誉}{xin4yu4}{9,13}{⼈,⾔}[HSK 7-9]
  \definition{s.}{crédito; prestígio; reputação}
  \seealsoref{信用}{xin4yong4}
  \synonymref{光荣}{guang1rong2}
  \synonymref{名誉}{ming2yu4}
  \synonymref{诺言}{nuo4yan2}
  \synonymref{荣誉}{rong2yu4}
  \synonymref{声誉}{sheng1yu4}
  \synonymref{信用}{xin4yong4}
  \antonymref{虚假}{xu1jia3}
\end{EntryWithPhonetic}

%%%%%%%%%% 兴 %%%%%%%%%%
\subsection*{兴}\addcontentsline{loh}{figure}{兴 \dpy{xing1}}

\begin{EntryWithPhonetic}{兴}{xing1}{6}{⼋}
  \definition*{s.}{Sobrenome: Xing}
  \definition{adj.}{próspero; florescente}
  \definition{adv.}{Dialeto: talvez}
  \definition{v.}{ascender; prosperar; prevalecer; tornar"-se popular | promover; encorajar; fazer prevalecer | começar; iniciar; lançar; mobilizar | erguer"-se; levantar"-se | (usualmente no negativo) permitir; deixar}
  \seeref{xing4}
  \antonymref{衰}{cui1}
  \antonymref{废}{fei4}
  \antonymref{亡}{wang2}
\end{EntryWithPhonetic}

\begin{EntryWithPhonetic}{兴办}{xing1ban4}{6,4}{⼋,⼒}
  \definition{v.}{iniciar; configurar; estabelecer (um negócio)}
  \seealsoref{兴建}{xing1jian4}
  \synonymref{成立}{cheng2li4}
  \synonymref{创办}{chuang4ban4}
  \synonymref{创立}{chuang4li4}
  \synonymref{创造}{chuang4zao4}
  \synonymref{建立}{jian4li4}
  \synonymref{建设}{jian4she4}
  \synonymref{建造}{jian4zao4}
  \synonymref{开发}{kai1fa5}
  \synonymref{设备}{she4bei4}
  \synonymref{作战}{zuo4zhan4}
\end{EntryWithPhonetic}

\begin{EntryWithPhonetic}{兴奋}{xing1fen4}{6,8}{⼋,⼤}[HSK 4]
  \definition{adj.}{animado; excitante; empolgante}
  \definition{s.}{excitação; empolgação}
  \definition{v.}{excitar; intoxicar}
\end{EntryWithPhonetic}

\begin{EntryWithPhonetic}{兴奋剂}{xing1fen4ji4}{6,8,8}{⼋,⼤,⼑}[HSK 7-9]
  \definition{s.}{excitante; estimulante; droga | analéptico; tônico; agonista | uma injeção de ânimo}
\end{EntryWithPhonetic}

\begin{EntryWithPhonetic}{兴建}{xing1jian4}{6,8}{⼋,⼵}[HSK 7-9]
  \definition{v.}{construir; edificar; iniciar a construção (geralmente referindo"-se a projetos de grande escala)}
  \seealsoref{兴办}{xing1ban4}
  \antonymref{破坏}{po4huai4}
\end{EntryWithPhonetic}

\begin{EntryWithPhonetic}{兴起}{xing1qi3}{6,10}{⼋,⾛}[HSK 7-9]
  \definition{v.}{surgir; brotar; crescer; estar em ascensão; começar a aparecer e a florescer | ficar excitado; aumentar a excitação}
  \synonymref{崛起}{jue2qi3}
  \antonymref{灭亡}{mie4wang2}
\end{EntryWithPhonetic}

\begin{EntryWithPhonetic}{兴旺}{xing1wang4}{6,8}{⼋,⽇}[HSK 6]
  \definition{adj.}{próspero; propício; favorável; auspicioso}
\end{EntryWithPhonetic}

%%%%%%%%%% 星 %%%%%%%%%%
\subsection*{星}\addcontentsline{loh}{figure}{星 \dpy{xing1}}

\begin{EntryWithPhonetic}{星}{xing1}{9}{⽇}
  \definition*{s.}{Xing, a vigésima quinta das vinte e oito constelações em que a esfera celeste era dividida na antiga astronomia chinesa, consistindo em sete estrelas em Hydra}
  \definition[颗]{s.}{estrela | (astronomia) corpo celeste | partícula | pequenas marcas no braço de uma balança romana indicando jin e suas frações | artista famoso (estrela de cinema, estrela de jogos de bola, etc.) | satélite (artificial) | pequena quantidade}
\end{EntryWithPhonetic}

\begin{EntryWithPhonetic}{星表}{xing1biao3}{9,8}{⽇,⾐}
  \definition{s.}{catálogo de estrelas}
\end{EntryWithPhonetic}

\begin{EntryWithPhonetic}{星辰}{xing1chen2}{9,7}{⽇,⾠}
  \definition{s.}{as estrelas; estrelas e constelações; um termo geral para estrelas}
\end{EntryWithPhonetic}

\begin{EntryWithPhonetic}{星火}{xing1huo3}{9,4}{⽇,⽕}
  \definition{s.}{faísca | estrela cadente; meteoro; a luz de um meteoro | foco de incêndio; um pequeno incêndio}
\end{EntryWithPhonetic}

\begin{EntryWithPhonetic}{星际迷航}{xing1ji4 mi2hang2}{9,7,9,10}{⽇,⾩,⾡,⾈}
  \definition*{s.}{Jornada nas Estrelas; \emph{Star Trek}}
\end{EntryWithPhonetic}

\begin{EntryWithPhonetic}{星期}{xing1qi1}{9,12}{⽇,⽉}[HSK 1]
  \definition[个]{s.}{semana | dias da semana; usado em conjunto com 日, 一, 二, 三, 四, 五, 六, 天, indica um determinado dia da semana | abreviação de domingo}
  \seealsoref{星期二}{xing1qi1er4}
  \seealsoref{星期六}{xing1qi1liu4}
  \seealsoref{星期日}{xing1qi1ri4}
  \seealsoref{星期三}{xing1qi1san1}
  \seealsoref{星期四}{xing1qi1si4}
  \seealsoref{星期天}{xing1qi1tian1}
  \seealsoref{星期五}{xing1qi1wu3}
  \seealsoref{星期一}{xing1qi1yi1}
  \synonymref{礼拜}{li3bai4}
  \synonymref{周期}{zhou1qi1}
  \antonymref{工作日}{gong1zuo4ri4}
\end{EntryWithPhonetic}

\begin{EntryWithPhonetic}{星期二}{xing1qi1er4}{9,12,2}{⽇,⽉,⼆}
  \definition{s.}{terça"-feira}
\end{EntryWithPhonetic}

\begin{EntryWithPhonetic}{星期六}{xing1qi1liu4}{9,12,4}{⽇,⽉,⼋}
  \definition{s.}{sábado}
\end{EntryWithPhonetic}

\begin{EntryWithPhonetic}{星期日}{xing1qi1ri4}{9,12,4}{⽇,⽉,⽇}[HSK 1]
  \definition[个]{s.}{domingo}
  \seealsoref{星期天}{xing1qi1tian1}
  \synonymref{星期天}{xing1qi1tian1}
\end{EntryWithPhonetic}

\begin{EntryWithPhonetic}{星期三}{xing1qi1san1}{9,12,3}{⽇,⽉,⼀}
  \definition{s.}{quarta"-feira}
\end{EntryWithPhonetic}

\begin{EntryWithPhonetic}{星期四}{xing1qi1si4}{9,12,5}{⽇,⽉,⼞}
  \definition{s.}{quinta"-feira}
\end{EntryWithPhonetic}

\begin{EntryWithPhonetic}{星期天}{xing1qi1tian1}{9,12,4}{⽇,⽉,⼤}[HSK 1]
  \definition[个]{s.}{domingo}
  \seealsoref{星期日}{xing1qi1ri4}
  \synonymref{星期日}{xing1qi1ri4}
  \synonymref{周末}{zhou1mo4}
\end{EntryWithPhonetic}

\begin{EntryWithPhonetic}{星期五}{xing1qi1wu3}{9,12,4}{⽇,⽉,⼆}
  \definition{s.}{sexta"-feira}
\end{EntryWithPhonetic}

\begin{EntryWithPhonetic}{星期一}{xing1qi1yi1}{9,12,1}{⽇,⽉,⼀}
  \definition{s.}{segunda"-feira}
\end{EntryWithPhonetic}

\begin{EntryWithPhonetic}{星球大战}{xing1qiu2 da4zhan4}{9,11,3,9}{⽇,⽟,⼤,⼽}
  \definition*{s.}{Guerra nas Estrelas; \emph{Star Wars} | outro nome para a Iniciativa de Defesa Estratégica (SDI), lançada pelo presidente dos EUA Ronald Reagan}
\end{EntryWithPhonetic}

\begin{EntryWithPhonetic}{星星}{xing1xing5}{9,9}{⽇,⽇}[HSK 2]
  \definition[颗,群,片]{s.}{estrela; em astronomia, refere"-se aos corpos celestes luminosos no universo, como as estrelas que brilham no céu noturno | estrela; uma metáfora para alguém ou algo que se destaca em um determinado campo e atrai atenção | objetos em forma de estrela}
  \synonymref{星辰}{xing1chen2}
\end{EntryWithPhonetic}

\begin{EntryWithPhonetic}{星座}{xing1zuo4}{9,10}{⽇,⼴}[HSK 7-9]
  \definition[张]{s.}{Astronomia: constelação, cada uma das várias divisões (regiões) do céu noturno}
\end{EntryWithPhonetic}

%%%%%%%%%% 猩 %%%%%%%%%%
\subsection*{猩}\addcontentsline{loh}{figure}{猩 \dpy{xing1}}

\begin{EntryWithPhonetic}{猩}{xing1}{12}{⽝}
  \definition[只]{s.}{orangotango}
\end{EntryWithPhonetic}

\begin{EntryWithPhonetic}{猩猩}{xing1xing5}{12,12}{⽝,⽝}[HSK 7-9]
  \definition{s.}{orangotango;  \emph{Pongo satyrus}; \emph{Pongo pygnaeus}}
\end{EntryWithPhonetic}

%%%%%%%%%% 腥 %%%%%%%%%%
\subsection*{腥}\addcontentsline{loh}{figure}{腥 \dpy{xing1}}

\begin{EntryWithPhonetic}{腥}{xing1}{13}{⾁}[HSK 7-9]
  \definition{adj.}{ter cheiro de peixe, frutos do mar, etc.}
  \definition{s.}{carne ou peixe crus | cheiro de peixe}
\end{EntryWithPhonetic}

%%%%%%%%%% 刑 %%%%%%%%%%
\subsection*{刑}\addcontentsline{loh}{figure}{刑 \dpy{xing2}}

\begin{EntryWithPhonetic}{刑}{xing2}{6}{⼑}
  \definition*{s.}{Sobrenome: Xing}
  \definition{s.}{punição (aplicada por um crime); pena; sentença | tortura; castigo corporal | castigo}
\end{EntryWithPhonetic}

\begin{EntryWithPhonetic}{刑法}{xing2fa3}{6,8}{⼑,⽔}[HSK 7-9]
  \definition{s.}{castigo corporal; tortura | código penal; direito penal; \emph{jus criminale}}
\end{EntryWithPhonetic}

%%%%%%%%%% 行 %%%%%%%%%%
\subsection*{行}\addcontentsline{loh}{figure}{行 \dpy{xing2}}

\begin{EntryWithPhonetic}{行}{xing2}{6}{⾏}[HSK 1][Kangxi 144]
  \definition*{s.}{Sobrenome: Xing}
  \definition{adj.}{de viajar; relacionado a viagens | temporário; improvisado; provisório | capaz; competente}
  \definition{adv.}{em breve}
  \definition{s.}{comportamento; conduta | caligrafia cursiva (na caligrafia chinesa); escrita cursiva}
  \definition{v.}{ir | fazer uma viagem | estar em voga; prevalecer; circular | fazer; executar; realizar; envolver"-se em | estar tudo bem; O.K. | indica a realização de uma determinada atividade (usado principalmente antes de verbos dissilábicos) | (em medicina) fazer efeito}
  \seeref{hang2}
  \seeref{hang4}
  \seeref{heng2}
  \synonymref{动}{dong4}
  \synonymref{走}{zou3}
  \antonymref{止}{zhi3}
\end{EntryWithPhonetic}

\begin{EntryWithPhonetic}{行程}{xing2cheng2}{6,12}{⾏,⽲}[HSK 6]
  \definition{s.}{rota ou distância de viagem; distância; jornada | curso; progresso; processo | curso; deslocamento; viagem}[活塞行程有点不对劲。===Há algo errado com o curso do pistão.]
\end{EntryWithPhonetic}

\begin{EntryWithPhonetic}{行动}{xing2dong4}{6,6}{⾏,⼒}[HSK 2]
  \definition[次,场,项]{s.}{ação; operação; comportamento}
  \definition{v.}{circular; mover"-se; andar | agir; tomar medidas; atividades para atingir um determinado propósito}
  \seealsoref{行为}{xing2wei2}
  \synonymref{动作}{dong4zuo4}
  \synonymref{活动}{huo2dong5}
  \synonymref{举措}{ju3cuo4}
  \synonymref{举止}{ju3zhi3}
  \synonymref{行为}{xing2wei2}
  \synonymref{运动}{yun4dong5}
  \synonymref{作为}{zuo4wei2}
  \antonymref{等待}{deng3dai4}
  \antonymref{口头}{kou3tou2}
  \antonymref{思考}{si1kao3}
  \antonymref{停歇}{ting2xie1}
\end{EntryWithPhonetic}

\begin{EntryWithPhonetic}{行进}{xing2jin4}{6,7}{⾏,⾡}
  \definition{s.}{(geralmente de tropas) marchar para a frente; avançar | progredir; marchar para a frente; caminhar para a frente (usado principalmente em grupos)}
  \synonymref{迈进}{mai4jin4}
  \synonymref{前进}{qian2jin4}
  \synonymref{行走}{xing2zou3}
  \antonymref{后退}{hou4tui4}
  \antonymref{静止}{jing4zhi3}
  \antonymref{退却}{tui4que4}
  \antonymref{退缩}{tui4suo1}
\end{EntryWithPhonetic}

\begin{EntryWithPhonetic}{行礼}{xing2li3}{6,5}{⾏,⽰}
  \definition{v.}{saudar | Dialeto: dar presentes | prestar continência}
  \synonymref{敬礼}{jing4/li3}
\end{EntryWithPhonetic}

\begin{EntryWithPhonetic}{行李}{xing2li5}{6,7}{⾏,⽊}[HSK 3]
  \definition[点,个]{s.}{bagagem, malas, cestas de vime, etc. que você leva quando sai de casa}
\end{EntryWithPhonetic}

\begin{EntryWithPhonetic}{行人}{xing2ren2}{6,2}{⾏,⼈}[HSK 2]
  \definition[个]{s.}{pedestre; transeunte; viajante à pé; pessoas caminhando na estrada}
  \synonymref{路人}{lu4ren2}
\end{EntryWithPhonetic}

\begin{EntryWithPhonetic}{行使}{xing2shi3}{6,8}{⾏,⼈}[HSK 7-9]
  \definition{v.}{exercitar; realizar; executar (deveres); usar (poderes)}
  \synonymref{利用}{li4yong4}
  \synonymref{使用}{shi3yong4}
  \synonymref{应用}{ying4yong4}
  \synonymref{运用}{yun4yong4}
  \antonymref{放弃}{fang4qi4}
\end{EntryWithPhonetic}

\begin{EntryWithPhonetic}{行驶}{xing2shi3}{6,8}{⾏,⾺}[HSK 5]
  \definition{v.}{ir; navegar; viajar (utilizando um veículo, navio, etc.)}
\end{EntryWithPhonetic}

\begin{EntryWithPhonetic}{行为}{xing2wei2}{6,4}{⾏,⼂}[HSK 2]
  \definition[个,种,类]{s.}{ação; comportamento; conduta; atividades que são controladas por pensamentos e manifestadas externamente}
  \seealsoref{行动}{xing2dong4}
  \synonymref{动作}{dong4zuo4}
  \synonymref{活动}{huo2dong5}
  \synonymref{举动}{ju3dong4}
  \synonymref{举止}{ju3zhi3}
  \synonymref{手脚}{shou3jiao3}
  \synonymref{行动}{xing2dong4}
  \synonymref{作为}{zuo4wei2}
  \antonymref{口头}{kou3tou2}
  \antonymref{内心}{nei4xin1}
  \antonymref{思想}{si1xiang3}
\end{EntryWithPhonetic}

\begin{EntryWithPhonetic}{行星}{xing2xing1}{6,9}{⾏,⽇}
  \definition[颗]{s.}{planeta}
  \seealsoref{惑星}{huo4xing1}
\end{EntryWithPhonetic}

\begin{EntryWithPhonetic}{行凶}{xing2/xiong1}{6,4}{⾏,⼐}
  \definition{v.+compl.}{cometer agressão física ou homicídio; praticar violência}
  \synonymref{打架}{da3/jia4}
\end{EntryWithPhonetic}

\begin{EntryWithPhonetic}{行政}{xing2zheng4}{6,9}{⾏,⽁}[HSK 7-9]
  \definition{s.}{administração; refere"-se ao trabalho de gestão interna de agências governamentais, empresas, organizações, etc.}
  \definition{v.}{administrar; exercício do poder estatal}
\end{EntryWithPhonetic}

\begin{EntryWithPhonetic}{行走}{xing2zou3}{6,7}{⾏,⾛}[HSK 7-9]
  \definition{v.}{andar; caminhar}
  \synonymref{徒步}{tu2bu4}
  \antonymref{住宿}{zhu4su4}
\end{EntryWithPhonetic}

%%%%%%%%%% 邢 %%%%%%%%%%
\subsection*{邢}\addcontentsline{loh}{figure}{邢 \dpy{xing2}}

\begin{EntryWithPhonetic}{邢}{xing2}{6}{⾢}
  \definition*{s.}{Usado em nomes de lugares, por exemplo XingTai (邢台) | Sobrenome: Xing}
  \seealsoref{邢台}{xing2tai2}
\end{EntryWithPhonetic}

\begin{EntryWithPhonetic}{邢台}{xing2tai2}{6,5}{⾢,⼝}
  \definition*{s.}{Xingtai, uma cidade na província de Hebei (河北)}
  \seealsoref{河北}{he2bei3}
\end{EntryWithPhonetic}

%%%%%%%%%% 形 %%%%%%%%%%
\subsection*{形}\addcontentsline{loh}{figure}{形 \dpy{xing2}}

\begin{EntryWithPhonetic}{形}{xing2}{7}{⼺}[HSK 6]
  \definition{s.}{forma; formato | corpo; entidade}
  \definition{v.}{aparecer; revelar; mostrar | comparar; contrastar}
\end{EntryWithPhonetic}

\begin{EntryWithPhonetic}{形成}{xing2cheng2}{7,6}{⼺,⼽}[HSK 3]
  \definition{v.}{moldar; formar; tomar forma; tornar"-se algo ou surgir uma situação após mudanças e desenvolvimentos}
\end{EntryWithPhonetic}

\begin{EntryWithPhonetic}{形而上学}{xing2'er2shang4xue2}{7,6,3,8}{⼺,⽽,⼀,⼦}
  \definition{s.}{metafísica}
\end{EntryWithPhonetic}

\begin{EntryWithPhonetic}{形容}{xing2rong2}{7,10}{⼺,⼧}[HSK 4]
  \definition{s.}{aparência; semblante}
  \definition{v.}{descrever}
\end{EntryWithPhonetic}

\begin{EntryWithPhonetic}{形式}{xing2shi4}{7,6}{⼺,⼷}[HSK 3]
  \definition[种,个]{s.}{forma; formato; modalidade; a aparência, estrutura ou estado das coisas, etc.}
\end{EntryWithPhonetic}

\begin{EntryWithPhonetic}{形势}{xing2shi4}{7,8}{⼺,⼒}[HSK 4]
  \definition[个,种]{s.}{terreno; características topográficas; situação geográfica, principalmente de uma perspectiva militar | situação; circunstâncias; a situação geral, a tendência de como as coisas estão se desenvolvendo e mudando | geralmente não é usado em situações pessoais}
\end{EntryWithPhonetic}

\begin{EntryWithPhonetic}{形态}{xing2tai4}{7,8}{⼺,⼼}[HSK 5]
  \definition[种]{s.}{forma; forma como as coisas se apresentam | forma; padrão; postura | morfologia; forma; Gramática: refere"-se às formas internas de mudança das palavras, incluindo a formação de palavras e as mudanças morfológicas}
\end{EntryWithPhonetic}

\begin{EntryWithPhonetic}{形象}{xing2xiang4}{7,11}{⼺,⾗}[HSK 3]
  \definition{adj.}{vívido; expressão concreta e vívida}
  \definition[个,种]{s.}{imagem; forma; figura; formas ou posturas específicas que podem despertar pensamentos ou emoções nas pessoas | imagem literária; imagem artística; pessoas ou coisas com características diferentes criadas na literatura, no cinema e em outras artes}
\end{EntryWithPhonetic}

\begin{EntryWithPhonetic}{形形色色}{xing2xing2se4se4}{7,7,6,6}{⼺,⼺,⾊,⾊}[HSK 7-9]
  \definition{expr.}{de todas as cores; de todos os tipos de homens ou coisas; de todas as formas; de todas as tonalidades; de todas as descrições; de grande variedade e diversidade; vários}
  \synonymref{各式各样}{ge4shi4-ge4yang4}
  \synonymref{五花八门}{wu3hua1-ba1men2}
  \synonymref{应有尽有}{ying1you3-jin4you3}
\end{EntryWithPhonetic}

\begin{EntryWithPhonetic}{形影不离}{xing2ying3-bu4li2}{7,15,4,10}{⼺,⼺,⼀,⼇}[HSK 7-9]
  \definition{expr.}{inseparáveis; unidos pelo quadril; descrevendo um relacionamento próximo}
\end{EntryWithPhonetic}

\begin{EntryWithPhonetic}{形状}{xing2zhuang4}{7,7}{⼺,⽝}[HSK 3]
  \definition[个,种]{s.}{forma; aparência ; aspecto; a aparência de um objeto ou figura, representada pela combinação de superfícies ou linhas externas}
\end{EntryWithPhonetic}

%%%%%%%%%% 型 %%%%%%%%%%
\subsection*{型}\addcontentsline{loh}{figure}{型 \dpy{xing2}}

\begin{EntryWithPhonetic}{型}{xing2}{9}{⼟}[HSK 4]
  \definition{s.}{molde; modelo | modelo; tipo; padrão}
\end{EntryWithPhonetic}

\begin{EntryWithPhonetic}{型号}{xing2hao4}{9,5}{⼟,⼝}[HSK 4]
  \definition[个,种]{s.}{modelo; tipo; refere"-se ao desempenho, às especificações e ao tamanho de aeronaves, máquinas, implementos agrícolas, etc.}
\end{EntryWithPhonetic}

%%%%%%%%%% 省 %%%%%%%%%%
\subsection*{省}\addcontentsline{loh}{figure}{省 \dpy{xing3}}

\begin{EntryWithPhonetic}{省}{xing3}{9}{⽬}
  \definition{v.}{examinar"-se criticamente; verificar (os próprios pensamentos, palavras e ações) | visitar (especialmente os pais ou pessoas mais velhas) | estar ciente; tornar"-se consciente; compreender; tomar consciência | examinar minuciosamente; inspecionar; escrutinar}
  \seeref{sheng3}
\end{EntryWithPhonetic}

\begin{EntryWithPhonetic}{省悟}{xing3wu4}{9,10}{⽬,⼼}
  \definition{v.}{chegar a perceber (ou ver) a verdade, o próprio erro, etc.; despertar para a realidade | cair na real | perceber | ver a verdade}
  \synonymref{觉悟}{jue2wu4}
  \synonymref{觉醒}{jue2xing3}
  \synonymref{领悟}{ling3wu4}
  \synonymref{醒悟}{xing3wu4}
  \antonymref{沉迷}{chen2mi2}
  \antonymref{糊涂}{hu2tu5}
  \antonymref{迷惑}{mi2huo4}
\end{EntryWithPhonetic}

%%%%%%%%%% 醒 %%%%%%%%%%
\subsection*{醒}\addcontentsline{loh}{figure}{醒 \dpy{xing3}}

\begin{EntryWithPhonetic}{醒}{xing3}{16}{⾣}[HSK 4]
  \definition{adj.}{impressionante; notável; admirável; atraente; chamativo}
  \definition{v.}{ficar sóbrio; voltar a si; recuperar a consciência; retornar à normalidade após intoxicação, anestesia ou coma | despertar; estar acordado | ter a mente clara; mover a consciência da confusão para a compreensão | vir a entender; tornar"-se ciente de; tomar consciência de}
\end{EntryWithPhonetic}

\begin{EntryWithPhonetic}{醒来}{xing3lai5}{16,7}{⾣,⽊}[HSK 7-9]
  \definition{v.}{acordar; o estado de sono terminou ou ocorreu antes de adormecer}
  \synonymref{觉醒}{jue2xing3}
  \synonymref{起床}{qi3/chuang2}
  \synonymref{起来}{qi3/lai2}
  \antonymref{睡觉}{shui4/jiao4}
\end{EntryWithPhonetic}

\begin{EntryWithPhonetic}{醒目}{xing3mu4}{16,5}{⾣,⽬}[HSK 7-9]
  \definition{adj.}{chamativo; que chama a atenção; que atrai o olhar (de palavras escritas ou imagens); (texto, imagens, etc.) as imagens são nítidas e fáceis de visualizar}
  \synonymref{精明}{jing1ming2}
  \synonymref{精通}{jing1tong1}
  \synonymref{能干}{neng2gan4}
  \synonymref{显眼}{xian3yan3}
  \synonymref{耀眼}{yao4yan3}
  \antonymref{懵懂}{meng3dong3}
  \antonymref{模糊}{mo2hu5}
\end{EntryWithPhonetic}

\begin{EntryWithPhonetic}{醒悟}{xing3wu4}{16,10}{⾣,⼼}[HSK 7-9]
  \definition{v.}{chegar a perceber a verdade (o próprio erro, etc.); despertar para a realidade; despertar e compreender, libertando"-se da confusão e do deslumbramento}
  \seealsoref{觉悟}{jue2wu4}
  \synonymref{觉悟}{jue2wu4}
  \synonymref{觉醒}{jue2xing3}
\end{EntryWithPhonetic}

%%%%%%%%%% 兴 %%%%%%%%%%
\subsection*{兴}\addcontentsline{loh}{figure}{兴 \dpy{xing4}}

\begin{EntryWithPhonetic}{兴}{xing4}{6}{⼋}
  \definition{s.}{estado de espírito ou desejo de fazer algo; interesse; entusiasmo}
  \seeref{xing1}
\end{EntryWithPhonetic}

\begin{EntryWithPhonetic}{兴高采烈}{xing4gao1-cai3lie4}{6,10,8,10}{⼋,⾼,⾤,⽕}[HSK 7-9]
  \definition{expr.}{feliz e animado (expressão idiomática); originalmente, referia"-se a artigos com ideais elevados e linguagem incisiva; mais tarde, passou a ser usado para descrever alguém que está de bom humor e cheio de energia; de bom humor; animado; jubilante; eufórico}
  \antonymref{愁眉苦脸}{chou2mei2-ku3lian3}
  \antonymref{无精打采}{wu2jing1-da3cai3}
\end{EntryWithPhonetic}

\begin{EntryWithPhonetic}{兴趣}{xing4qu4}{6,15}{⼋,⾛}[HSK 4]
  \definition[个,种,点,股,份]{s.}{interesse (desejo de conhecer sobre alguma coisa ou coisa no qual está interessado) | \emph{hobby}}
\end{EntryWithPhonetic}

\begin{EntryWithPhonetic}{兴致}{xing4zhi4}{6,10}{⼋,⾄}[HSK 7-9]
  \definition{s.}{interesse; disposição para desfrutar}
  \seealsoref{兴趣}{xing4qu4}
  \synonymref{趣味}{qu4wei4}
  \synonymref{玩味}{wan2wei4}
  \synonymref{兴趣}{xing4qu4}
\end{EntryWithPhonetic}

%%%%%%%%%% 姓 %%%%%%%%%%
\subsection*{姓}\addcontentsline{loh}{figure}{姓 \dpy{xing4}}

\begin{EntryWithPhonetic}{姓}{xing4}{8}{⼥}[HSK 2]
  \definition[个]{s.}{sobrenome; nome de família; um caractere que representa um sistema familiar, os chineses colocam o sobrenome em primeiro lugar e o nome em segundo}
  \definition{v.}{ter como sobrenome; tratar um ou mais caracteres como sobrenome}
  \seealsoref{贵姓}{gui4xing4}
  \seealsoref{名字}{ming2zi5}
  \synonymref{贵姓}{gui4xing4}
  \synonymref{名字}{ming2zi5}
\end{EntryWithPhonetic}

\begin{EntryWithPhonetic}{姓名}{xing4ming2}{8,6}{⼥,⼝}[HSK 2]
  \definition{s.}{nome; nome completo; sobrenome e nome próprio}
  \synonymref{名字}{ming2zi5}
\end{EntryWithPhonetic}

\begin{EntryWithPhonetic}{姓氏}{xing4shi4}{8,4}{⼥,⽒}[HSK 7-9]
  \definition{s.}{sobrenome; nome de família; sobrenome e nome de clã; o sobrenome tinha origem na linhagem materna, enquanto o nome de clã na linhagem paterna; mais tarde, os dois tornaram-"se indistinguíveis e passaram a ser referidos coletivamente como sobrenome ou nome de clã}
  \synonymref{姓名}{xing4ming2}
\end{EntryWithPhonetic}

%%%%%%%%%% 幸 %%%%%%%%%%
\subsection*{幸}\addcontentsline{loh}{figure}{幸 \dpy{xing4}}

\begin{EntryWithPhonetic}{幸}{xing4}{8}{⼲}
  \definition*{s.}{Sobrenome: Xing}
  \definition{adj.}{feliz}
  \definition{adv.}{afortunadamente; felizmente}
  \definition{s.}{felicidade}
  \definition{v.}{alegrar"-se; sentir"-se feliz e contente | favorecer; patrocinar | vir; chegar; antigamente, referia"-se à chegada de um monarca a um determinado lugar}
\end{EntryWithPhonetic}

\begin{EntryWithPhonetic}{幸存}{xing4cun2}{8,6}{⼲,⼦}[HSK 7-9]
  \definition{v.}{sobreviver}
  \synonymref{幸免}{xing4mian3}
  \antonymref{遇难}{yu4/nan4}
\end{EntryWithPhonetic}

\begin{EntryWithPhonetic}{幸福}{xing4fu2}{8,13}{⼲,⽰}[HSK 3]
  \definition{adj.}{feliz; a vida, a família e outras circunstâncias deixam as pessoas satisfeitas e felizes}
  \definition{s.}{felicidade; bem estar; sensação ou experiência satisfatória e feliz, etc.}
\end{EntryWithPhonetic}

\begin{EntryWithPhonetic}{幸好}{xing4hao3}{8,6}{⼲,⼥}[HSK 7-9]
  \definition{adv.}{felizmente; por sorte; afortunadamente}
  \synonymref{多亏}{duo1kui1}
  \synonymref{好在}{hao3zai4}
  \synonymref{幸亏}{xing4kui1}
  \antonymref{可惜}{ke3xi1}
\end{EntryWithPhonetic}

\begin{EntryWithPhonetic}{幸亏}{xing4kui1}{8,3}{⼲,⼆}[HSK 7-9]
  \definition{adv.}{felizmente; por sorte; por obra do destino; isso indica que consequências adversas foram evitadas por acaso devido a certas condições}
  \seealsoref{多亏}{duo1kui1}
  \seealsoref{幸好}{xing4hao3}
  \synonymref{多亏}{duo1kui1}
  \synonymref{好在}{hao3zai4}
  \synonymref{幸好}{xing4hao3}
\end{EntryWithPhonetic}

\begin{EntryWithPhonetic}{幸免}{xing4mian3}{8,7}{⼲,⼉}[HSK 7-9]
  \definition{v.}{escapar por pura sorte; ter um escape por pouco; evitar por acaso}
  \synonymref{幸存}{xing4cun2}
\end{EntryWithPhonetic}

\begin{EntryWithPhonetic}{幸运}{xing4yun4}{8,7}{⼲,⾡}[HSK 3]
  \definition{adj.}{sortudo; feliz; afortunado}
  \definition[个,点,丝]{s.}{boa sorte; boa fortuna}
\end{EntryWithPhonetic}

\begin{EntryWithPhonetic}{幸运抽奖}{xing4yun4 chou1jiang3}{8,7,8,9}{⼲,⾡,⼿,⼤}
  \definition{s.}{loteria | sorteio}
\end{EntryWithPhonetic}

\begin{EntryWithPhonetic}{幸运儿}{xing4yun4'er2}{8,7,2}{⼲,⾡,⼉}
  \definition{s.}{o favorito da fortuna; sujeito sortudo | cão sortudo; pessoa sortuda}
\end{EntryWithPhonetic}

%%%%%%%%%% 性 %%%%%%%%%%
\subsection*{性}\addcontentsline{loh}{figure}{性 \dpy{xing4}}

\begin{EntryWithPhonetic}{性}{xing4}{8}{⼼}[HSK 3]
  \definition[个]{s.}{natureza; caráter; personalidade | propriedade; qualidade; natureza e características das coisas | sexo; gênero | sexualidade; relacionado com a reprodução e a sexualidade | caráter; temperamento}
  \definition{suf.}{indica uma determinada propriedade ou característica de algo; segue um substantivo, verbo ou adjetivo, formando um substantivo abstrato ou um adjetivo que expressa uma propriedade}
\end{EntryWithPhonetic}

\begin{EntryWithPhonetic}{性别}{xing4bie2}{8,7}{⼼,⼑}[HSK 3]
  \definition[种]{s.}{sexo; gênero}
\end{EntryWithPhonetic}

\begin{EntryWithPhonetic}{性格}{xing4ge2}{8,10}{⼼,⽊}[HSK 3]
  \definition[种,个]{s.}{caráter; temperamento; as características psicológicas manifestadas na atitude e no comportamento em relação às pessoas e às coisas}
\end{EntryWithPhonetic}

\begin{EntryWithPhonetic}{性价比}{xing4jia4bi3}{8,6,4}{⼼,⼈,⽐}[HSK 7-9]
  \definition{s.}{relação custo"-benefício; a relação entre o valor de desempenho de um produto e seu valor de preço é uma medida quantitativa que reflete a acessibilidade do produto | desempenho de custo}
\end{EntryWithPhonetic}

\begin{EntryWithPhonetic}{性命}{xing4ming4}{8,8}{⼼,⼝}[HSK 7-9]
  \definition[条]{s.}{vida (de um homem ou animal)}
  \seealsoref{生命}{sheng1ming4}
  \synonymref{生命}{sheng1ming4}
\end{EntryWithPhonetic}

\begin{EntryWithPhonetic}{性能}{xing4neng2}{8,10}{⼼,⾁}[HSK 5]
  \definition{s.}{natureza; propriedade; desempenho; função (de uma máquina, etc.); grau de conformidade dos produtos mecânicos ou outros produtos industriais com os requisitos de projeto}
\end{EntryWithPhonetic}

\begin{EntryWithPhonetic}{性侵}{xing4qin1}{8,9}{⼼,⼈}
  \definition{s.}{agressão sexual}
  \definition{v.}{agredir sexualmente}
\end{EntryWithPhonetic}

\begin{EntryWithPhonetic}{性情}{xing4qing2}{8,11}{⼼,⼼}[HSK 7-9]
  \definition[种]{s.}{disposição; temperamento; personalidade; disposição e caráter | sentimento; emoções; natureza interior; caráter emocional; pensamentos e sentimentos}
  \seealsoref{性格}{xing4ge2}
  \synonymref{本性}{ben3xing4}
  \synonymref{个性}{ge4xing4}
  \synonymref{脾气}{pi2qi5}
  \synonymref{特性}{te4xing4}
  \synonymref{天性}{tian1xing4}
  \synonymref{性格}{xing4ge2}
\end{EntryWithPhonetic}

\begin{EntryWithPhonetic}{性生活}{xing4sheng1huo2}{8,5,9}{⼼,⽣,⽔}
  \definition{s.}{vida sexual}
\end{EntryWithPhonetic}

\begin{EntryWithPhonetic}{性质}{xing4zhi4}{8,8}{⼼,⾙}[HSK 4]
  \definition[个,种,类]{s.}{natureza; qualidade; caráter; propriedade; propriedade fundamental que distingue uma coisa de outra}
\end{EntryWithPhonetic}

%%%%%%%%%% 凶 %%%%%%%%%%
\subsection*{凶}\addcontentsline{loh}{figure}{凶 \dpy{xiong1}}

\begin{EntryWithPhonetic}{凶}{xiong1}{4}{⼐}[HSK 6]
  \definition{adj.}{sinistro; desfavorável; azarado | ruim para as colheitas; improdutivo; ameaçado pela fome | feroz; vicioso; cruel | medroso; terrível}
  \definition[个]{s.}{mal; assassinato; ato de violência; atos de matar ou ferir pessoas | assassino; malfeitor; criminoso; pessoa má; pessoa violenta}
  \definition{v.}{ser feroz; tratar cruelmente}
  \antonymref{吉}{ji2}
\end{EntryWithPhonetic}

\begin{EntryWithPhonetic}{凶残}{xiong1can2}{4,9}{⼐,⽍}[HSK 7-9]
  \definition{adj.}{feroz e cruel; selvagem e impiedoso; implacável | cruel}
  \synonymref{残忍}{can2ren3}
  \synonymref{粗暴}{cu1bao4}
  \synonymref{凶恶}{xiong1'e4}
  \synonymref{凶狠}{xiong1hen3}
  \antonymref{慈善}{ci2shan4}
  \antonymref{和蔼}{he2'ai3}
  \antonymref{和气}{he2qi5}
  \antonymref{平和}{ping2he2}
  \antonymref{仁慈}{ren2ci2}
  \antonymref{善良}{shan4liang2}
\end{EntryWithPhonetic}

\begin{EntryWithPhonetic}{凶恶}{xiong1'e4}{4,10}{⼐,⼼}[HSK 7-9]
  \definition{adj.}{feroz; vicioso; impiedoso; diabólico; rufião; malévolo}
  \seealsoref{凶狠}{xiong1hen3}
  \synonymref{残忍}{can2ren3}
  \synonymref{粗暴}{cu1bao4}
  \synonymref{厉害}{li4hai5}
  \antonymref{慈善}{ci2shan4}
  \antonymref{慈祥}{ci2xiang2}
  \antonymref{和蔼}{he2'ai3}
  \antonymref{和气}{he2qi5}
  \antonymref{平和}{ping2he2}
  \antonymref{善良}{shan4liang2}
\end{EntryWithPhonetic}

\begin{EntryWithPhonetic}{凶狠}{xiong1hen3}{4,9}{⼐,⽝}[HSK 7-9]
  \definition{adj.}{feroz e malicioso; malévolo | poderoso; vigoroso | cruel | vingativo | vicioso}
  \seealsoref{凶恶}{xiong1'e4}
  \synonymref{残酷}{can2ku4}
  \synonymref{残忍}{can2ren3}
  \synonymref{粗暴}{cu1bao4}
  \synonymref{邪恶}{xie2'e4}
  \synonymref{凶恶}{xiong1'e4}
  \antonymref{和蔼}{he2'ai3}
  \antonymref{和气}{he2qi5}
  \antonymref{平和}{ping2he2}
  \antonymref{仁慈}{ren2ci2}
  \antonymref{善良}{shan4liang2}
  \antonymref{温柔}{wen1rou2}
\end{EntryWithPhonetic}

\begin{EntryWithPhonetic}{凶猛}{xiong1meng3}{4,11}{⼐,⽝}[HSK 7-9]
  \definition{adj.}{violento; terrível; feroz; descreve o poder destrutivo de animais, vento, chuva ou outras situações}
  \synonymref{粗暴}{cu1bao4}
  \synonymref{厉害}{li4hai5}
  \synonymref{利害}{li4hai5}
  \synonymref{猛烈}{meng3lie4}
  \synonymref{凶恶}{xiong1'e4}
  \antonymref{温和}{wen1he2}
  \antonymref{温顺}{wen1shun4}
\end{EntryWithPhonetic}

\begin{EntryWithPhonetic}{凶手}{xiong1shou3}{4,4}{⼐,⼿}[HSK 6]
  \definition[名]{s.}{assassino; homicida | agressor (que causou ferimentos a alguém)}
\end{EntryWithPhonetic}

%%%%%%%%%% 兄 %%%%%%%%%%
\subsection*{兄}\addcontentsline{loh}{figure}{兄 \dpy{xiong1}}

\begin{EntryWithPhonetic}{兄}{xiong1}{5}{⼉}
  \definition{s.}{irmão mais velho | parente mais velho do sexo masculino da mesma geração | uma forma cortês de tratamento entre amigos homens; um título respeitoso para amigos homens}
\end{EntryWithPhonetic}

\begin{EntryWithPhonetic}{兄弟}{xiong1di5}{5,7}{⼉,⼸}[HSK 4]
  \definition{adj.}{fraternal}
  \definition{pron.}{eu, me (termo de uso humilde por homens em discurso público)}
  \definition[个,位]{s.}{irmãos; irmão}
\end{EntryWithPhonetic}

%%%%%%%%%% 匈 %%%%%%%%%%
\subsection*{匈}\addcontentsline{loh}{figure}{匈 \dpy{xiong1}}

\begin{EntryWithPhonetic}{匈}{xiong1}{6}{⼓}
  \definition*{s.}{Hungria, abreviação de 匈牙利}
  \definition{s.}{peito; seio; tórax}
  \seealsoref{匈牙利}{xiong1ya2li4}
\end{EntryWithPhonetic}

\begin{EntryWithPhonetic}{匈奴}{xiong1nu2}{6,5}{⼓,⼥}
  \definition*{s.}{Xiongnu, um povo da estepe oriental que criou um império que floresceu na época das dinastias Qin e Han}
\end{EntryWithPhonetic}

\begin{EntryWithPhonetic}{匈牙利}{xiong1ya2li4}{6,4,7}{⼓,⽛,⼑}
  \definition*{s.}{Hungria}
\end{EntryWithPhonetic}

%%%%%%%%%% 汹 %%%%%%%%%%
\subsection*{汹}\addcontentsline{loh}{figure}{汹 \dpy{xiong1}}

\begin{EntryWithPhonetic}{汹}{xiong1}{7}{⽔}
  \definition{adj.}{turbulento; tempestuoso | rugindo; estrondoso | tumultuado}
\end{EntryWithPhonetic}

\begin{EntryWithPhonetic}{汹涌}{xiong1yong3}{7,10}{⽔,⽔}[HSK 7-9]
  \definition{adj.}{turbulento}
  \definition{adj.}{turbulento; tempestuoso; descreve uma quantidade imensa de água, em constante agitação e borbulhamento; também pode ser usada metaforicamente para descrever a chegada em larga escala de multidões, emoções ou mudanças sociais}
  \antonymref{平静}{ping2jing4}
\end{EntryWithPhonetic}

%%%%%%%%%% 胸 %%%%%%%%%%
\subsection*{胸}\addcontentsline{loh}{figure}{胸 \dpy{xiong1}}

\begin{EntryWithPhonetic}{胸}{xiong1}{10}{⾁}
  \definition{s.}{peito; seio; tórax | mente; coração; o coração humano}
\end{EntryWithPhonetic}

\begin{EntryWithPhonetic}{胸部}{xiong1bu4}{10,10}{⾁,⾢}[HSK 4]
  \definition{s.}{peito; tórax; seios}
\end{EntryWithPhonetic}

\begin{EntryWithPhonetic}{胸怀}{xiong1huai2}{10,7}{⾁,⼼}
  \definition[片,副]{s.}{mente; coração | peito}
  \definition{v.}{apreciar; guardar na memória; ter em mente}
  \seealsoref{心胸}{xin1xiong1}
\end{EntryWithPhonetic}

\begin{EntryWithPhonetic}{胸膛}{xiong1tang2}{10,15}{⾁,⾁}[HSK 7-9]
  \definition{s.}{peito; tórax; parte do tronco, entre o pescoço e o abdômen}
  \synonymref{胸部}{xiong1bu4}
\end{EntryWithPhonetic}

\begin{EntryWithPhonetic}{胸有成竹}{xiong1you3cheng2zhu2}{10,6,6,6}{⾁,⽉,⼽,⽵}[HSK 7-9]
  \definition{expr.}{ter um plano bem elaborado; planejar com antecedência; ao pintar bambu, a imagem do bambu já está na mente, isso significa, metaforicamente, ter um plano bem elaborado antes de fazer algo ou ter muita confiança na própria capacidade de realizar algo; implica ter considerado tudo minuciosamente antes de agir}
  \antonymref{不知所措}{bu4zhi1-suo3cuo4}
\end{EntryWithPhonetic}

%%%%%%%%%% 雄 %%%%%%%%%%
\subsection*{雄}\addcontentsline{loh}{figure}{雄 \dpy{xiong2}}

\begin{EntryWithPhonetic}{雄}{xiong2}{12}{⾫}
  \definition*{s.}{Sobrenome: Xiong}
  \definition{adj.}{masculino | grandioso; imponente; audacioso | poderoso}
  \definition{s.}{uma pessoa ou país com grande poder e influência}
\end{EntryWithPhonetic}

\begin{EntryWithPhonetic}{雄厚}{xiong2hou4}{12,9}{⾫,⼚}[HSK 7-9]
  \definition{adj.}{sólido; forte; abundante; (recursos humanos, recursos materiais, recursos financeiros, etc.) são mais do que suficientes}
  \synonymref{丰富}{feng1fu4}
  \synonymref{健壮}{jian4zhuang4}
  \synonymref{强壮}{qiang2zhuang4}
  \antonymref{薄弱}{bo2ruo4}
  \antonymref{单薄}{dan1bo2}
\end{EntryWithPhonetic}

\begin{EntryWithPhonetic}{雄伟}{xiong2wei3}{12,6}{⾫,⼈}[HSK 5]
  \definition{adj.}{magnífico; magnificente | imponente; magnífico}
\end{EntryWithPhonetic}

%%%%%%%%%% 熊 %%%%%%%%%%
\subsection*{熊}\addcontentsline{loh}{figure}{熊 \dpy{xiong2}}

\begin{EntryWithPhonetic}{熊}{xiong2}{14}{⽕}[HSK 5]
  \definition*{s.}{Sobrenome: Xiong}
  \definition[头,只]{s.}{urso}
  \definition{v.}{repreender; censurar}
\end{EntryWithPhonetic}

\begin{EntryWithPhonetic}{熊猫}{xiong2mao1}{14,11}{⽕,⽝}
  \definition[只,头]{s.}{panda gigante}
  \seealsoref{猫熊}{mao1xiong2}
\end{EntryWithPhonetic}

%%%%%%%%%% 休 %%%%%%%%%%
\subsection*{休}\addcontentsline{loh}{figure}{休 \dpy{xiu1}}

\begin{EntryWithPhonetic}{休}{xiu1}{6}{⼈}
  \definition{adj.}{feliz; alegre; festivo}
  \definition{adv.}{não; indica proibição ou dissuasão, equivalente a 别 ou 不要}
  \definition{s.}{fortuna e infortúnio; bom e mau}
  \definition{v.}{parar; cessar | descansar | abandonar a esposa e mandá"-la para casa; antigamente, o marido mandava a esposa de volta para a casa dos pais e rompia o relacionamento conjugal}
  \seealsoref{别}{bie2}
  \seealsoref{不要}{bu2yao4}
  \synonymref{息}{xi1}
  \synonymref{歇}{xie1}
\end{EntryWithPhonetic}

\begin{EntryWithPhonetic}{休兵}{xiu1bing1}{6,7}{⼈,⼋}
  \definition{s.}{armistício; cessar fogo}
  \definition{v.}{cessar fogo}
\end{EntryWithPhonetic}

\begin{EntryWithPhonetic}{休假}{xiu1/jia4}{6,11}{⼈,⼈}[HSK 2]
  \definition{v.+compl.}{ter um feriado; tirar férias; sair de férias}
  \synonymref{度假}{du4jia4}
  \synonymref{放假}{fang4/jia4}
  \synonymref{休憩}{xiu1qi4}
  \synonymref{休息}{xiu1xi5}
  \antonymref{工作}{gong1zuo4}
  \antonymref{开学}{kai1xue2}
  \antonymref{上班}{shang4/ban1}
\end{EntryWithPhonetic}

\begin{EntryWithPhonetic}{休克}{xiu1ke4}{6,7}{⼈,⼗}[HSK 7-9]
  \definition{s.}{choque; trata"-se de uma síndrome de hipóxia celular aguda causada por trauma grave, hemorragia maciça, infecção grave, envenenamento, etc.; os principais sintomas incluem queda da pressão arterial, palidez, calafrios, fraqueza e até coma}
  \definition{v.}{entrar em estado de choque; entar em choque}
  \synonymref{昏迷}{hun1mi2}
  \synonymref{窒息}{zhi4xi1}
\end{EntryWithPhonetic}

\begin{EntryWithPhonetic}{休眠}{xiu1mian2}{6,10}{⼈,⽬}[HSK 7-9]
  \definition{v.}{permanecer dormente | estar inativo (vulcão) | Biologia: estar dormente | Computação: hibernar}
  \synonymref{睡眠}{shui4mian2}
\end{EntryWithPhonetic}

\begin{EntryWithPhonetic}{休憩}{xiu1qi4}{6,16}{⼈,⼼}
  \definition{v.}{Literário: fazer uma pausa; descansar; relaxar}
  \synonymref{停息}{ting2xi1}
  \synonymref{停歇}{ting2xie1}
  \synonymref{歇息}{xie1xi5}
  \synonymref{休整}{xiu1zheng3}
  \synonymref{休息}{xiu1xi5}
  \synonymref{暂停}{zan4ting2}
  \antonymref{干活}{gan4/huo2}
  \antonymref{工作}{gong1zuo4}
  \antonymref{劳动}{lao2dong5}
  \antonymref{忙碌}{mang2lu4}
\end{EntryWithPhonetic}

\begin{EntryWithPhonetic}{休息室}{xiu1xi1shi4}{6,10,9}{⼈,⼼,⼧}
  \definition{s.}{saguão; \emph{lounge}; \emph{lobby}}
\end{EntryWithPhonetic}

\begin{EntryWithPhonetic}{休息}{xiu1xi5}{6,10}{⼈,⼼}[HSK 1]
  \definition{s.}{descanço}
  \definition{v.}{descansar; descansar um pouco; fazer uma pausa; interromper o trabalho, os estudos ou as atividades para recuperar as energias | dormir}
  \seealsoref{歇息}{xie1xi5}
  \synonymref{停息}{ting2xi1}
  \synonymref{停歇}{ting2xie1}
  \synonymref{歇息}{xie1xi5}
  \synonymref{休憩}{xiu1qi4}
  \synonymref{休养}{xiu1yang3}
  \synonymref{暂停}{zan4ting2}
  \antonymref{工作}{gong1zuo4}
  \antonymref{劳动}{lao2dong5}
  \antonymref{忙碌}{mang2lu4}
\end{EntryWithPhonetic}

\begin{EntryWithPhonetic}{休闲}{xiu1xian2}{6,7}{⼈,⾨}[HSK 5]
  \definition{s.}{ócio; lazer; tempo livre}
  \definition{v.}{desfrutar do lazer; sair de férias; aproveitar o tempo livre; parar de trabalhar ou estudar, estar em um estado de lazer e descontração | ficar ocioso}
\end{EntryWithPhonetic}

\begin{EntryWithPhonetic}{休想}{xiu1xiang3}{6,13}{⼈,⼼}[HSK 7-9]
  \definition{v.}{pare de sonhar; nunca pensar nisso; não imaginar que seja possível; não fantasiar}
  \synonymref{妄想}{wang4xiang3}
  \antonymref{但愿}{dan4yuan4}
  \antonymref{可以}{ke3yi3}
\end{EntryWithPhonetic}

\begin{EntryWithPhonetic}{休养}{xiu1yang3}{6,9}{⼈,⼋}[HSK 7-9]
  \definition{v.}{recuperar; convalescer | (força nacional) recuperar; revitalizar; reabilitar}
  \synonymref{疗养}{liao2yang3}
  \antonymref{繁忙}{fan2mang2}
\end{EntryWithPhonetic}

\begin{EntryWithPhonetic}{休整}{xiu1zheng3}{6,16}{⼈,⽁}
  \definition{v.}{(de tropas, etc.) descansar e passar por reorganização | descansar e reorganizar (militar); recarregar}
  \synonymref{调整}{tiao2zheng3}
  \synonymref{停歇}{ting2xie1}
  \synonymref{休憩}{xiu1qi4}
  \synonymref{休息}{xiu1xi5}
  \synonymref{整顿}{zheng3dun4}
\end{EntryWithPhonetic}

%%%%%%%%%% 修 %%%%%%%%%%
\subsection*{修}\addcontentsline{loh}{figure}{修 \dpy{xiu1}}

\begin{EntryWithPhonetic}{修}{xiu1}{9}{⼈}[HSK 3]
  \definition*{s.}{Sobrenome: Xiu}
  \definition{adj.}{comprido; alto e esbelto}
  \definition{s.}{revisionismo}
  \definition{v.}{embelezar; decorar | consertar; reparar; reformar | escrever; redigir; compilar | estudar; cultivar; aprender e praticar para aperfeiçoar ou melhorar (o caráter e o conhecimento) | construir; edificar | cortar ou aparar, para deixar bonito e arrumado | dedicar"-se à prática da religião}
\end{EntryWithPhonetic}

\begin{EntryWithPhonetic}{修补}{xiu1bu3}{9,7}{⼈,⾐}[HSK 7-9]
  \definition{v.}{consertar; remendar; reformar; reparar; renovar; consertar coisas quebradas para torná"-las inteiras | reparar; quando os tecidos de um organismo são danificados, as proteínas do corpo os repõem}
  \synonymref{修理}{xiu1li3}
  \antonymref{破坏}{po4huai4}
  \antonymref{损坏}{sun3huai4}
  \antonymref{损伤}{sun3shang1}
\end{EntryWithPhonetic}

\begin{EntryWithPhonetic}{修长}{xiu1chang2}{9,4}{⼈,⾧}[HSK 7-9]
  \definition{adj.}{alto e magro; esbelto; delgado | magro}
  \synonymref{苗条}{miao2tiao5}
  \synonymref{悠久}{you1jiu3}
  \antonymref{矮胖}{ai3pang4}
\end{EntryWithPhonetic}

\begin{EntryWithPhonetic}{修车}{xiu1che1}{9,4}{⼈,⾞}[HSK 6]
  \definition{v.}{consertar uma bicicleta (carro etc.)}[我打算明天去修车。===Pretendo consertar meu carro amanhã.]
\end{EntryWithPhonetic}

\begin{EntryWithPhonetic}{修订}{xiu1ding4}{9,4}{⼈,⾔}[HSK 7-9]
  \definition{v.}{revisar; reformular; revisar e corrigir}
\end{EntryWithPhonetic}

\begin{EntryWithPhonetic}{修复}{xiu1fu4}{9,9}{⼈,⼢}[HSK 5]
  \definition{v.}{reparar; restaurar; renovar | reparar; melhorar e restaurar (o relacionamento)}
\end{EntryWithPhonetic}

\begin{EntryWithPhonetic}{修改}{xiu1gai3}{9,7}{⼈,⽁}[HSK 3]
  \definition{v.}{revisar; retocar; corrigir erros e falhas em artigos, planos, etc.}
\end{EntryWithPhonetic}

\begin{EntryWithPhonetic}{修规}{xiu1gui1}{9,8}{⼈,⾒}
  \definition{s.}{plano de construção}
\end{EntryWithPhonetic}

\begin{EntryWithPhonetic}{修建}{xiu1jian4}{9,8}{⼈,⼵}[HSK 5]
  \definition{v.}{construir; erguer; animar; edificar; construir com tijolos, telhas, madeira, cimento, areia, etc.}
\end{EntryWithPhonetic}

\begin{EntryWithPhonetic}{修理}{xiu1li3}{9,11}{⼈,⽟}[HSK 4]
  \definition{v.}{consertar; reparar; restaurar algo danificado à sua forma ou função original | aparar; podar; cortar com tesouras e outras ferramentas para deixar árvores, flores, cabelos, etc. arrumados | culpar; punir; criticar ou punir uma pessoa para mostrar que ela está errada}
\end{EntryWithPhonetic}

\begin{EntryWithPhonetic}{修路}{xiu1/lu4}{9,13}{⼈,⾜}[HSK 7-9]
  \definition{v.+compl.}{reparar uma estrada; construir estradas}
\end{EntryWithPhonetic}

\begin{EntryWithPhonetic}{修养}{xiu1yang3}{9,9}{⼈,⼋}[HSK 5]
  \definition[种]{s.}{treinamento; domínio; realização; refere"-se a um determinado nível em termos de teoria, conhecimento, arte, pensamento, etc. | auto"-cultivo; refere"-se à atitude e ao comportamento cultivados ao longo do tempo, em conformidade com as exigências sociais}
  \seealsoref{教养}{jiao4yang3}
  \synonymref{教养}{jiao4yang3}
  \synonymref{素养}{su4yang3}
\end{EntryWithPhonetic}

\begin{EntryWithPhonetic}{修正}{xiu1zheng4}{9,5}{⼈,⽌}[HSK 7-9]
  \definition{v.}{revisar; alterar; corrigir; atualizar; corrigir erros de comportamento, pontos de vista, teorias, etc., para torná"-los corretos | emendar; alterar certas leis}
  \seealsoref{修改}{xiu1gai3}
  \synonymref{改进}{gai3jin4}
  \synonymref{改良}{gai3liang2}
  \synonymref{改正}{gai3zheng4}
  \synonymref{矫正}{jiao3zheng4}
  \synonymref{纠正}{jiu1zheng4}
  \synonymref{修改}{xiu1gai3}
\end{EntryWithPhonetic}

%%%%%%%%%% 羞 %%%%%%%%%%
\subsection*{羞}\addcontentsline{loh}{figure}{羞 \dpy{xiu1}}

\begin{EntryWithPhonetic}{羞}{xiu1}{10}{⽺}
  \definition{s.}{histórias de vergonha; coisa vergonhosa; coisas embaraçosas | comida deliciosa}
  \definition{v.}{envergonhar; fazer alguém se sentir envergonhado; fazer as pessoas ficarem constrangidas | sentir vergonha; sentir culpa; sentir"-se culpado | ter medo de ser ridicularizado pelos outros; sentir"-se envergonhado, constrangido}
  \seealsoref{馐}{xiu1}
\end{EntryWithPhonetic}

\begin{EntryWithPhonetic}{羞愧}{xiu1kui4}{10,12}{⽺,⼼}[HSK 7-9]
  \definition{adj.}{envergonhado; constrangido}
  \seealsoref{惭愧}{can2kui4}
  \synonymref{惭愧}{can2kui4}
  \synonymref{窘迫}{jiong3po4}
  \synonymref{以下}{yi3xia4}
  \synonymref{自责}{zi4ze2}
  \antonymref{高傲}{gao1'ao4}
  \antonymref{骄傲}{jiao1'ao4}
  \antonymref{自豪}{zi4hao2}
\end{EntryWithPhonetic}

%%%%%%%%%% 馐 %%%%%%%%%%
\subsection*{馐}\addcontentsline{loh}{figure}{馐 \dpy{xiu1}}

\begin{EntryWithPhonetic}{馐}{xiu1}{13}{⾷}
  \definition{s.}{iguaria; guloseima; comida deliciosa}
\end{EntryWithPhonetic}

%%%%%%%%%% 宿 %%%%%%%%%%
\subsection*{宿}\addcontentsline{loh}{figure}{宿 \dpy{xiu3}}

\begin{EntryWithPhonetic}{宿}{xiu3}{11}{⼧}
  \definition{s.}{usado para calcular a noite}[谈了半宿。===Conversamos por metade da noite.]
  \seeref{su4}
  \seeref{xiu4}
\end{EntryWithPhonetic}

%%%%%%%%%% 秀 %%%%%%%%%%
\subsection*{秀}\addcontentsline{loh}{figure}{秀 \dpy{xiu4}}

\begin{EntryWithPhonetic}{秀}{xiu4}{7}{⽲}
  \definition*{s.}{Sobrenome: Xiu}
  \definition{adj.}{elegante; belo | esperto; inteligente; astuto | excelente}
  \definition[场]{s.}{pessoa excelente; pessoa talentosa | Empréstimo Linguístico: \emph{show}}
  \definition{v.}{(culturas de grãos, etc.) produzir flores ou espigas}
\end{EntryWithPhonetic}

\begin{EntryWithPhonetic}{秀丽}{xiu4li4}{7,7}{⽲,⼀}[HSK 7-9]
  \definition{adj.}{bonito; belo; atraente; delicado e belo}
  \synonymref{灿烂}{can4lan4}
  \synonymref{美丽}{mei3li4}
  \synonymref{鲜艳}{xian1yan4}
  \synonymref{秀美}{xiu4mei3}
  \synonymref{艳丽}{yan4li4}
  \antonymref{丑陋}{chou3lou4}
\end{EntryWithPhonetic}

\begin{EntryWithPhonetic}{秀美}{xiu4mei3}{7,9}{⽲,⽺}[HSK 7-9]
  \definition{adj.}{elegante; delicado; gracioso; delicado e belo}
  \synonymref{俊俏}{jun4qiao4}
  \synonymref{美丽}{mei3li4}
  \synonymref{秀丽}{xiu4li4}
  \synonymref{艳丽}{yan4li4}
\end{EntryWithPhonetic}

%%%%%%%%%% 绣 %%%%%%%%%%
\subsection*{绣}\addcontentsline{loh}{figure}{绣 \dpy{xiu4}}

\begin{EntryWithPhonetic}{绣}{xiu4}{10}{⽷}[HSK 7-9]
  \definition[幅]{s.}{bordado}
  \definition{v.}{bordar; usar fios de seda, veludo ou algodão coloridos para criar padrões, imagens ou palavras em seda, tecido e outros materiais têxteis}
\end{EntryWithPhonetic}

%%%%%%%%%% 臭 %%%%%%%%%%
\subsection*{臭}\addcontentsline{loh}{figure}{臭 \dpy{xiu4}}

\begin{EntryWithPhonetic}{臭}{xiu4}{10}{⾃}
  \definition{s.}{odor; cheiro}
  \definition{v.}{cheirar; farejar; o mesmo que 嗅}
  \seeref{chou4}
  \seealsoref{嗅}{xiu4}
\end{EntryWithPhonetic}

%%%%%%%%%% 袖 %%%%%%%%%%
\subsection*{袖}\addcontentsline{loh}{figure}{袖 \dpy{xiu4}}

\begin{EntryWithPhonetic}{袖}{xiu4}{10}{⾐}
  \definition[个]{s.}{manga (de camisa, de camiseta, etc.)}
  \definition{v.}{dobrar para dentro da manga}
  \seealsoref{袖儿}{xiu4r5}
\end{EntryWithPhonetic}

\begin{EntryWithPhonetic}{袖儿}{xiu4r5}{10,2}{⾐,⼉}
  \definition{s.}{manga (de camisa, camiseta, etc.)}
\end{EntryWithPhonetic}

\begin{EntryWithPhonetic}{袖手旁观}{xiu4shou3-pang2guan1}{10,4,10,6}{⾐,⼿,⽅,⾒}[HSK 7-9]
  \definition{expr.}{``Observar de braços cruzados.''; olhar sem levantar um dedo; essa metáfora descreve alguém que se mantém longe de problemas ou não ajuda os outros; observar (ou ficar de braços cruzados); observar com indiferença}
\end{EntryWithPhonetic}

\begin{EntryWithPhonetic}{袖珍}{xiu4zhen1}{10,9}{⾐,⽟}[HSK 6]
  \definition{adj.}{do tamanho do bolso; de bolso (livro, agenda, etc.)}
\end{EntryWithPhonetic}

%%%%%%%%%% 宿 %%%%%%%%%%
\subsection*{宿}\addcontentsline{loh}{figure}{宿 \dpy{xiu4}}

\begin{EntryWithPhonetic}{宿}{xiu4}{11}{⼧}
  \definition{s.}{(astronomia) um termo antigo para constelação}
  \seeref{su4}
  \seeref{xiu3}
\end{EntryWithPhonetic}

%%%%%%%%%% 锈 %%%%%%%%%%
\subsection*{锈}\addcontentsline{loh}{figure}{锈 \dpy{xiu4}}

\begin{EntryWithPhonetic}{锈}{xiu4}{12}{⾦}[HSK 7-9]
  \definition{s.}{ferrugem; óxidos em superfícies metálicas}
  \definition{v.}{enferrujar}
\end{EntryWithPhonetic}

%%%%%%%%%% 嗅 %%%%%%%%%%
\subsection*{嗅}\addcontentsline{loh}{figure}{嗅 \dpy{xiu4}}

\begin{EntryWithPhonetic}{嗅}{xiu4}{13}{⼝}
  \definition{v.}{cheirar; farejar; identificar odores pelo nariz}
\end{EntryWithPhonetic}

\begin{EntryWithPhonetic}{嗅觉}{xiu4jue2}{13,9}{⼝,⾒}[HSK 7-9]
  \definition{s.}{sentido do olfato; sensação olfativa; a sensação produzida quando a mucosa nasal entra em contato com moléculas gasosas de certas substâncias}
  \antonymref{触觉}{chu4jue2}
  \antonymref{视觉}{shi4jue2}
\end{EntryWithPhonetic}

%%%%%%%%%% 须 %%%%%%%%%%
\subsection*{须}\addcontentsline{loh}{figure}{须 \dpy{xu1}}

\begin{EntryWithPhonetic}{须}{xu1}{9}{⾴}[HSK 7-9]
  \definition*{s.}{Sobrenome: Xu}
  \definition{s.}{barba; bigode; originalmente referindo"-se à barba que crescia no queixo, posteriormente passou a se referir à barba em geral | estruturas semelhantes a barbas em plantas e animais}
  \definition{v.}{dever; ter que; precisar | aguardar; ter necessidade; dever}
\end{EntryWithPhonetic}

%%%%%%%%%% 虚 %%%%%%%%%%
\subsection*{虚}\addcontentsline{loh}{figure}{虚 \dpy{xu1}}

\begin{EntryWithPhonetic}{虚}{xu1}{11}{⾌}[HSK 7-9]
  \definition*{s.}{Xu, a décima primeira das vinte e oito constelações em que a esfera celeste foi dividida, consistindo de duas estrelas em linha reta, uma em Aquário e a outra em Equuleus | Xu, uma das mansões lunares}
  \definition{adj.}{ficção; falso | vazio; ausência; período de vazio | modesto; humilde | fraco; frágil | inseguro; tímido; hesitante; medroso; falta de coragem devido a sentimentos de vergonha ou incerteza}
  \definition{adv.}{em vão; por nada}
  \definition{s.}{moralidade; o conceito mais geral (ideologia, princípios, políticas); os princípios, diretrizes e políticas da política | lacunas; pontos fracos}
  \definition{v.}{reservar espaço}
  \seealsoref{实}{shi2}
  \antonymref{实}{shi2}
\end{EntryWithPhonetic}

\begin{EntryWithPhonetic}{虚构}{xu1gou4}{11,8}{⾌,⽊}[HSK 7-9]
  \definition{adj.}{fictício; imaginário}
  \definition{v.}{inventar; fabricar; criar; fabricar do nada}
  \synonymref{编造}{bian1zao4}
  \synonymref{虚拟}{xu1ni3}
\end{EntryWithPhonetic}

\begin{EntryWithPhonetic}{虚幻}{xu1huan4}{11,4}{⾌,⼳}[HSK 7-9]
  \definition{adj.}{irreal; ilusório; imaginário; vago e ilusório; falso e inverídico}
  \synonymref{错觉}{cuo4jue2}
  \synonymref{梦幻}{meng4huan4}
  \synonymref{虚假}{xu1jia3}
  \antonymref{事实}{shi4shi2}
  \antonymref{现实}{xian4shi2}
  \antonymref{真实}{zhen1shi2}
\end{EntryWithPhonetic}

\begin{EntryWithPhonetic}{虚假}{xu1jia3}{11,11}{⾌,⼈}[HSK 7-9]
  \definition{adj.}{falso; simulado; fictício; desonesto; inverdadeiro; inconsistente com a realidade}
  \seealsoref{虚伪}{xu1wei3}
  \synonymref{荒谬}{huang1miu4}
  \synonymref{虚幻}{xu1huan4}
  \synonymref{虚伪}{xu1wei3}
  \antonymref{诚实}{cheng2shi5}
  \antonymref{确实}{que4shi2}
  \antonymref{实力}{shi2li4}
  \antonymref{实在}{shi2zai5}
  \antonymref{信誉}{xin4yu4}
  \antonymref{由衷}{you2zhong1}
  \antonymref{真诚}{zhen1cheng2}
  \antonymref{真实}{zhen1shi2}
  \antonymref{真正}{zhen1zheng4}
  \antonymref{衷心}{zhong1xin1}
\end{EntryWithPhonetic}

\begin{EntryWithPhonetic}{虚拟}{xu1ni3}{11,7}{⾌,⼿}[HSK 7-9]
  \definition{adj.}{fictício; teórico; hipotético; não condizente com os fatos ou não necessariamente condizente com eles | Computação: virtual}
  \definition{v.}{imaginar; inventar | Computação: emular}
  \synonymref{编造}{bian1zao4}
  \synonymref{虚构}{xu1gou4}
  \antonymref{真实}{zhen1shi2}
\end{EntryWithPhonetic}

\begin{EntryWithPhonetic}{虚弱}{xu1ruo4}{11,10}{⾌,⼸}[HSK 7-9]
  \definition{adj.}{fraco; debilitado; com saúde precária | fraco; débil; flácido; (força nacional e força militar) são fracas e vazias}
  \synonymref{薄弱}{bo2ruo4}
  \synonymref{脆弱}{cui4ruo4}
  \synonymref{单薄}{dan1bo2}
  \synonymref{软弱}{ruan3ruo4}
  \synonymref{衰老}{shuai1lao3}
  \synonymref{衰弱}{shuai1ruo4}
  \synonymref{无力}{wu2li4}
  \antonymref{健康}{jian4kang1}
  \antonymref{健壮}{jian4zhuang4}
  \antonymref{结实}{jie1shi5}
  \antonymref{强大}{qiang2da4}
  \antonymref{强壮}{qiang2zhuang4}
\end{EntryWithPhonetic}

\begin{EntryWithPhonetic}{虚伪}{xu1wei3}{11,6}{⾌,⼈}[HSK 7-9]
  \definition{adj.}{falso; não confiável; fabricado; hipócrita}
  \seealsoref{虚假}{xu1jia3}
  \synonymref{荒谬}{huang1miu4}
  \synonymref{卖弄}{mai4nong5}
  \synonymref{冒充}{mao4chong1}
  \synonymref{虚假}{xu1jia3}
  \antonymref{诚恳}{cheng2ken3}
  \antonymref{诚挚}{cheng2zhi4}
  \antonymref{诚实}{cheng2shi5}
  \antonymref{老实}{lao3shi5}
  \antonymref{事实}{shi4shi2}
  \antonymref{天真}{tian1zhen1}
  \antonymref{真诚}{zhen1cheng2}
  \antonymref{真实}{zhen1shi2}
  \antonymref{真正}{zhen1zheng4}
  \antonymref{衷心}{zhong1xin1}
\end{EntryWithPhonetic}

\begin{EntryWithPhonetic}{虚心}{xu1xin1}{11,4}{⾌,⼼}[HSK 5]
  \definition{adj.}{modesto; humilde; de mente aberta; não ser presunçoso, ser capaz de aceitar as opiniões dos outros}
\end{EntryWithPhonetic}

%%%%%%%%%% 需 %%%%%%%%%%
\subsection*{需}\addcontentsline{loh}{figure}{需 \dpy{xu1}}

\begin{EntryWithPhonetic}{需}{xu1}{14}{⾬}[HSK 7-9]
  \definition*{s.}{Sobrenome: Xu}
  \definition{s.}{necessidades; bens de primeira necessidade}
  \definition{v.}{precisar; querer; exigir}
  \seealsoref{需要}{xu1yao4}
\end{EntryWithPhonetic}

\begin{EntryWithPhonetic}{需求}{xu1qiu2}{14,7}{⾬,⽔}[HSK 3]
  \definition[种]{s.}{necessidades; demanda; exigência; solicitações decorrentes de necessidades}
\end{EntryWithPhonetic}

\begin{EntryWithPhonetic}{需要}{xu1yao4}{14,9}{⾬,⾑}[HSK 3]
  \definition[种]{s.}{necessidade; desejo ou exigência em relação a algo}
  \definition{v.}{precisar; querer; exigir; demandar; solicitar}
\end{EntryWithPhonetic}

%%%%%%%%%% 徐 %%%%%%%%%%
\subsection*{徐}\addcontentsline{loh}{figure}{徐 \dpy{xu2}}

\begin{EntryWithPhonetic}{徐}{xu2}{10}{⼻}
  \definition*{s.}{Sobrenome: Xu}
  \definition{adv.}{lentamente; suavemente}
\end{EntryWithPhonetic}

\begin{EntryWithPhonetic}{徐徐}{xu2xu2}{10,10}{⼻,⼻}[HSK 7-9]
  \definition{adv.}{lentamente; suavemente}
  \seealsoref{缓}{huan3}
  \seealsoref{缓缓}{huan3huan3}
  \seealsoref{慢}{man4}
  \seealsoref{慢慢}{man4man4}
  \synonymref{迟迟}{chi2chi2}
  \synonymref{缓缓}{huan3huan3}
  \synonymref{缓慢}{huan3man4}
  \synonymref{渐渐}{jian4jian4}
  \synonymref{慢慢}{man4man4}
  \antonymref{匆匆}{cong1cong1}
\end{EntryWithPhonetic}

%%%%%%%%%% 许 %%%%%%%%%%
\subsection*{许}\addcontentsline{loh}{figure}{许 \dpy{xu3}}

\begin{EntryWithPhonetic}{许}{xu3}{6}{⾔}[HSK 7-9]
  \definition*{s.}{Xu, um estado da Dinastia Zhou | Sobrenome: Xu}
  \definition{adv.}{um pouco;  talvez; expressa especulação ou estimativa}
  \definition{part.}{cerca de; aproximadamente; usado depois de certos numerais, frases de quantidade ou 些 ou 少 para indicar um número próximo a um certo número}
  \definition{pron.}{muitos; um monte de}
  \definition{v.}{elogiar; aprovar | prometer; prometer dar antecipadamente; dedicar | permitir; concordar; aprovar | (uma menina) estar prometida a; refere"-se especificamente ao noivado}
  \seealsoref{或者}{huo4zhe3}
  \seealsoref{可能}{ke3neng2}
  \seealsoref{少}{shao3}
  \seealsoref{些}{xie1}
\end{EntryWithPhonetic}

\begin{EntryWithPhonetic}{许多}{xu3duo1}{6,6}{⾔,⼣}[HSK 2]
  \definition{num.}{muitos; muito; numerosos; uma grande quantidade de}
  \seealsoref{多}{duo1}
  \synonymref{不少}{bu4shao3}
  \synonymref{大量}{da4liang4}
  \synonymref{多}{duo1}
  \synonymref{好多}{hao3duo1}
  \synonymref{众多}{zhong4duo1}
  \antonymref{稀少}{xi1shao3}
  \antonymref{些许}{xie1xu3}
\end{EntryWithPhonetic}

\begin{EntryWithPhonetic}{许可}{xu3ke3}{6,5}{⾔,⼝}[HSK 5]
  \definition{v.}{permitir; autorizar}
\end{EntryWithPhonetic}

\begin{EntryWithPhonetic}{许可证}{xu3ke3zheng4}{6,5,7}{⾔,⼝,⾔}[HSK 7-9]
  \definition{s.}{concessão; licença; permissão; permissão por escrito para fazer algo}
\end{EntryWithPhonetic}

%%%%%%%%%% 旭 %%%%%%%%%%
\subsection*{旭}\addcontentsline{loh}{figure}{旭 \dpy{xu4}}

\begin{EntryWithPhonetic}{旭}{xu4}{6}{⽇}
  \definition*{s.}{Sobrenome: Xu}
  \definition{s.}{o sol nascente; os primeiros raios de sol}
\end{EntryWithPhonetic}

\begin{EntryWithPhonetic}{旭日}{xu4ri4}{6,4}{⽇,⽇}[HSK 7-9]
  \definition{s.}{o sol nascente; o sol da manhã}
  \antonymref{落日}{luo4ri4}
  \antonymref{夕阳}{xi1yang2}
\end{EntryWithPhonetic}

%%%%%%%%%% 序 %%%%%%%%%%
\subsection*{序}\addcontentsline{loh}{figure}{序 \dpy{xu4}}

\begin{EntryWithPhonetic}{序}{xu4}{7}{⼴}[HSK 7-9]
  \definition*{s.}{Sobrenome: Xu}
  \definition{adj.}{introdutório; inicial; no início; antes do conteúdo principal}
  \definition[篇]{s.}{ordem; sequência | prelúdio; prefácio | escolas administradas por autoridades locais na antiguidade}
  \definition{v.}{organizar em ordem}
\end{EntryWithPhonetic}

\begin{EntryWithPhonetic}{序幕}{xu4mu4}{7,13}{⼴,⼱}[HSK 7-9]
  \definition{s.}{prelúdio; prólogo; cena de abertura de uma peça teatral; a cena que precede o primeiro ato serve para introduzir ou preparar o terreno para toda a apresentação (ópera, balé, peça teatral, etc.)}
  \antonymref{尾声}{wei3sheng1}
\end{EntryWithPhonetic}

%%%%%%%%%% 叙 %%%%%%%%%%
\subsection*{叙}\addcontentsline{loh}{figure}{叙 \dpy{xu4}}

\begin{EntryWithPhonetic}{叙}{xu4}{9}{⼜}
  \definition{s.}{introdução; prefácio}
  \definition{v.}{conversar; bater papo | narrar; recontar; relatar | avaliar; analisar | classificar e avaliar}
  \synonymref{道}{dao4}
  \synonymref{讲}{jiang3}
  \synonymref{说}{shuo1}
  \synonymref{谈}{tan2}
\end{EntryWithPhonetic}

\begin{EntryWithPhonetic}{叙述}{xu4shu4}{9,8}{⼜,⾡}[HSK 7-9]
  \definition{v.}{relatar; narrar; recontar; dar um relato de; registrar ou relatar os eventos conforme ocorreram}
  \seealsoref{阐述}{chan3shu4}
  \synonymref{阐述}{chan3shu4}
\end{EntryWithPhonetic}

%%%%%%%%%% 恤 %%%%%%%%%%
\subsection*{恤}\addcontentsline{loh}{figure}{恤 \dpy{xu4}}

\begin{EntryWithPhonetic}{恤}{xu4}{9}{⼼}
  \definition{v.}{ter pena; simpatizar | dar alívio; compensar}
\end{EntryWithPhonetic}

%%%%%%%%%% 畜 %%%%%%%%%%
\subsection*{畜}\addcontentsline{loh}{figure}{畜 \dpy{xu4}}

\begin{EntryWithPhonetic}{畜}{xu4}{10}{⽥}
  \definition{v.}{criar (animais domésticos)}
  \seeref{chu4}
\end{EntryWithPhonetic}

%%%%%%%%%% 续 %%%%%%%%%%
\subsection*{续}\addcontentsline{loh}{figure}{续 \dpy{xu4}}

\begin{EntryWithPhonetic}{续}{xu4}{11}{⽷}[HSK 7-9]
  \definition*{s.}{Sobrenome: Xu}
  \definition{v.}{continuar; estender; unir | adicionar; fornecer mais; adicionar à}
  \antonymref{断}{duan4}
  \antonymref{绝}{jue2}
\end{EntryWithPhonetic}

%%%%%%%%%% 酗 %%%%%%%%%%
\subsection*{酗}\addcontentsline{loh}{figure}{酗 \dpy{xu4}}

\begin{EntryWithPhonetic}{酗}{xu4}{11}{⾣}
  \definition{v.}{ter tendência ao consumo excessivo de álcool}
\end{EntryWithPhonetic}

\begin{EntryWithPhonetic}{酗酒}{xu4jiu3}{11,10}{⾣,⾣}[HSK 7-9]
  \definition{v.}{beber muito; beber em excesso; entregar"-se ao consumo excessivo de álcool; um comportamento caracterizado pelo consumo frequente de álcool em excesso e pela perda de controle, geralmente por um longo período; também se refere à perda de controle após a embriaguez}
\end{EntryWithPhonetic}

%%%%%%%%%% 絮 %%%%%%%%%%
\subsection*{絮}\addcontentsline{loh}{figure}{絮 \dpy{xu4}}

\begin{EntryWithPhonetic}{絮}{xu4}{12}{⽷}
  \definition{adj.}{entediado; farto | prolixo; tagarela; falante; prolixo}
  \definition[瓣]{s.}{(algodão) enchimento; adicionando | Obsoleto: fio de seda grosso | amentilho | lgo semelhante ao algodão}
  \definition{v.}{encher com algodão | tagarelar; ser prolixo; ser loquaz}
\end{EntryWithPhonetic}

\begin{EntryWithPhonetic}{絮叨}{xu4dao5}{12,5}{⽷,⼝}[HSK 7-9]
  \definition{v.}{falar sem parar sem chegar ao ponto; ser prolixo; ser tagarela; ser verborrágico}
  \synonymref{唠叨}{lao2dao5}
\end{EntryWithPhonetic}

%%%%%%%%%% 宣 %%%%%%%%%%
\subsection*{宣}\addcontentsline{loh}{figure}{宣 \dpy{xuan1}}

\begin{EntryWithPhonetic}{宣}{xuan1}{9}{⼧}
  \definition*{s.}{Sobrenome: Xuan}
  \definition{v.}{declarar; proclamar; anunciar; falar publicamente | drenar (líquidos)}
\end{EntryWithPhonetic}

\begin{EntryWithPhonetic}{宣布}{xuan1bu4}{9,5}{⼧,⼱}[HSK 3]
  \definition{v.}{declarar; proclamar; pronunciar; anunciar; informar oficialmente a todos sobre as últimas decisões e situações}
\end{EntryWithPhonetic}

\begin{EntryWithPhonetic}{宣称}{xuan1cheng1}{9,10}{⼧,⽲}[HSK 7-9]
  \definition{v.}{afirmar; declarar; professar | afirmar; pronunciar; alegar}
  \synonymref{传播}{chuan2bo1}
  \synonymref{鼓吹}{gu3chui1}
  \synonymref{声称}{sheng1cheng1}
  \synonymref{宣传}{xuan1chuan2}
  \synonymref{宣扬}{xuan1yang2}
  \synonymref{转播}{zhuan3bo1}
\end{EntryWithPhonetic}

\begin{EntryWithPhonetic}{宣传}{xuan1chuan2}{9,6}{⼧,⼈}[HSK 3]
  \definition[个]{v.}{propagar; divulgar; fazer propaganda; explicar e esclarecer às pessoas, para que elas acreditem e sigam as ações}
\end{EntryWithPhonetic}

\begin{EntryWithPhonetic}{宣读}{xuan1du2}{9,10}{⼧,⾔}[HSK 7-9]
  \definition{v.}{ler em voz alta; dizer em voz alta}
  \synonymref{发表}{fa1biao3}
  \synonymref{朗读}{lang3du2}
  \antonymref{默读}{mo4du2}
\end{EntryWithPhonetic}

\begin{EntryWithPhonetic}{宣告}{xuan1gao4}{9,7}{⼧,⼝}[HSK 7-9]
  \definition{v.}{declarar; proclamar; pronunciar; anunciar, informar solenemente o público; transmitir um tom solene}
  \seealsoref{宣布}{xuan1bu4}
  \synonymref{颁布}{ban1bu4}
  \synonymref{颁发}{ban1fa1}
  \synonymref{发表}{fa1biao3}
  \synonymref{公告}{gong1gao4}
  \synonymref{揭晓}{jie1xiao3}
  \synonymref{宣布}{xuan1bu4}
\end{EntryWithPhonetic}

\begin{EntryWithPhonetic}{宣誓}{xuan1/shi4}{9,14}{⼧,⾔}[HSK 7-9]
  \definition{v.+compl.}{fazer um juramento; jurar; fazer um voto (ou promessa) | prestar juramento; fazer um juramento; fazer uma promessa; ao assumir uma tarefa ou ingressar em uma organização, expressar publicamente a própria determinação em uma cerimônia específica}
  \synonymref{发誓}{fa1/shi4}
  \synonymref{宣言}{xuan1yan2}
\end{EntryWithPhonetic}

\begin{EntryWithPhonetic}{宣泄}{xuan1xie4}{9,8}{⼧,⽔}[HSK 7-9]
  \definition{v.}{drenar; escoar (líquidos); deixar a água fluir | desabafar; tirar algo do peito; relaxar; revelar}
  \synonymref{暴露}{bao4lu4}
  \synonymref{发泄}{fa1xie4}
  \synonymref{透露}{tou4lu4}
  \synonymref{泄漏}{xie4lou4}
  \synonymref{泄露}{xie4lou4}
\end{EntryWithPhonetic}

\begin{EntryWithPhonetic}{宣言}{xuan1yan2}{9,7}{⼧,⾔}[HSK 7-9]
  \definition[篇,份]{s.}{declaração; manifesto; declarações emitidas por (uma nação, partido político ou organização) expressando publicamente suas opiniões sobre questões importantes com o objetivo de publicidade e apelo}
  \definition{v.}{declarar; proclamar; afirmar}
  \synonymref{承诺}{cheng2nuo4}
  \synonymref{诺言}{nuo4yan2}
  \synonymref{宣誓}{xuan1/shi4}
\end{EntryWithPhonetic}

\begin{EntryWithPhonetic}{宣扬}{xuan1yang2}{9,6}{⼧,⼿}[HSK 7-9]
  \definition{v.}{defender; anunciar; divulgar; propagar; divulgar e espalhar amplamente a notícia}
  \seealsoref{宣传}{xuan1chuan2}
  \synonymref{传播}{chuan2bo1}
  \synonymref{鼓动}{gu3dong4}
  \synonymref{流传}{liu2chuan2}
  \synonymref{流转}{liu2zhuan3}
  \synonymref{散步}{san4/bu4}
  \synonymref{散布}{san4bu4}
  \synonymref{宣称}{xuan1cheng1}
  \synonymref{宣传}{xuan1chuan2}
  \synonymref{张扬}{zhang1yang2}
  \synonymref{转播}{zhuan3bo1}
\end{EntryWithPhonetic}

%%%%%%%%%% 喧 %%%%%%%%%%
\subsection*{喧}\addcontentsline{loh}{figure}{喧 \dpy{xuan1}}

\begin{EntryWithPhonetic}{喧}{xuan1}{12}{⼝}
  \definition{adj.}{barulhento; ruidoso}
  \definition{v.}{fazer barulho}
  \seealsoref{吵架}{chao3/jia4}
  \seealsoref{吵闹}{chao3nao4}
  \seealsoref{争吵}{zheng1chao3}
  \synonymref{吵}{chao3}
\end{EntryWithPhonetic}

\begin{EntryWithPhonetic}{喧哗}{xuan1hua2}{12,9}{⼝,⼝}[HSK 7-9]
  \definition{adj.}{barulhento; ruidoso}
  \definition{s.}{alvoroço; tumulto; confusão; ruído confuso; som alto e caótico}
  \definition{v.}{causar alvoroço; fazer alarde}
  \synonymref{吵闹}{chao3nao4}
  \synonymref{忙乱}{mang2luan4}
  \synonymref{热闹}{re4nao5}
  \synonymref{喧闹}{xuan1nao4}
  \antonymref{安静}{an1jing4}
  \antonymref{寂静}{ji4jing4}
  \antonymref{寂寥}{ji4liao2}
  \antonymref{宁静}{ning2jing4}
  \antonymref{僻静}{pi4jing4}
\end{EntryWithPhonetic}

\begin{EntryWithPhonetic}{喧闹}{xuan1nao4}{12,8}{⼝,⾾}[HSK 7-9]
  \definition{adj.}{barulhento; agitado; tumultuoso; estridente}
  \definition{v.}{agitar"-se; clamar; fazer barulho; causar alvoroço}
  \synonymref{吵闹}{chao3nao4}
  \synonymref{热闹}{re4nao5}
  \synonymref{喧哗}{xuan1hua2}
  \synonymref{争吵}{zheng1chao3}
  \antonymref{安静}{an1jing4}
  \antonymref{寂静}{ji4jing4}
  \antonymref{寂寞}{ji4mo4}
  \antonymref{宁静}{ning2jing4}
  \antonymref{平静}{ping2jing4}
  \antonymref{悄悄}{qiao1qiao1}
\end{EntryWithPhonetic}

%%%%%%%%%% 玄 %%%%%%%%%%
\subsection*{玄}\addcontentsline{loh}{figure}{玄 \dpy{xuan2}}

\begin{EntryWithPhonetic}{玄}{xuan2}{5}{⽞}[HSK 7-9][Kangxi 95]
  \definition{adj.}{preto; escuro | profundo; abstruso; escondido | não confiável; irrealista; não confiável}
\end{EntryWithPhonetic}

\begin{EntryWithPhonetic}{玄机}{xuan2ji1}{5,6}{⽞,⽊}[HSK 7-9]
  \definition{s.}{Religião: verdade arcana | segredo obscuro (ou profundo) | princípios misteriosos | teoria profunda (no taoísmo e no budismo)}
\end{EntryWithPhonetic}

\begin{EntryWithPhonetic}{玄学}{xuan2xue2}{5,8}{⽞,⼦}
  \definition{s.}{Escola Philosófica Wei e Jin amalgamando os ideais daoísta e confucionistas | tradução da metafísica (形而上学) | Obsoleto: metafísica}
  \seealsoref{形而上学}{xing2'er2shang4xue2}
\end{EntryWithPhonetic}

%%%%%%%%%% 悬 %%%%%%%%%%
\subsection*{悬}\addcontentsline{loh}{figure}{悬 \dpy{xuan2}}

\begin{EntryWithPhonetic}{悬}{xuan2}{11}{⼼}[HSK 6]
  \definition{adj.}{pendente; não resolvido; sem nenhum resultado | distante; a distância é grande; a diferença é grande | Dialeto: perigoso}
  \definition{v.}{pendurar; suspender | levantar; elevar | sentir"-se ansioso; ser solícito | imaginar}
\end{EntryWithPhonetic}

\begin{EntryWithPhonetic}{悬挂}{xuan2gua4}{11,9}{⼼,⼿}[HSK 7-9]
  \definition{v.}{pendurar; pender; suspender; prender um objeto em um ou mais pontos em algum lugar com a ajuda de uma corda, gancho, prego, etc.}
\end{EntryWithPhonetic}

\begin{EntryWithPhonetic}{悬念}{xuan2nian4}{11,8}{⼼,⼼}[HSK 7-9]
  \definition[处]{s.}{suspense (sentido à medida que uma história, peça teatral, etc., se desenvolve até um clímax); o sentimento de preocupação com o desenvolvimento da história e o destino dos personagens ao apreciar uma peça de teatro, um filme ou outra obra de arte}
  \definition{v.}{preocupar"-se com; ficar pensando em; manter em suspense}
  \synonymref{担心}{dan1/xin1}
  \synonymref{惦记}{dian4ji4}
  \synonymref{顾虑}{gu4lv4}
  \synonymref{挂念}{gua4nian4}
  \synonymref{怀念}{huai2nian4}
  \synonymref{牵挂}{qian1gua4}
  \synonymref{思念}{si1nian4}
  \synonymref{想念}{xiang3nian4}
  \antonymref{放心}{fang4/xin1}
\end{EntryWithPhonetic}

\begin{EntryWithPhonetic}{悬殊}{xuan2shu1}{11,10}{⼼,⽍}[HSK 7-9]
  \definition{adj.}{vastamente diferentes; polos opostos}
  \synonymref{差距}{cha1ju4}
  \antonymref{相当}{xiang1dang1}
\end{EntryWithPhonetic}

\begin{EntryWithPhonetic}{悬心吊胆}{xuan2xin1-diao4dan3}{11,4,6,9}{⼼,⼼,⼝,⾁}
  \definition{expr.}{de tirar o fôlego; estar ansioso e com medo; estar à beira de um ataque de nervos; ter o coração na boca; estar em suspense}
  \synonymref{提心吊胆}{ti2xin1-diao4dan3}
\end{EntryWithPhonetic}

\begin{EntryWithPhonetic}{悬崖}{xuan2ya2}{11,11}{⼼,⼭}[HSK 7-9]
  \definition[处]{s.}{penhasco; precipício; escarpa}
\end{EntryWithPhonetic}

%%%%%%%%%% 旋 %%%%%%%%%%
\subsection*{旋}\addcontentsline{loh}{figure}{旋 \dpy{xuan2}}

\begin{EntryWithPhonetic}{旋}{xuan2}{11}{⽅}
  \definition*{s.}{Sobrenome: Xuan}
  \definition{adv.}{em breve; rapidamente}
  \definition{s.}{redemoinho; turbilhão; vórtice}
  \definition{v.}{girar; circular; rodar | retornar; voltar}
\end{EntryWithPhonetic}

\begin{EntryWithPhonetic}{旋律}{xuan2lv4}{11,9}{⽅,⼻}[HSK 7-9]
  \definition[段]{s.}{canto; melodia; na música, refere"-se à combinação regular e rítmica de várias notas musicais, sendo o principal meio de expressar o conteúdo, o estilo, o gênero e as características nacionais da música | melodia; metáfora para o movimento harmonioso das coisas}
\end{EntryWithPhonetic}

\begin{EntryWithPhonetic}{旋涡}{xuan2wo1}{11,10}{⽅,⽔}[HSK 7-9]
  \definition{s.}{redemoinho; vórtice; turbilhão | vorticidade; vórtices | espiral}
\end{EntryWithPhonetic}

\begin{EntryWithPhonetic}{旋转}{xuan2zhuan3}{11,8}{⽅,⾞}[HSK 6]
  \definition{v.}{girar; rodar; revolver; rodopiar; o movimento circular de um objeto em torno de um ponto ou eixo}
\end{EntryWithPhonetic}

%%%%%%%%%% 选 %%%%%%%%%%
\subsection*{选}\addcontentsline{loh}{figure}{选 \dpy{xuan3}}

\begin{EntryWithPhonetic}{选}{xuan3}{9}{⾡}[HSK 2]
  \definition{s.}{pessoa ou coisa selecionada | seleções; antologia; trabalhos selecionados e compilados}
  \definition{v.}{selecionar; escolher | eleger}
  \seealsoref{当选}{dang1xuan3}
  \seealsoref{挑}{tiao1}
  \seealsoref{挑选}{tiao1xuan3}
  \seealsoref{推选}{tui1xuan3}
  \seealsoref{选举}{xuan3ju3}
  \synonymref{当选}{dang1xuan3}
  \synonymref{拣}{jian3}
  \synonymref{挑}{tiao1}
  \synonymref{挑选}{tiao1xuan3}
  \synonymref{推选}{tui1xuan3}
  \synonymref{选拔}{xuan3ba2}
  \synonymref{选举}{xuan3ju3}
\end{EntryWithPhonetic}

\begin{EntryWithPhonetic}{选拔}{xuan3ba2}{9,8}{⾡,⼿}[HSK 6]
  \definition{v.}{selecionar; escolher}
\end{EntryWithPhonetic}

\begin{EntryWithPhonetic}{选举}{xuan3ju3}{9,9}{⾡,⼂}[HSK 6]
  \definition[次,个]{s.}{eleição; as eleições são o processo pelo qual os cidadãos escolhem os seus representantes ou líderes através do voto}
  \definition{v.}{votar; eleger; eleger representantes ou responsáveis votando ou levantando as mãos}
\end{EntryWithPhonetic}

\begin{EntryWithPhonetic}{选民}{xuan3min2}{9,5}{⾡,⽒}[HSK 7-9]
  \definition{s.}{eleitor; refere"-se a um indivíduo | circunscrição eleitoral; eleitorado; refere"-se ao todo}
  \synonymref{公民}{gong1min2}
\end{EntryWithPhonetic}

\begin{EntryWithPhonetic}{选手}{xuan3shou3}{9,4}{⾡,⼿}[HSK 3]
  \definition[位,名,个,些]{s.}{jogador; (selecionado) competidor; atleta selecionado para uma competição esportiva; participantes selecionados entre um grande número de candidatos}
\end{EntryWithPhonetic}

\begin{EntryWithPhonetic}{选项}{xuan3xiang4}{9,9}{⾡,⾴}[HSK 7-9]
  \definition{s.}{uma escolha; uma opção; uma alternativa; itens disponíveis para seleção ou escolha}
  \definition{v.}{fazer uma escolha (entre várias alternativas)}
\end{EntryWithPhonetic}

\begin{EntryWithPhonetic}{选修}{xuan3xiu1}{9,9}{⾡,⼈}[HSK 5]
  \definition{v.}{fazer como disciplina eletiva; escolher entre uma seleção de cursos disponíveis}
\end{EntryWithPhonetic}

\begin{EntryWithPhonetic}{选用}{xuan3yong4}{9,5}{⾡,⽤}[HSK 7-9]
  \definition{v.}{selecionar para emprego ou para uso | escolher para algum propósito | selecionar e usar}
  \seealsoref{选择}{xuan3ze2}
  \synonymref{采纳}{cai3na4}
  \synonymref{采取}{cai3qu3}
  \synonymref{采用}{cai3yong4}
  \synonymref{选择}{xuan3ze2}
\end{EntryWithPhonetic}

\begin{EntryWithPhonetic}{选择}{xuan3ze2}{9,8}{⾡,⼿}[HSK 4]
  \definition[个,种,次]{s.}{escolha; opção; resultado da escolha; possibilidade de escolha}
  \definition{v.}{selecionar; escolher}
\end{EntryWithPhonetic}

%%%%%%%%%% 炫 %%%%%%%%%%
\subsection*{炫}\addcontentsline{loh}{figure}{炫 \dpy{xuan4}}

\begin{EntryWithPhonetic}{炫}{xuan4}{9}{⽕}
  \definition{adj.}{deslumbrante; ofuscante}
  \definition{v.}{exibir; ostentar}
\end{EntryWithPhonetic}

\begin{EntryWithPhonetic}{炫耀}{xuan4yao4}{9,20}{⽕,⽻}[HSK 7-9]
  \definition{v.}{ostentar; exibir; fazer uma demonstração de; fazer as pessoas perceberem que elas têm muito dinheiro, grandes habilidades ou um status elevado | deslumbrar}
  \synonymref{高调}{gao1diao4}
  \synonymref{夸耀}{kua1yao4}
  \synonymref{卖弄}{mai4nong5}
  \antonymref{谦虚}{qian1xu1}
  \antonymref{谦逊}{qian1xun4}
\end{EntryWithPhonetic}

%%%%%%%%%% 削 %%%%%%%%%%
\subsection*{削}\addcontentsline{loh}{figure}{削 \dpy{xue1}}

\begin{EntryWithPhonetic}{削}{xue1}{9}{⼑}
  \definition{v.}{aparar; talhar; cortar; é usado exclusivamente em palavras compostas, como 剥削, 减削, 弱弱}
  \seeref{xiao1}
  \seealsoref{剥削}{bo1xue1}
  \seealsoref{减削}{jian3xue1}
  \seealsoref{弱弱}{ruo4 ruo4}
\end{EntryWithPhonetic}

\begin{EntryWithPhonetic}{削弱}{xue1ruo4}{9,10}{⼑,⼸}[HSK 7-9]
  \definition{v.}{enfraquecer; incapacitar; prejudicar; causar enfraquecimento}
  \seealsoref{减弱}{jian3ruo4}
  \synonymref{减弱}{jian3ruo4}
  \synonymref{减少}{jian3shao3}
  \synonymref{衰弱}{shuai1ruo4}
  \antonymref{充实}{chong1shi2}
  \antonymref{巩固}{gong3gu4}
  \antonymref{加强}{jia1qiang2}
  \antonymref{强化}{qiang2hua4}
  \antonymref{增强}{zeng1qiang2}
  \antonymref{壮大}{zhuang4da4}
\end{EntryWithPhonetic}

%%%%%%%%%% 靴 %%%%%%%%%%
\subsection*{靴}\addcontentsline{loh}{figure}{靴 \dpy{xue1}}

\begin{EntryWithPhonetic}{靴}{xue1}{13}{⾰}
  \definition[双]{s.}{botas}
\end{EntryWithPhonetic}

\begin{EntryWithPhonetic}{靴子}{xue1zi5}{13,3}{⾰,⼦}[HSK 7-9]
  \definition[只,双,打]{s.}{botas; sapatos com cabedal que se estende acima do osso do tornozelo}
\end{EntryWithPhonetic}

%%%%%%%%%% 薛 %%%%%%%%%%
\subsection*{薛}\addcontentsline{loh}{figure}{薛 \dpy{xue1}}

\begin{EntryWithPhonetic}{薛}{xue1}{16}{⾋}
  \definition*{s.}{Estado vassalo durante a Dinastia Zhou (1046--256 a.C.) | Sobrenome: Xue}
  \definition{s.}{erva semelhante ao absinto (clássico)}
\end{EntryWithPhonetic}

\begin{EntryWithPhonetic}{薛稷}{xue1 ji4}{16,15}{⾋,⽲}
  \definition*{s.}{Xue Ji (649--713), um dos quatro grandes calígrafos do início da dinastia Tang, 唐初四大家}
  \seealsoref{唐初四大家}{tang2chu1 si4 da4jia1}
\end{EntryWithPhonetic}

%%%%%%%%%% 穴 %%%%%%%%%%
\subsection*{穴}\addcontentsline{loh}{figure}{穴 \dpy{xue2}}

\begin{EntryWithPhonetic}{穴}{xue2}{5}{⽳}[Kangxi 116]
  \definition*{s.}{Sobrenome: Xue}
  \definition{s.}{caverna; toca; buraco; cavidade; oco | cova de caixão | Medicina Chinesa: acupressão; ponto de acupuntura; local | túmulo; sepultura; vala comum; cova}
\end{EntryWithPhonetic}

\begin{EntryWithPhonetic}{穴位}{xue2wei4}{5,7}{⽳,⼈}[HSK 7-9]
  \definition{s.}{ponto de acupuntura; acupressão; em medicina, os pontos de acupuntura no corpo humano são áreas onde as terminações nervosas estão densamente agrupadas ou por onde passam fibras nervosas espessas}
\end{EntryWithPhonetic}

%%%%%%%%%% 学 %%%%%%%%%%
\subsection*{学}\addcontentsline{loh}{figure}{学 \dpy{xue2}}

\begin{EntryWithPhonetic}{学}{xue2}{8}{⼦}[HSK 1]
  \definition[所]{s.}{aprendizagem; conhecimento; sabedoria; erudição | objeto de estudo; ramo do conhecimento | escola; faculdade | teoria; doutrina}
  \definition{v.}{estudar; aprender | imitar; copiar}
  \seealsoref{学习}{xue2xi2}
  \synonymref{仿}{fang3}
  \synonymref{闻}{wen2}
  \synonymref{习}{xi2}
  \synonymref{学习}{xue2xi2}
  \synonymref{知}{zhi1}
  \synonymref{智}{zhi4}
  \antonymref{教}{jiao1}
  \antonymref{教}{jiao4}
\end{EntryWithPhonetic}

\begin{EntryWithPhonetic}{学费}{xue2fei4}{8,9}{⼦,⾙}[HSK 3]
  \definition[笔]{s.}{mensalidade (taxa); prêmio; taxas que os alunos devem pagar para estudar na escola, conforme estabelecido pela escola | preço pelo que se aprendeu ao custo do próprio bolso; a metáfora do preço a pagar para obter uma determinada experiência | custo; preço; todas as despesas necessárias durante o período de estudos do aluno}
\end{EntryWithPhonetic}

\begin{EntryWithPhonetic}{学分}{xue2fen1}{8,4}{⼦,⼑}[HSK 4]
  \definition{s.}{créditos de um curso; uma unidade de medida do peso e do tempo do curso no ensino superior; cada curso vale um crédito para uma aula por semana durante um semestre; alunos devem concluir o número necessário de créditos para se formar}
\end{EntryWithPhonetic}

\begin{EntryWithPhonetic}{学好}{xue2hao3}{8,6}{⼦,⼥}
  \definition{v.}{aprender com bons exemplos; imitar os bons | seguir bons exemplos; tomar as boas ações como exemplos e seguir o exemplo delas}
\end{EntryWithPhonetic}

\begin{EntryWithPhonetic}{学会}{xue2hui4}{8,6}{⼦,⼈}[HSK 6]
  \definition[个]{s.}{sociedade; instituto; sociedade científica; um grupo acadêmico composto por pessoas que estudam um determinado assunto, como a Sociedade de Física, a Sociedade de Biologia, etc.}
  \definition{v.}{aprender; dominar; aprender e aplicar}
\end{EntryWithPhonetic}

\begin{EntryWithPhonetic}{学科}{xue2ke1}{8,9}{⼦,⽲}[HSK 5]
  \definition[门,级]{s.}{ramo do aprendizado; disciplina | disciplina escolar; curso de estudo | cursos teóricos oferecidos em treinamento militar ou físico | disciplina acadêmica | curso | assunto; tema}
  \antonymref{术科}{shu4ke1}
\end{EntryWithPhonetic}

\begin{EntryWithPhonetic}{学历}{xue2li4}{8,4}{⼦,⼚}[HSK 7-9]
  \definition{s.}{credenciais acadêmicas; histórico de escolaridade formal; formação educacional; a experiência educacional, incluindo as escolas em que a pessoa se formou, geralmente se refere especificamente ao tipo de escola em que ela se graduou}
  \synonymref{文凭}{wen2ping2}
\end{EntryWithPhonetic}

\begin{EntryWithPhonetic}{学年}{xue2nian2}{8,6}{⼦,⼲}[HSK 4]
  \definition{s.}{ano letivo; ano acadêmico}
\end{EntryWithPhonetic}

\begin{EntryWithPhonetic}{学期}{xue2qi1}{8,12}{⼦,⽉}[HSK 2]
  \definition[个,段]{s.}{semestre; período escolar; um ano acadêmico é dividido em dois semestres, um semestre do início do outono até as férias de inverno e um semestre do início da primavera até as férias de verão}
\end{EntryWithPhonetic}

\begin{EntryWithPhonetic}{学生}{xue2sheng5}{8,5}{⼦,⽣}[HSK 1]
  \definition[个,名,群]{s.}{aluno; estudante; pupilo}
  \synonymref{弟子}{di4zi3}
  \synonymref{学员}{xue2yuan2}
  \synonymref{学子}{xue2zi3}
  \antonymref{教师}{jiao4shi1}
  \antonymref{教授}{jiao4shou4}
  \antonymref{老师}{lao3shi1}
  \antonymref{师长}{shi1zhang3}
  \antonymref{师傅}{shi1fu5}
  \antonymref{先生}{xian1sheng5}
\end{EntryWithPhonetic}

\begin{EntryWithPhonetic}{学生证}{xue2sheng5zheng4}{8,5,7}{⼦,⽣,⾔}
  \definition{s.}{documento (cartão) de identidade de estudante}
\end{EntryWithPhonetic}

\begin{EntryWithPhonetic}{学时}{xue2shi2}{8,7}{⼦,⽇}[HSK 4]
  \definition{s.}{hora"-aula; hora de aula; período}
\end{EntryWithPhonetic}

\begin{EntryWithPhonetic}{学士}{xue2shi4}{8,3}{⼦,⼠}[HSK 7-9]
  \definition{s.}{acadêmico | bacharel; o nível mais baixo de escolaridade, concedido pela universidade após a graduação}
  \synonymref{博士}{bo2shi4}
\end{EntryWithPhonetic}

\begin{EntryWithPhonetic}{学术}{xue2shu4}{8,5}{⼦,⽊}[HSK 4]
  \definition[种]{s.}{aprendizagem; aprendizado; ciências; aprendizado sistemático e especializado}
\end{EntryWithPhonetic}

\begin{EntryWithPhonetic}{学说}{xue2shuo1}{8,9}{⼦,⾔}[HSK 7-9]
  \definition[种]{s.}{teoria; doutrina; proposições ou percepções acadêmicas sistemáticas}
  \synonymref{观点}{guan1dian3}
  \synonymref{见解}{jian4jie3}
  \synonymref{学术}{xue2shu4}
\end{EntryWithPhonetic}

\begin{EntryWithPhonetic}{学堂}{xue2tang2}{8,11}{⼦,⼟}[HSK 7-9]
  \definition{s.}{escola (antigo) | faculdade}
  \synonymref{学校}{xue2xiao4}
\end{EntryWithPhonetic}

\begin{EntryWithPhonetic}{学位}{xue2wei4}{8,7}{⼦,⼈}[HSK 5]
  \definition[个]{s.}{grau; grau acadêmico; título concedido com base no nível acadêmico profissional, como doutorado, mestrado, etc.}
\end{EntryWithPhonetic}

\begin{EntryWithPhonetic}{学问}{xue2wen5}{8,6}{⼦,⾨}[HSK 4]
  \definition[门,种,个,项]{s.}{aprendizado, conhecimento, erudição; a compreensão correta do mundo objetivo que alguém tem | conhecimento; aprendizado sistemático; conhecimento sistemático sobre algo ou uma ciência que pode ser aprendido em um livro ou em uma experiência prática}
\end{EntryWithPhonetic}

\begin{EntryWithPhonetic}{学习}{xue2xi2}{8,3}{⼦,⼄}[HSK 1]
  \definition{s.}{estudo}
  \definition{v.}{estudar; aprender; adquirir conhecimentos ou habilidades através da leitura, da audição, da pesquisa e da prática}
  \seealsoref{学}{xue2}
  \synonymref{进修}{jin4xiu1}
  \synonymref{练习}{lian4xi2}
  \synonymref{效仿}{xiao4fang3}
  \synonymref{学}{xue2}
\end{EntryWithPhonetic}

\begin{EntryWithPhonetic}{学校}{xue2xiao4}{8,10}{⼦,⽊}[HSK 1]
  \definition[所,个]{s.}{escola; instituição de ensino}
  \synonymref{学堂}{xue2tang2}
\end{EntryWithPhonetic}

\begin{EntryWithPhonetic}{学业}{xue2ye4}{8,5}{⼦,⼀}[HSK 7-9]
  \definition{s.}{escolaridade; os estudos de alguém; conhecimento; aprendizagem | estudo; tarefa de casa; trabalho de curso}
  \synonymref{课程}{ke4cheng2}
\end{EntryWithPhonetic}

\begin{EntryWithPhonetic}{学艺}{xue2yi4}{8,4}{⼦,⾋}[HSK 7-9]
  \definition{s.}{conhecimentos e habilidades}
  \definition{v.}{aprender um ofício ou profissão | aprender uma habilidade ou arte}
\end{EntryWithPhonetic}

\begin{EntryWithPhonetic}{学员}{xue2yuan2}{8,7}{⼦,⼝}[HSK 6]
  \definition[位,名,批,个]{s.}{estudante; estagiário; geralmente se refere a pessoas que estudam em escolas ou cursos de treinamento diferentes de faculdades, escolas de ensino médio e escolas primárias}
\end{EntryWithPhonetic}

\begin{EntryWithPhonetic}{学院}{xue2yuan4}{8,9}{⼦,⾩}[HSK 1]
  \definition[个,所,家]{s.}{academia; instituto; um tipo de instituição de ensino superior que se concentra em uma determinada área de especialização, como faculdades de engenharia, faculdades de música, faculdades de educação, etc.}
\end{EntryWithPhonetic}

\begin{EntryWithPhonetic}{学者}{xue2zhe3}{8,8}{⼦,⽼}[HSK 5]
  \definition[位]{s.}{erudito; homem culto; pessoas que fazem pesquisas acadêmicas geralmente se referem àquelas que alcançaram certo sucesso acadêmico}
\end{EntryWithPhonetic}

\begin{EntryWithPhonetic}{学子}{xue2zi3}{8,3}{⼦,⼦}[HSK 7-9]
  \definition{s.}{Literário: estudante | acadêmico}
\end{EntryWithPhonetic}

%%%%%%%%%% 雪 %%%%%%%%%%
\subsection*{雪}\addcontentsline{loh}{figure}{雪 \dpy{xue3}}

\begin{EntryWithPhonetic}{雪}{xue3}{11}{⾬}[HSK 2]
  \definition*{s.}{Sobrenome: Xue}
  \definition[场,层]{s.}{neve | algo parecido com neve}
  \definition{v.}{limpar; enxugar; remover}
  \seealsoref{雪花}{xue3hua1}
  \synonymref{白}{bai2}
  \synonymref{除}{chu2}
  \synonymref{雪花}{xue3hua1}
\end{EntryWithPhonetic}

\begin{EntryWithPhonetic}{雪板}{xue3ban3}{11,8}{⾬,⽊}
  \definition{s.}{esqui; prancha de \emph{snowboard}}
  \definition{v.}{esquiar; praticar \textit{snowboard}}
\end{EntryWithPhonetic}

\begin{EntryWithPhonetic}{雪糕}{xue3gao1}{11,16}{⾬,⽶}
  \definition{s.}{neve; picolé; sorvete; um tipo de sobremesa fria; sua textura é semelhante à de sorvete e é servida em cubos}
  \synonymref{棒冰}{bang4bing1}
  \synonymref{冰糕}{bing1gao1}
  \synonymref{冰棍}{bing1gun4}
\end{EntryWithPhonetic}

\begin{EntryWithPhonetic}{雪花}{xue3hua1}{11,7}{⾬,⾋}
  \definition[片]{s.}{floco de neve}
  \seealsoref{雪}{xue3}
\end{EntryWithPhonetic}

\begin{EntryWithPhonetic}{雪葩}{xue3pa1}{11,12}{⾬,⾋}
  \definition{s.}{Empréstimo Linguístico: sorvete}
\end{EntryWithPhonetic}

\begin{EntryWithPhonetic}{雪人}{xue3ren2}{11,2}{⾬,⼈}
  \definition*{s.}{\emph{Yeti}; Abominável Homem das Neves}
  \definition{s.}{boneco de neve; uma figura feita de neve}
\end{EntryWithPhonetic}

\begin{EntryWithPhonetic}{雪山}{xue3shan1}{11,3}{⾬,⼭}[HSK 7-9]
  \definition{s.}{montanha com o pico coberto de neve; montanha nevada | montanha de neve; montanhas cobertas de neve durante todo o ano}
\end{EntryWithPhonetic}

\begin{EntryWithPhonetic}{雪上加霜}{xue3shang4-jia1shuang1}{11,3,5,17}{⾬,⼀,⼒,⾬}[HSK 7-9]
  \definition{expr.}{para piorar ainda mais a situação; agravando a situação; piorando ainda mais as coisas; jogando sal na ferida; essa metáfora descreve uma situação em que desastres são sofridos repetidamente e os danos se tornam cada vez mais graves}
\end{EntryWithPhonetic}

\begin{EntryWithPhonetic}{雪鞋}{xue3xie2}{11,15}{⾬,⾰}
  \definition[双]{s.}{sapatos de neve; são calçados usados para caminhar na neve}
\end{EntryWithPhonetic}

%%%%%%%%%% 血 %%%%%%%%%%
\subsection*{血}\addcontentsline{loh}{figure}{血 \dpy{xue4}}

\begin{EntryWithPhonetic}{血}{xue4}{6}{⾎}[HSK 3][Kangxi 143]
  \definition[滴,袋,口,毫升]{s.}{sangue | parente consanguíneo; com laços de parentesco | pessoa ativa e animada; metáfora para uma personalidade ou espírito forte e sincero | medicina tradicional chinesa refere"-se à menstruação}
  \seeref{xie3}
\end{EntryWithPhonetic}

\begin{EntryWithPhonetic}{血管}{xue4guan3}{6,14}{⾎,⽵}[HSK 6]
  \definition[根,条,种]{s.}{vaso; vaso sanguíneo; os canais tubulares pelos quais o sangue circula são divididos em três tipos: artérias, veias e capilares}
\end{EntryWithPhonetic}

\begin{EntryWithPhonetic}{血汗}{xue4han4}{6,6}{⾎,⽔}
  \definition{s.}{sangue e suor; suor e trabalho árduo}
  \synonymref{心血}{xin1xue4}
\end{EntryWithPhonetic}

\begin{EntryWithPhonetic}{血脉}{xue4mai4}{6,9}{⾎,⾁}[HSK 7-9]
  \definition{s.}{linhagem sanguínea; laços sanguíneos}
  \synonymref{脉络}{mai4luo4}
  \synonymref{血管}{xue4guan3}
  \synonymref{血缘}{xue4yuan2}
\end{EntryWithPhonetic}

\begin{EntryWithPhonetic}{血栓}{xue4shuan1}{6,10}{⾎,⽊}[HSK 7-9]
  \definition{s.}{coágulo sanguíneo; trombo | trombose}
\end{EntryWithPhonetic}

\begin{EntryWithPhonetic}{血压}{xue4ya1}{6,6}{⾎,⼚}[HSK 7-9]
  \definition{s.}{pressão arterial; a pressão do sangue contra as paredes dos vasos sanguíneos}
\end{EntryWithPhonetic}

\begin{EntryWithPhonetic}{血液}{xue4ye4}{6,11}{⾎,⽔}[HSK 6]
  \definition[毫升]{s.}{sangue | linha de vida; sangue vital; uma metáfora para o importante componente ou força que mantém a vitalidade coletiva}
\end{EntryWithPhonetic}

\begin{EntryWithPhonetic}{血缘}{xue4yuan2}{6,12}{⾎,⽷}[HSK 7-9]
  \definition{s.}{linhagem sanguínea; consanguinidade; laços sanguíneos; parentesco de sangue; refere"-se às relações naturais que se formam entre membros da família e do clã devido à reprodução}
\end{EntryWithPhonetic}

%%%%%%%%%% 勋 %%%%%%%%%%
\subsection*{勋}\addcontentsline{loh}{figure}{勋 \dpy{xun1}}

\begin{EntryWithPhonetic}{勋}{xun1}{9}{⼒}
  \definition{s.}{mérito; serviço meritório; conquista}
\end{EntryWithPhonetic}

\begin{EntryWithPhonetic}{勋章}{xun1zhang1}{9,11}{⼒,⾳}[HSK 7-9]
  \definition{s.}{medalha; condecoração; ordem (no Reino Unido)}
\end{EntryWithPhonetic}

%%%%%%%%%% 熏 %%%%%%%%%%
\subsection*{熏}\addcontentsline{loh}{figure}{熏 \dpy{xun1}}

\begin{EntryWithPhonetic}{熏}{xun1}{14}{⽕}[HSK 7-9]
  \definition{adj.}{quente; suave; ameno}
  \definition{v.}{fumar; fumigar; expor à fumaça ou vapor; (fumaça, vapores, etc.) que entram em contato com um objeto fazem com que ele mude de cor ou adquira um odor | defumar (alimentos); curar por meio da fumaça; métodos de processamento de alimentos que utilizam fumaça e fogo para grelhar os alimentos, conferindo"-lhes um sabor especial | permear; influenciar; estar exposto a; afetado por exposição prolongada}
  \seeref{xun4}
\end{EntryWithPhonetic}

\begin{EntryWithPhonetic}{熏陶}{xun1tao2}{14,10}{⽕,⾩}[HSK 7-9]
  \definition{v.}{edificar; nutrir; exercer uma influência edificante sobre; as pessoas com quem você convive ao longo do tempo exercem gradualmente uma influência positiva em seu estilo de vida, pensamentos, comportamento, caráter e conhecimento}
  \synonymref{教会}{jiao1hui4}
  \synonymref{教导}{jiao4dao3}
  \synonymref{教授}{jiao4shou4}
  \synonymref{教学}{jiao4xue2}
  \synonymref{教养}{jiao4yang3}
  \synonymref{教育}{jiao4yu4}
  \synonymref{陶冶}{tao2ye3}
\end{EntryWithPhonetic}

\begin{EntryWithPhonetic}{熏香}{xun1xiang1}{14,9}{⽕,⾹}
  \definition{s.}{incenso}
  \definition{v.}{exalar fumaça; fumigar}
\end{EntryWithPhonetic}

%%%%%%%%%% 寻 %%%%%%%%%%
\subsection*{寻}\addcontentsline{loh}{figure}{寻 \dpy{xun2}}

\begin{EntryWithPhonetic}{寻}{xun2}{6}{⼨}[HSK 7-9]
  \definition*{s.}{Sobrenome: Xun}
  \definition{clas.}{uma unidade antiga de comprimento, igual a 8 尺}
  \definition{v.}{procurar; buscar}
  \seeref{xin2}
  \seealsoref{尺}{chi3}
  \synonymref{找}{zhao3}
\end{EntryWithPhonetic}

\begin{EntryWithPhonetic}{寻常}{xun2chang2}{6,11}{⼨,⼱}[HSK 7-9]
  \definition{adj.}{comum; usual; ordinário}
  \synonymref{平常}{ping2chang2}
  \synonymref{普通}{pu3tong1}
  \synonymref{日常}{ri4chang2}
  \synonymref{通常}{tong1chang2}
  \antonymref{别致}{bie2zhi4}
  \antonymref{特殊}{te4shu1}
\end{EntryWithPhonetic}

\begin{EntryWithPhonetic}{寻觅}{xun2mi4}{6,8}{⼨,⾒}[HSK 7-9]
  \definition{v.}{procurar; buscar; procurar por toda parte, na esperança de ver alguém ou conseguir alguma coisa}
  \synonymref{失传}{shi1chuan2}
  \synonymref{探求}{tan4qiu2}
  \synonymref{探索}{tan4suo3}
  \synonymref{寻求}{xun2qiu2}
  \synonymref{寻找}{xun2zhao3}
  \synonymref{找寻}{zhao3xun2}
  \synonymref{追求}{zhui1qiu2}
  \antonymref{埋藏}{mai2cang2}
\end{EntryWithPhonetic}

\begin{EntryWithPhonetic}{寻求}{xun2qiu2}{6,7}{⼨,⽔}[HSK 5]
  \definition{v.}{procurar; perseguir; explorar; ir em busca de}
\end{EntryWithPhonetic}

\begin{EntryWithPhonetic}{寻找}{xun2zhao3}{6,7}{⼨,⼿}[HSK 4]
  \definition{v.}{buscar; procurar; pesquisar; encontrar, que pode ser usado tanto para coisas concretas quanto para coisas abstratas}
\end{EntryWithPhonetic}

%%%%%%%%%% 巡 %%%%%%%%%%
\subsection*{巡}\addcontentsline{loh}{figure}{巡 \dpy{xun2}}

\begin{EntryWithPhonetic}{巡}{xun2}{6}{⾡}
  \definition{clas.}{rodada de bebidas | usado para servir vinho a todos}
  \definition{v.}{patrulhar; fazer rondas; fazer uma excursão de inspeção}
\end{EntryWithPhonetic}

\begin{EntryWithPhonetic}{巡逻}{xun2luo2}{6,11}{⾡,⾡}[HSK 7-9]
  \definition{s.}{patrulha}
  \definition{v.}{patrulhar; sair em patrulha}
  \synonymref{察看}{cha2kan4}
\end{EntryWithPhonetic}

%%%%%%%%%% 询 %%%%%%%%%%
\subsection*{询}\addcontentsline{loh}{figure}{询 \dpy{xun2}}

\begin{EntryWithPhonetic}{询}{xun2}{8}{⾔}
  \definition{v.}{perguntar; indagar; reunir informações | consultar; buscar conselho}
\end{EntryWithPhonetic}

\begin{EntryWithPhonetic}{询问}{xun2wen4}{8,6}{⾔,⾨}[HSK 5]
  \definition{v.}{indagar; perguntar sobre; pedir conselho}
\end{EntryWithPhonetic}

%%%%%%%%%% 循 %%%%%%%%%%
\subsection*{循}\addcontentsline{loh}{figure}{循 \dpy{xun2}}

\begin{EntryWithPhonetic}{循}{xun2}{12}{⼻}
  \definition{v.}{seguir; cumprir; cumprir com}
\end{EntryWithPhonetic}

\begin{EntryWithPhonetic}{循环}{xun2huan2}{12,8}{⼻,⽟}[HSK 6]
  \definition{s.}{ciclo; circulação}
  \definition{v.}{circular; as coisas se movem ou mudam em um ciclo}
\end{EntryWithPhonetic}

\begin{EntryWithPhonetic}{循序渐进}{xun2xu4-jian4jin4}{12,7,11,7}{⼻,⼴,⽔,⾡}[HSK 7-9]
  \definition{expr.}{fazer progresso constante e incremental; avançar gradualmente na devida ordem; fazer as coisas passo a passo; isso se refere ao aprendizado ou trabalho de uma maneira que envolve o aprofundamento ou aprimoramento gradual da compreensão por meio de determinadas etapas}
\end{EntryWithPhonetic}

%%%%%%%%%% 训 %%%%%%%%%%
\subsection*{训}\addcontentsline{loh}{figure}{训 \dpy{xun4}}

\begin{EntryWithPhonetic}{训}{xun4}{5}{⾔}[HSK 7-9]
  \definition{s.}{instrução; ensinamento; ensino | padrão; modelo; exemplo; regra; diretriz | explicação ou interpretação crítica de um texto | treinamento; exercício}
  \definition{v.}{instruir; admoestar; dar uma palestra a alguém; ensinar | explicar; instruir; explicação do significado da palavra | treinar}
\end{EntryWithPhonetic}

\begin{EntryWithPhonetic}{训诂}{xun4gu3}{5,7}{⾔,⾔}
  \definition{s.}{estudos exegéticos (de textos antigos); exegese}
  \definition{v.}{explicação de palavras e frases em livros antigos | interpretar e elaborar glossários e comentários sobre textos clássicos}
\end{EntryWithPhonetic}

\begin{EntryWithPhonetic}{训练}{xun4lian4}{5,8}{⾔,⽷}[HSK 3]
  \definition{v.}{treinar; exercitar; planejar e executar de forma sistemática o desenvolvimento de habilidades ou competências específicas}
\end{EntryWithPhonetic}

%%%%%%%%%% 迅 %%%%%%%%%%
\subsection*{迅}\addcontentsline{loh}{figure}{迅 \dpy{xun4}}

\begin{EntryWithPhonetic}{迅}{xun4}{6}{⾡}
  \definition{adj.}{rápido; veloz}
  \definition{adv.}{rapidamente; velozmente}
\end{EntryWithPhonetic}

\begin{EntryWithPhonetic}{迅速}{xun4su4}{6,10}{⾡,⾡}[HSK 4]
  \definition{adv.}{rapidamente; velozmente; prontamente}
\end{EntryWithPhonetic}

%%%%%%%%%% 驯 %%%%%%%%%%
\subsection*{驯}\addcontentsline{loh}{figure}{驯 \dpy{xun4}}

\begin{EntryWithPhonetic}{驯}{xun4}{6}{⾺}[HSK 7-9]
  \definition{adj.}{manso e dócil; submisso; gentil}
  \definition{v.}{domesticar; amansar; tornar obediente}
\end{EntryWithPhonetic}

%%%%%%%%%% 逊 %%%%%%%%%%
\subsection*{逊}\addcontentsline{loh}{figure}{逊 \dpy{xun4}}

\begin{EntryWithPhonetic}{逊}{xun4}{9}{⾡}
  \definition*{s.}{Sobrenome: Xun}
  \definition{adj.}{modesto; humilde | Literário: inferior; não tão bom quanto; inferior a}
  \definition{v.}{abdicar; desistir}
\end{EntryWithPhonetic}

\begin{EntryWithPhonetic}{逊色}{xun4se4}{9,6}{⾡,⾊}[HSK 7-9]
  \definition{adj.}{inferior; indivíduos com habilidades ou condições inferiores não podem ser comparados}
  \antonymref{出色}{chu1se4}
  \antonymref{媲美}{pi4mei3}
\end{EntryWithPhonetic}

%%%%%%%%%% 熏 %%%%%%%%%%
\subsection*{熏}\addcontentsline{loh}{figure}{熏 \dpy{xun4}}

\begin{EntryWithPhonetic}{熏}{xun4}{14}{⽕}
  \definition{adj.}{Literário: quente; ameno}
  \definition{v.}{expor à fumaça ou vapores; fumigar | tratar (carnes, peixes, etc.) com fumaça; defumar | perfumar com incenso, etc. | sufocar; (gás) pode causar asfixia e envenenamento}
  \seeref{xun1}
\end{EntryWithPhonetic}

%%%%% EOF %%%%%

